% Sampling and Estimation — Pure Mathematics with Statistics Upper Sixth — Cours signé OnBuch+
% Official MINESEC syllabus. Compiler avec : tectonic PMS15-sampling-and-estimation.tex
\newcommand{\DOCMATIERE}{Pure Mathematics with Statistics}
\newcommand{\DOCNIVEAU}{Upper Sixth}
\newcommand{\DOCMODULE}{Upper Sixth — Pure mathematics with statistics}
\newcommand{\DOCLECON}{15}
\newcommand{\DOCTITRE}{Sampling and Estimation}
% OnBuch+ — préambule « cours signés » ENRICHI (compilé avec tectonic / XeLaTeX)
% Système visuel « Pop » d'OnBuch+ : fond crème (jamais blanc), bordures encre
% épaisses, ombres « dures » décalées, titres Archivo Black en capitales.
% Version MATHÉMATIQUES : graphes (pgfplots/tikz), boîtes pédagogiques
% (définition, propriété, méthode, attention, le savais-tu, exercice résolu,
% exercice, TP, bilan), page de garde et signature OnBuch+ ancrée en bas.
%
% Macros attendues dans main.tex AVANT \input{preamble} :
%   \DOCMATIERE  \DOCNIVEAU  \DOCMODULE  \DOCLECON (numéro)  \DOCTITRE
\documentclass[11pt]{article}

\usepackage{fontspec}
\usepackage[english]{babel}
\usepackage[a4paper,margin=2cm,top=3.4cm,bottom=2.8cm,headheight=1.4cm,footskip=1.3cm]{geometry}
\usepackage{amsmath,amssymb,amsfonts}
\usepackage{mathtools} % \xmapsto, \coloneqq... souvent utilisés par les modèles
\usepackage{cancel} % simplifications barrées (bilans de forces)
% \abs{z} (module, valeur absolue) et \norm{u} (norme), utilisés par les modèles
\providecommand{\abs}{}\let\abs\relax\DeclarePairedDelimiter{\abs}{\lvert}{\rvert}
\providecommand{\norm}{}\let\norm\relax\DeclarePairedDelimiter{\norm}{\lVert}{\rVert}
\usepackage[table]{xcolor} % [table] : fournit aussi \rowcolors (alternance de lignes)
\usepackage{pagecolor}
\usepackage{tikz}
\usetikzlibrary{arrows.meta,positioning,calc,shapes.geometric,shapes.misc,shapes.arrows,decorations.pathmorphing,decorations.pathreplacing,decorations.markings,patterns,shadows,mindmap,trees,fit,backgrounds,matrix,intersections,angles,quotes}
\usepackage{pgfplots}
\usepackage[version=4]{mhchem} % équations nucléaires \ce{^{235}_{92}U}
\usepackage[european,siunitx]{circuitikz} % schémas électriques
\pgfplotsset{compat=1.18}
\usepgfplotslibrary{fillbetween,groupplots} % aires entre courbes (calcul intégral)
\usepackage{enumitem}
\usepackage{fancyhdr}
% [most] charge toutes les bibliothèques tcolorbox (listings, minted, raster,
% documentation...) et gonfle inutilement la mémoire TeX sur un document long
% (~300 boîtes) : on ne charge que ce qui est réellement utilisé.
\usepackage[skins,breakable]{tcolorbox}
\usepackage{graphicx}
\usepackage{colortbl}
\usepackage{booktabs}
\usepackage{tabularx}
\usepackage{diagbox} % case d'en-tête diagonale des tableaux à double entrée
% Colonnes extensibles centrée / gauche / droite (C, L, R), souvent utilisées par les modèles
\newcolumntype{C}{>{\centering\arraybackslash}X}
\newcolumntype{L}{>{\raggedright\arraybackslash}X}
\newcolumntype{R}{>{\raggedleft\arraybackslash}X}
\usepackage{array}
\usepackage{multicol}
\usepackage{siunitx}
% parse-numbers=false : accepte aussi \SI{2,5 \times 10^{-3}}{...}
\sisetup{locale=UK,output-decimal-marker={.},group-minimum-digits=4,parse-numbers=false}
\DeclareSIUnit\ohm{\ensuremath{\symup{\Omega}}} % U+2126 absent de la police
\DeclareSIUnit\division{div} % réglages d'oscilloscope (ms/div, V/div)
\DeclareSIUnit\tr{tr} % tours (vitesse de rotation, ex. \SI{1500}{\tr\per\minute})
\usepackage{qrcode}
% \degree : commande fréquemment utilisée par les modèles pour un angle ou une
% température en dehors de \SI{}{} (non fournie par les paquets chargés).
\providecommand{\degree}{^{\circ}}
% \qed (fin de démonstration), utilisé par les modèles sans amsthm
\providecommand{\qed}{\hfill$\square$}
% \barre : double barre des valeurs interdites dans un tableau de variations (tkz-tab)
\providecommand{\barre}{\ensuremath{\|}}
% environnement proof (amsthm non chargé) utilisé par les modèles
\ifdefined\proof\else\newenvironment{proof}[1][Proof]{\par\noindent\textit{#1.}\ }{\qed\par}\fi
\usepackage{lastpage}
\usepackage{hyperref}
\hypersetup{hidelinks}

% ============================================================
% Palette « Pop » OnBuch+ — mêmes hex que la classe P (plus_ui.dart)
% ============================================================
\definecolor{poporange}{HTML}{F59321}
\definecolor{popdark}{HTML}{E07A0C}
\definecolor{popcream}{HTML}{FFF6E8}
\definecolor{popink}{HTML}{1C1714}
\definecolor{poppurple}{HTML}{7A5AE0}
\definecolor{popgreen}{HTML}{2E9E6A}
\definecolor{poppink}{HTML}{E0457B}
\definecolor{popblue}{HTML}{2F7DE1}
\definecolor{popgold}{HTML}{C98A12}
\definecolor{popmuted}{HTML}{8A7F76}
\definecolor{popline}{HTML}{EADFD2}
\definecolor{popgray}{HTML}{8A7F76}
\colorlet{popgrey}{popgray}
% Alias tolérants (le modèle nomme parfois les couleurs différemment)
\colorlet{popwhite}{white}
\colorlet{popblack}{popink}
\colorlet{popred}{poppink}
\colorlet{popyellow}{popgold}
% Teintes claires pour fonds de tableaux / zones
\colorlet{popblueL}{popblue!10!white}
\colorlet{poporangeL}{poporange!14!white}
\colorlet{popgreenL}{popgreen!12!white}
\colorlet{poppurpleL}{poppurple!12!white}
\colorlet{poppinkL}{poppink!12!white}
\colorlet{popgoldL}{popgold!14!white}

\pagecolor{popcream}

% ============================================================
% Typographie
% ============================================================
\defaultfontfeatures{Ligatures=TeX}
\setmainfont{PlusJakartaSans-Medium.ttf}[Path=../../../fonts/,BoldFont=PlusJakartaSans-Bold.ttf]
% Maths en sans-serif (Fira Math) avec chiffres et lettres droites de Plus
% Jakarta, pour que les formules se fondent dans le texte.
\usepackage[math-style=ISO,bold-style=ISO]{unicode-math}
\setmathfont{FiraMath-Regular.otf}[Path=../../../fonts/]
\setmathfont{PlusJakartaSans-Medium.ttf}[Path=../../../fonts/,range={up/num,up/{Latin,latin},\mathup}]
% (icomma retiré : décimales avec point en anglais)
% Caractères mathématiques Unicode absents des polices de texte (Plus Jakarta,
% Archivo Black) : rendus par leur équivalent mathématique, en texte comme en
% formule — sinon ils s'affichent en carré vide (ex. « Arithmétique dans ℤ »).
\usepackage{newunicodechar}
\newunicodechar{ℝ}{\ensuremath{\mathbb{R}}}
\newunicodechar{ℤ}{\ensuremath{\mathbb{Z}}}
\newunicodechar{ℂ}{\ensuremath{\mathbb{C}}}
\newunicodechar{ℕ}{\ensuremath{\mathbb{N}}}
\newunicodechar{ℚ}{\ensuremath{\mathbb{Q}}}
\newunicodechar{𝔻}{\ensuremath{\mathbb{D}}}
\newunicodechar{α}{\ensuremath{\alpha}}
\newunicodechar{β}{\ensuremath{\beta}}
\newunicodechar{γ}{\ensuremath{\gamma}}
\newunicodechar{δ}{\ensuremath{\delta}}
\newunicodechar{ε}{\ensuremath{\varepsilon}}
\newunicodechar{θ}{\ensuremath{\theta}}
\newunicodechar{λ}{\ensuremath{\lambda}}
\newunicodechar{μ}{\ensuremath{\mu}}
\newunicodechar{σ}{\ensuremath{\sigma}}
\newunicodechar{τ}{\ensuremath{\tau}}
\newunicodechar{φ}{\ensuremath{\varphi}}
\newunicodechar{ψ}{\ensuremath{\psi}}
\newunicodechar{ω}{\ensuremath{\omega}}
\newunicodechar{Σ}{\ensuremath{\Sigma}}
% majuscules produites par \MakeUppercase dans les titres (α-aminés -> Α)
\newunicodechar{Α}{\ensuremath{\alpha}}
\newunicodechar{Β}{\ensuremath{\beta}}
\newunicodechar{Γ}{\ensuremath{\Gamma}}
\newunicodechar{Δ}{\ensuremath{\Delta}}
\newunicodechar{Ε}{\ensuremath{\varepsilon}}
\newunicodechar{Θ}{\ensuremath{\Theta}}
\newunicodechar{Λ}{\ensuremath{\Lambda}}
\newunicodechar{Μ}{\ensuremath{\mu}}
\newunicodechar{Τ}{\ensuremath{\tau}}
\newunicodechar{Φ}{\ensuremath{\Phi}}
\newunicodechar{Ψ}{\ensuremath{\Psi}}
\newunicodechar{Ω}{\ensuremath{\Omega}}
\newunicodechar{↦}{\ensuremath{\mapsto}}
\newunicodechar{∈}{\ensuremath{\in}}
\newunicodechar{∉}{\ensuremath{\notin}}
\newunicodechar{∧}{\ensuremath{\wedge}}
\newunicodechar{⇔}{\ensuremath{\Leftrightarrow}}
\newunicodechar{⟺}{\ensuremath{\Longleftrightarrow}}
\newunicodechar{⇒}{\ensuremath{\Rightarrow}}
\newunicodechar{⟹}{\ensuremath{\Longrightarrow}}
\newunicodechar{∘}{\ensuremath{\circ}}
\newunicodechar{∪}{\ensuremath{\cup}}
\newunicodechar{∩}{\ensuremath{\cap}}
\newunicodechar{≡}{\ensuremath{\equiv}}
\newunicodechar{⊂}{\ensuremath{\subset}}
\newunicodechar{⊥}{\ensuremath{\perp}}
\newunicodechar{∃}{\ensuremath{\exists}}
\newunicodechar{∀}{\ensuremath{\forall}}
\newunicodechar{∛}{\ensuremath{\sqrt[3]{\,}}}
\newunicodechar{∜}{\ensuremath{\sqrt[4]{\,}}}
\newunicodechar{⃗}{\ensuremath{{}^{\to}}}
\newunicodechar{ⁿ}{\ifmmode^{n}\else\textsuperscript{\ensuremath{n}}\fi}
\newunicodechar{ˣ}{\ifmmode^{x}\else\textsuperscript{\ensuremath{x}}\fi}
\newunicodechar{ᵇ}{\ifmmode^{b}\else\textsuperscript{\ensuremath{b}}\fi}
\newunicodechar{ʳ}{\ifmmode^{r}\else\textsuperscript{\ensuremath{r}}\fi}
\newunicodechar{ᵘ}{\ifmmode^{u}\else\textsuperscript{\ensuremath{u}}\fi}
\newunicodechar{ʸ}{\ifmmode^{y}\else\textsuperscript{\ensuremath{y}}\fi}
\newunicodechar{ᵃ}{\ifmmode^{a}\else\textsuperscript{\ensuremath{a}}\fi}
\newunicodechar{ᵉ}{\ifmmode^{e}\else\textsuperscript{\ensuremath{e}}\fi}
\newunicodechar{ᵏ}{\ifmmode^{k}\else\textsuperscript{\ensuremath{k}}\fi}
\newunicodechar{ᵖ}{\ifmmode^{p}\else\textsuperscript{\ensuremath{p}}\fi}
\newunicodechar{ᴺ}{\ifmmode^{N}\else\textsuperscript{\ensuremath{N}}\fi}
\newunicodechar{ᶜ}{\ifmmode^{c}\else\textsuperscript{\ensuremath{c}}\fi}
\newunicodechar{ʰ}{\ifmmode^{h}\else\textsuperscript{\ensuremath{h}}\fi}
\newunicodechar{ᵠ}{\ifmmode^{q}\else\textsuperscript{\ensuremath{q}}\fi}
\newunicodechar{ₙ}{\ifmmode_{n}\else\textsubscript{\ensuremath{n}}\fi}
\newunicodechar{ₐ}{\ifmmode_{a}\else\textsubscript{\ensuremath{a}}\fi}
\newunicodechar{ᵢ}{\ifmmode_{i}\else\textsubscript{\ensuremath{i}}\fi}
\newunicodechar{ᵣ}{\ifmmode_{r}\else\textsubscript{\ensuremath{r}}\fi}
\newunicodechar{ₖ}{\ifmmode_{k}\else\textsubscript{\ensuremath{k}}\fi}
\newunicodechar{ᵦ}{\ifmmode_{\beta}\else\textsubscript{\ensuremath{\beta}}\fi}
\newunicodechar{ₓ}{\ifmmode_{x}\else\textsubscript{\ensuremath{x}}\fi}
\newfontfamily\popdisplay{ArchivoBlack-Regular.ttf}[Path=../../../fonts/]
\newfontfamily\poplogo{SpaceGrotesk-Bold.ttf}[Path=../../../fonts/]
\newfontfamily\popsemi{PlusJakartaSans-SemiBold.ttf}[Path=../../../fonts/]
\newfontfamily\popxbold{PlusJakartaSans-ExtraBold.ttf}[Path=../../../fonts/]
\newfontfamily\popxboldls{PlusJakartaSans-ExtraBold.ttf}[Path=../../../fonts/,LetterSpace=3.0]

\setlist[enumerate]{leftmargin=1.7em,itemsep=5pt,topsep=4pt}
\setlist[itemize]{leftmargin=1.7em,itemsep=4pt,topsep=4pt,label={\color{poporange}\textbullet}}
\setlength{\parindent}{0pt}
\setlength{\parskip}{5pt}
\renewcommand{\arraystretch}{1.25}

% pgfplots : style Pop par défaut
\pgfplotsset{
  every axis/.append style={axis line style={popink,line width=1pt},tick style={popink},
    grid style={popline,line width=0.5pt},label style={font=\small},tick label style={font=\footnotesize}},
  popcurve/.style={poporange,line width=1.8pt,no markers,smooth},
}
% Même style disponible dans un \draw TikZ simple (hors pgfplots)
\tikzset{popcurve/.style={poporange,line width=1.8pt}}
\tikzset{popgrid/.style={popline,very thin}}
\colorlet{popcurve}{poporange} % le nom du style est parfois utilisé comme couleur

% ============================================================
% Pill / squiggle
% ============================================================
\newcommand{\poppill}[3]{%
  \tikz[baseline=-0.55ex]\node[draw=popink,line width=1.2pt,rounded corners=2.6pt,%
    fill=#2,inner xsep=6.5pt,inner ysep=3pt]{%
    \popxboldls\fontsize{7}{7}\selectfont\color{#3}\MakeUppercase{#1}};%
}
\newcommand{\popsquiggle}[1]{%
  \tikz[baseline=-0.4ex]{
    \draw[#1,line width=2.6pt,line cap=round]
      plot [smooth] coordinates {(0,0.09) (0.34,0.01) (0.68,0.21) (1.02,0.01) (1.36,0.10)};
  }%
}
% Mot-clé surligné (marqueur orange sous le texte)
\newcommand{\cle}[1]{{\popxbold\color{popdark}#1}}

% ============================================================
% En-tête / pied
% ============================================================
\pagestyle{fancy}
\fancyhf{}
\renewcommand{\headrulewidth}{0pt}
\renewcommand{\footrulewidth}{0pt}
\lhead{\raisebox{-0.15ex}{\poplogo\fontsize{13}{13}\selectfont\color{popink}OnBuch\color{poporange}+}}
\rhead{\poppill{\DOCMATIERE\ \textbullet\ \DOCNIVEAU\ \textbullet\ Chapter \DOCLECON}{white}{popink}}
\lfoot{\popxbold\fontsize{7.6}{7.6}\selectfont\color{popmuted}\MakeUppercase{Signed course OnBuch+}\ \textbullet\ {\popsemi\DOCTITRE}}
\rfoot{\tikz[baseline=-0.6ex]\node[rounded corners=8.5pt,fill=popink,minimum height=17pt,minimum width=17pt,inner xsep=5pt,inner sep=0]{\popxbold\fontsize{8}{8}\selectfont\color{popcream}\thepage\,/\,\pageref*{LastPage}};}

% ============================================================
% Titres
% ============================================================
\newcommand{\chaptitre}[2]{%
  \begin{tcolorbox}[enhanced,colback=poporange,colframe=popink,boxrule=2.6pt,arc=11pt,
      drop shadow=popink,left=15pt,right=15pt,top=12pt,bottom=11pt,before skip=3pt,after skip=18pt]
    \poppill{#2}{popink}{popcream}\par\vspace{9pt}
    {\raggedright\hyphenpenalty=10000\exhyphenpenalty=10000\popdisplay\fontsize{22}{24}\selectfont\color{popcream}\MakeUppercase{#1}\par}
  \end{tcolorbox}
}
% Section numérotée : pastille orange avec le numéro + titre Archivo
\newcounter{popsec}
\newcommand{\coursec}[1]{%
  \stepcounter{popsec}\setcounter{popsub}{0}%
  \par\vspace{16pt}\noindent
  \begin{minipage}{\linewidth}
  \tikz[baseline=-0.7ex]\node[draw=popink,line width=1.6pt,fill=poporange,rounded corners=4pt,minimum size=22pt,inner sep=0]{\popdisplay\fontsize{12}{12}\selectfont\color{popcream}\thepopsec};%
  \hspace{8pt}{\popdisplay\fontsize{15}{17}\selectfont\color{popink}\MakeUppercase{#1}}\par
  \vspace{4pt}\noindent{\color{poporange}\rule{36pt}{3.6pt}}
  \end{minipage}\par\nopagebreak\vspace{8pt}
}
\newcounter{popsub}
\newcommand{\courssub}[1]{%
  \stepcounter{popsub}%
  \par\vspace{9pt}\noindent{\popxbold\fontsize{11.5}{13}\selectfont\color{popink}{\color{poporange}\thepopsec.\thepopsub}\hspace{6pt}#1}\par\nopagebreak\vspace{4pt}
}

% ============================================================
% Boîtes pédagogiques — même recette : fond blanc, bordure épaisse colorée,
% ombre dure encre, badge plein. Titre optionnel [#1].
% ============================================================
\tcbset{popbase/.style={breakable,enhanced,colback=white,boxrule=2.3pt,arc=9pt,
  shadow={3pt}{-3pt}{0pt}{popink},left=14pt,right=14pt,top=11pt,bottom=11pt,before skip=11pt,after skip=15pt,
  fontupper=\popsemi\fontsize{10.6}{14.5}\selectfont\color{popink},
  fontlower=\popsemi\fontsize{10.6}{14.5}\selectfont\color{popink}}}
% Coche dessinée (le glyphe ✓ n'existe pas dans les polices du document)
\newcommand{\popcheck}[1]{\tikz[baseline=-0.5ex]\draw[#1,line width=1.6pt,line cap=round,line join=round](-0.13,0.0)--(-0.04,-0.09)--(0.13,0.1);}
\providecommand{\coche}{\popcheck{popgreen}}
\providecommand{\resultat}[1]{\par\smallskip\noindent\fcolorbox{popgreen}{popgreenL}{\parbox{\dimexpr\linewidth-2\fboxsep-2\fboxrule\relax}{\textbf{Result.} #1}}\par\smallskip}
\newcommand{\popboxhead}[3]{\poppill{#1}{#2}{white}\ifx\relax#3\relax\else\hspace{6pt}{\popxbold\fontsize{10.6}{12}\selectfont\color{popink}#3}\fi\par\vspace{7pt}}

\newtcolorbox{aretenir}[1][]{popbase,colframe=popgreen,before upper={\popboxhead{Key points}{popgreen}{#1}}}
\newtcolorbox{exemplebox}[1][]{popbase,colframe=poporange,before upper={\popboxhead{Exemple}{poporange}{#1}}}
\newtcolorbox{definition}[1][]{popbase,colframe=popblue,before upper={\popboxhead{Definition}{popblue}{#1}}}
\newtcolorbox{propriete}[1][]{popbase,colframe=poppurple,before upper={\popboxhead{Property}{poppurple}{#1}}}
\newtcolorbox{methode}[1][]{popbase,colframe=popgold,before upper={\popboxhead{Method}{popgold}{#1}}}
\newtcolorbox{attention}[1][]{popbase,colframe=poppink,before upper={\popboxhead{Warning}{poppink}{#1}}}
\newtcolorbox{experience}[1][]{popbase,colframe=popblue,colback=popblueL,before upper={\popboxhead{Experiment}{popblue}{#1}}}
% « Le savais-tu ? » : carte violette pleine (contexte, culture, Cameroun)
\newtcolorbox{savaistu}[1][]{popbase,colback=poppurple,colframe=popink,
  fontupper=\popsemi\fontsize{10.4}{14.2}\selectfont\color{white},
  before upper={\poppill{Did you know?}{white}{poppurple}\ifx\relax#1\relax\else\hspace{6pt}{\popxbold\color{white}#1}\fi\par\vspace{7pt}}}
% Exercice résolu : énoncé en haut, \tcblower puis la solution sur fond crème
\newcounter{popexo}\newcounter{popexr}
\newtcolorbox{exoresolu}[1][]{popbase,colframe=poporange,colbacklower=popcream,
  segmentation style={popink,line width=1.2pt,dashed},
  before upper={\stepcounter{popexr}\popboxhead{Worked example \thepopexr}{poporange}{#1}},
  before lower={\poppill{Solution}{popink}{popcream}\par\vspace{6pt}}}
% Exercice (non résolu en place) — la correction est dans la partie Corrigés
\newtcolorbox{exercice}[1][]{popbase,colframe=popink,
  before upper={\stepcounter{popexo}\popboxhead{Exercise \thepopexo}{popink}{#1}}}
% Corrigé
\newtcolorbox{corrige}[1][]{popbase,colframe=popgreen,colback=popgreenL,
  before upper={\popboxhead{Solution}{popgreen}{#1}}}
% Objectifs / prérequis / bilan
\newtcolorbox{objectifs}{popbase,colframe=popink,colback=poporangeL,
  before upper={\popboxhead{Chapter objectives}{popink}{}}}
\newtcolorbox{prerequis}{popbase,colframe=popink,colback=popblueL,
  before upper={\popboxhead{Prerequisites}{popblue}{}}}
\newtcolorbox{bilan}{popbase,colframe=popink,colback=white,boxrule=2.8pt,
  before upper={\popboxhead{Fiche bilan}{poporange}{L'essentiel à connaître pour l'examen}}}
% Figure encadrée avec légende
\newtcolorbox{popfigure}[1][]{enhanced,breakable=false,colback=white,colframe=popline,boxrule=1.4pt,arc=8pt,
  left=10pt,right=10pt,top=10pt,bottom=8pt,before skip=10pt,after skip=12pt,halign=center,
  lower separated=false,halign lower=center,
  fontlower=\popsemi\fontsize{9}{11.5}\selectfont\color{popmuted},
  #1}
\newcommand{\legende}[1]{\par\vspace{6pt}{\popsemi\fontsize{9}{11.5}\selectfont\color{popmuted}#1}}

% ============================================================
% Signature OnBuch+
% ============================================================
\newcommand{\poplogoinline}[1]{{\poplogo\fontsize{#1}{#1}\selectfont\color{popink}OnBuch\color{poporange}+}}
\newcommand{\signatureonbuch}{%
  \begin{tcolorbox}[enhanced,colback=popink,colframe=popink,boxrule=2.2pt,arc=10pt,
      drop shadow=poporange,left=14pt,right=14pt,top=10pt,bottom=10pt,before skip=0pt,after skip=0pt]
    \tikz[baseline=-0.7ex]\node[circle,fill=popgreen,draw=popcream,line width=1.3pt,minimum size=22pt,inner sep=0]{\popcheck{white}};%
    \hspace{9pt}\begin{minipage}[c]{0.62\linewidth}
      {\popxbold\fontsize{12}{13}\selectfont\color{popcream}Signed\ \textbullet\ {\poplogo OnBuch\color{poporange}+}}\par\vspace{2pt}
      {\raggedright\popsemi\fontsize{8}{10.5}\selectfont\color{popline}\MakeUppercase{OnBuch+ teaching team}\\\MakeUppercase{Follows the official MINESEC syllabus}\par}
    \end{minipage}\hfill
    {\poplogo\fontsize{22}{22}\selectfont\color{popcream}OnBuch\color{poporange}+}
  \end{tcolorbox}%
}

% ============================================================
% Page de garde
% #1 titre · #2 matière · #3 niveau · #4 module · #5 n° leçon · #6 durée ·
% #7 description / objectifs (texte ou itemize)
% ============================================================
\newcommand{\pagegarde}[7]{%
  \thispagestyle{empty}
  \newgeometry{a4paper,margin=1.7cm,top=1.6cm,bottom=1.4cm}%
  \begin{tikzpicture}[remember picture,overlay]
    \fill[poporange!18] ([xshift=-3.2cm,yshift=-3.6cm]current page.north east) circle (4.6cm);
    \fill[poppurple!14] ([xshift=2.4cm,yshift=7.6cm]current page.south west) circle (3.8cm);
  \end{tikzpicture}%
  {\poplogo\fontsize{30}{30}\selectfont\color{popink}OnBuch\color{poporange}+}\hfill
  \raisebox{0.6ex}{\poppill{Signed course \textbullet\ Premium}{popink}{popcream}}\par
  \vspace{1pt}\popsquiggle{poppurple}\par
  \vspace{8pt}
  \begin{tcolorbox}[enhanced,colback=poporange,colframe=popink,boxrule=2.8pt,arc=13pt,
      drop shadow=popink,left=18pt,right=18pt,top=12pt,bottom=13pt,after skip=10pt]
    \poppill{#2 \textbullet\ #3}{popink}{popcream}\hspace{5pt}\poppill{#4}{white}{popink}\par\vspace{8pt}
    {\popdisplay\fontsize{14}{14}\selectfont\color{popink}CHAPTER #5}\par\vspace{4pt}
    {\raggedright\hyphenpenalty=10000\exhyphenpenalty=10000\popdisplay\fontsize{25}{27}\selectfont\color{popcream}\MakeUppercase{#1}\par}
    \vspace{8pt}
    {\popxbold\fontsize{9.5}{9.5}\selectfont\color{popink}\MakeUppercase{Suggested time: #6}}
  \end{tcolorbox}
  \begin{tcolorbox}[enhanced,colback=white,colframe=popink,boxrule=2.2pt,arc=10pt,
      drop shadow=popink,left=16pt,right=16pt,top=10pt,bottom=10pt,after skip=8pt]
    \poppill{In this chapter}{poppurple}{white}\par\vspace{5pt}
    \popsemi\fontsize{9.3}{12.6}\selectfont\color{popink}#7
  \end{tcolorbox}
  \noindent\begin{minipage}[t]{0.487\linewidth}
    \begin{tcolorbox}[enhanced,colback=popgreen,colframe=popink,boxrule=2.2pt,arc=10pt,
        drop shadow=popink,left=11pt,right=11pt,top=10pt,bottom=10pt,height=3.1cm,valign=center]
      \tcbox[colback=white,colframe=popink,boxrule=1.3pt,arc=5pt,boxsep=0pt,nobeforeafter,
        left=3pt,right=3pt,top=3pt,bottom=3pt,valign=center]{\qrcode[height=1.9cm]{https://play.google.com/store/apps/details?id=com.luvvix.onbuch&hl=en}}%
      \hspace{8pt}\begin{minipage}[c]{0.5\linewidth}
        {\popdisplay\fontsize{11}{12.5}\selectfont\color{white}\MakeUppercase{Download OnBuch}}\par\vspace{4pt}
        {\raggedright\popsemi\fontsize{8.4}{11}\selectfont\color{white}Free \textbullet\ Google Play\\AI tutor, past papers, results\par}
      \end{minipage}
    \end{tcolorbox}
  \end{minipage}\hfill
  \begin{minipage}[t]{0.487\linewidth}
    \begin{tcolorbox}[enhanced,colback=poppurple,colframe=popink,boxrule=2.2pt,arc=10pt,
        drop shadow=popink,left=11pt,right=11pt,top=10pt,bottom=10pt,height=3.1cm,valign=center]
      \tikz[baseline=-0.7ex]\node[circle,fill=white,draw=popink,line width=1.3pt,minimum size=1.6cm,inner sep=0]{\popdisplay\fontsize{24}{24}\selectfont\color{poppurple}+};%
      \hspace{8pt}\begin{minipage}[c]{0.6\linewidth}
        {\popdisplay\fontsize{11}{12.5}\selectfont\color{white}\MakeUppercase{Go OnBuch+}}\par\vspace{4pt}
        {\raggedright\popsemi\fontsize{8.4}{11}\selectfont\color{white}All signed courses, unlimited AI tutor, PDF sheets — OnBuch+ tab in the app\par}
      \end{minipage}
    \end{tcolorbox}
  \end{minipage}
  \vspace{8pt}
  \popcontenu
  \vfill
  \signatureonbuch
  \restoregeometry
  \newpage
}

% Rangée « Ce cours contient » (page de garde)
\newcommand{\poptile}[3]{% #1 couleur #2 gros texte #3 légende
  \begin{tcolorbox}[enhanced,colback=#1,colframe=popink,boxrule=2pt,arc=9pt,shadow={2.5pt}{-2.5pt}{0pt}{popink},
      left=6pt,right=6pt,top=9pt,bottom=9pt,halign=center,valign=center,height=1.75cm,width=\linewidth,nobeforeafter]
    {\popdisplay\fontsize{12.5}{13}\selectfont\color{white}\MakeUppercase{#2}}\par\vspace{3pt}
    {\centering\popsemi\fontsize{7.8}{9.5}\selectfont\color{white}#3\par}
  \end{tcolorbox}}
\newcommand{\popcontenu}{%
  \par\noindent{\popxboldls\fontsize{7.5}{7.5}\selectfont\color{popmuted}THIS SIGNED COURSE CONTAINS}\par\vspace{7pt}
  \noindent\begin{minipage}[t]{0.235\linewidth}\poptile{popblue}{Course}{complete, illustrated, step by step}\end{minipage}\hfill
  \begin{minipage}[t]{0.235\linewidth}\poptile{poporange}{Methods}{and worked examples}\end{minipage}\hfill
  \begin{minipage}[t]{0.235\linewidth}\poptile{poppink}{Exercises}{exam style + solutions}\end{minipage}\hfill
  \begin{minipage}[t]{0.235\linewidth}\poptile{popgreen}{Summary sheet}{the essentials to remember}\end{minipage}\par}

% Sommaire
\newtcolorbox{sommaire}{enhanced,colback=white,colframe=popink,boxrule=2.2pt,arc=10pt,
  drop shadow=popink,left=18pt,right=18pt,top=13pt,bottom=13pt,before skip=6pt,after skip=20pt,
  before upper={\poppill{Contents}{poppurple}{white}\par\vspace{9pt}\popsemi\fontsize{10.6}{14.5}\selectfont\color{popink}}}

% Signature de fin de cours — poussée tout en bas de la dernière page
\newcommand{\findecours}{%
  \par\vfill\nopagebreak
  \begin{center}\popsquiggle{poporange}\end{center}\vspace{4pt}
  \signatureonbuch
}
\AtBeginDocument{\def\llbracket{\mathopen{[\mkern-3.5mu[}}\def\rrbracket{\mathclose{]\mkern-3.5mu]}}} % intervalles d'entiers (glyphe ⟦ absent de la police)
% Glyphes absents de Fira Math : redéfinis à partir de glyphes disponibles
\AtBeginDocument{%
  \def\setminus{\mathbin{\backslash}}\let\smallsetminus\setminus
  \def\vdots{\mathord{\vbox{\baselineskip3.2pt\lineskiplimit0pt\kern4pt\hbox{.}\hbox{.}\hbox{.}}}}%
  \def\ddots{\mathinner{\mkern1mu\raise6.5pt\hbox{.}\mkern2mu\raise3.6pt\hbox{.}\mkern2mu\raise0.7pt\hbox{.}\mkern1mu}}%
  \def\triangle{\mathord{\scalebox{0.9}{$\symup{\Delta}$}}}%
  \def\nabla{\mathord{\raisebox{\depth}{\scalebox{1}[-1]{$\symup{\Delta}$}}}}%
  \def\checkmark{\popcheck{popdark}}%
  \def\star{\mathbin{\ast}}\def\bigstar{\mathord{\ast}}%
  \def\diamond{\mathbin{\scalebox{0.6}{$\Diamond$}}}%
  \def\bigcup{\mathop{\mathchoice{\raisebox{-0.3em}{\scalebox{1.7}{$\cup$}}}{\scalebox{1.25}{$\cup$}}{\cup}{\cup}}\displaylimits}%
  \def\bigcap{\mathop{\mathchoice{\raisebox{-0.3em}{\scalebox{1.7}{$\cap$}}}{\scalebox{1.25}{$\cap$}}{\cap}{\cap}}\displaylimits}%
  \def\lesseqqgtr{\mathrel{\lessgtr}}\def\bigoplus{\mathop{\oplus}\displaylimits}%
}
\providecommand{\ssi}{if and only if} % abréviation fréquente
\AtBeginDocument{\providecommand{\wideparen}{\overparen}} % arc de cercle

% Fonctions usuelles que les modèles utilisent sans que amsmath les fournisse
\providecommand{\arsinh}{\operatorname{arsinh}}\providecommand{\arcosh}{\operatorname{arcosh}}
\providecommand{\artanh}{\operatorname{artanh}}\providecommand{\arsech}{\operatorname{arsech}}
\providecommand{\arcsch}{\operatorname{arcsch}}\providecommand{\arcoth}{\operatorname{arcoth}}
\providecommand{\arcsinh}{\operatorname{arsinh}}\providecommand{\arccosh}{\operatorname{arcosh}}
\providecommand{\arctanh}{\operatorname{artanh}}\providecommand{\sech}{\operatorname{sech}}
\providecommand{\csch}{\operatorname{csch}}\providecommand{\arcsec}{\operatorname{arcsec}}
\providecommand{\arccsc}{\operatorname{arccsc}}\providecommand{\arccot}{\operatorname{arccot}}
\providecommand{\sgn}{\operatorname{sgn}}\providecommand{\lcm}{\operatorname{lcm}}
\providecommand{\Var}{\operatorname{Var}}\providecommand{\Cov}{\operatorname{Cov}}

%% FIGFIX-BEGIN v5 — correctifs globaux des figures (docs/onbuch-generation/tools/figfix)
\usepackage{adjustbox}\usepackage{etoolbox}\tcbuselibrary{raster}
% toute figure TikZ/pgfplots plus large que la ligne est réduite à la largeur disponible
\BeforeBeginEnvironment{tikzpicture}{\begin{adjustbox}{max width=\linewidth}}
\AfterEndEnvironment{tikzpicture}{\end{adjustbox}}
% légendes pgfplots lisibles par-dessus les courbes
\pgfplotsset{every axis/.append style={legend style={fill=white,fill opacity=0.92,text opacity=1,draw=black!15,inner sep=3pt}}}
% étiquettes d'arêtes (syntaxe edge["..."]) avec fond blanc
\tikzset{every edge quotes/.append style={fill=white,inner sep=1.2pt}}
%% FIGFIX-END
\begin{document}

\pagegarde{Sampling and Estimation}{Pure Mathematics with Statistics}{Upper Sixth}{Upper Sixth — Pure mathematics with statistics}{15}{12 periods}{This chapter equips students with the theory and practice of sampling and statistical estimation, enabling them to draw valid inferences about populations from samples — a core skill for the GCE Advanced Level Pure Mathematics with Statistics examination.
\vspace{4pt}
\begin{multicols}{2}\small
\begin{itemize}
\item Distinguish between population, sample, parameter and statistic and identify common sampling designs.
\item Explain the advantages and limitations of sampling compared with a census.
\item Define sampling error and differentiate it from non‑sampling error.
\item Construct and interpret point estimators for population mean and proportion.
\item State and apply the criteria for a good estimator (unbiasedness, consistency, efficiency, sufficiency).
\item Derive and interpret confidence intervals for a population mean (known and unknown variance) and for a population proportion.
\end{itemize}\end{multicols}}

\begin{sommaire}
\begin{multicols}{2}
\begin{enumerate}
\item 1. Foundations of Sampling
\item 2. Sampling Errors and Point Estimation
\item 3. Properties of Good Estimators
\item 4. Confidence Intervals for the Mean (Known Variance)
\item 5. Confidence Intervals for the Mean (Unknown Variance) and for a Proportion
\item 6. Margin of Error, Sample Size Determination and Applications
\item Integration activity
\item Methods and worked examples
\item Exercises
\item Solutions to the exercises
\item Summary sheet
\end{enumerate}
\end{multicols}
\end{sommaire}

\begin{objectifs}
\begin{itemize}
\item Distinguish between population, sample, parameter and statistic and identify common sampling designs.
\item Explain the advantages and limitations of sampling compared with a census.
\item Define sampling error and differentiate it from non‑sampling error.
\item Construct and interpret point estimators for population mean and proportion.
\item State and apply the criteria for a good estimator (unbiasedness, consistency, efficiency, sufficiency).
\item Derive and interpret confidence intervals for a population mean (known and unknown variance) and for a population proportion.
\item Calculate the margin of error and determine the required sample size for a desired precision.
\item Recognise common pitfalls in estimation and apply sampling techniques to real‑life Cameroonian contexts.
\end{itemize}
\end{objectifs}

\begin{prerequis}
\begin{itemize}
\item Probability distributions (binomial, normal) from Lower Sixth.
\item Basic descriptive statistics: mean, variance, standard deviation.
\item Concept of a random variable and its expectation.
\item Standard normal (Z) and Student’s t tables.
\item Algebraic manipulation of inequalities and square roots.
\end{itemize}
\end{prerequis}

\newpage

\coursec{Foundations of Sampling}

Imagine a cocoa cooperative in the South-West region of Cameroon that wants to estimate the average bean weight per pod for the new harvest season. There are tens of thousands of pods spread over many farms. Weighing every single pod would take weeks, cost a fortune in labour, and by the time the last pod was weighed the first results would already be out of date. Yet the cooperative must make decisions \emph{now}: what price to offer farmers, whether the quality meets export standards, how much drying capacity to prepare. The way out is elegant: select a \emph{well-designed sample} of pods, weigh only those, and use the result to draw a reliable conclusion about the whole harvest. This idea — learning about a large whole by carefully studying a small part — is the heart of \cle{sampling}, and it raises the central problem of this chapter: \textbf{how do we choose the part so that what we learn from it can be trusted?}

\courssub{Population, Sample, Parameter and Statistic}

Before any sampling can be done, we must be precise about \emph{what} we are studying and \emph{what} we are measuring. The vocabulary below is used throughout the GCE Advanced Level statistics papers, so learn it exactly.

\begin{definition}[Population, Sample, Parameter, Statistic]
\begin{itemize}
\item A \cle{population} (or \cle{target population}) is the complete collection of all individuals, items or measurements about which we want information. It may be finite (all cocoa pods of the cooperative this season) or effectively infinite (all bottles a factory will ever fill).
\item A \cle{sample} is a subset of the population that is actually observed or measured. The number of units in the sample is the \cle{sample size}, denoted $n$.
\item A \cle{parameter} is a numerical characteristic of the \emph{population}, e.g.\ the population mean $\mu$, the population variance $\sigma^2$, or the population proportion $p$. Parameters are usually \emph{unknown} constants and are what we wish to estimate.
\item A \cle{statistic} is a numerical characteristic computed from the \emph{sample}, e.g.\ the sample mean $\bar{x}$, the sample variance $s^2$, or the sample proportion $\hat{p}$. Statistics are \emph{known} once the sample is taken; they are random variables before the sample is drawn, and are used to estimate parameters.
\end{itemize}
\end{definition}

The fundamental relationship is:
\[
\underbrace{\bar{x},\ s^2,\ \hat{p}}_{\text{statistics (known, from sample)}} \quad \xrightarrow{\ \text{estimation}\ } \quad \underbrace{\mu,\ \sigma^2,\ p}_{\text{parameters (unknown, of population)}}
\]

\begin{exemplebox}[Identifying the four notions]
A school inspector in Bamenda wants to know the mean score of all Upper Sixth candidates in the North-West region in a mock mathematics examination. She selects 120 candidates and finds their mean score to be $54.2$ with standard deviation $11.3$.
\begin{itemize}
\item \textbf{Population:} all Upper Sixth candidates in the region.
\item \textbf{Sample:} the 120 selected candidates.
\item \textbf{Parameter:} the true mean score $\mu$ of all candidates (unknown).
\item \textbf{Statistic:} the sample mean $\bar{x} = 54.2$ (known).
\end{itemize}
\end{exemplebox}

\begin{attention}[Do not confuse parameter and statistic]
A very common examination error is to call $\bar{x}$ a "parameter". Remember: parameters belong to the \emph{population} and are denoted by Greek letters ($\mu$, $\sigma$, $p$); statistics belong to the \emph{sample} and are denoted by Latin letters ($\bar{x}$, $s$, $\hat{p}$).
\end{attention}

\courssub{Probability Sampling Methods}

A sampling method is called a \cle{probability sampling} method when chance enters the selection in a controlled, measurable way.

\begin{propriete}[Defining property of probability sampling]
In probability sampling, every unit of the population has a \emph{known, non-zero probability} of being selected. This property is what allows us to quantify the accuracy of estimates (through confidence intervals, studied later in this chapter) and to avoid systematic selection bias. (Statement examinable; no formal proof required.)
\end{propriete}

\textbf{Simple random sampling (SRS).} Every possible sample of size $n$ from a population of size $N$ has the same probability of being chosen. Equivalently, every unit has the same probability $n/N$ of being included, and selections are independent in the sense that the inclusion of one unit does not affect the probabilities for others (when sampling without replacement, the probabilities change slightly but the design remains SRS). This is the reference method against which all others are judged.

\textbf{Systematic sampling.} Order the population in a list (the \cle{sampling frame}), compute the sampling interval $k = N/n$ (usually rounded to an integer), choose a random start $r$ between $1$ and $k$, then select units $r,\ r+k,\ r+2k,\ \dots$ It is quick and spreads the sample evenly, but is dangerous if the list has a hidden periodic pattern coinciding with $k$ (e.g.\ a list of houses where every $k$-th house is a corner shop).

\textbf{Stratified sampling.} Divide the population into homogeneous subgroups called \cle{strata} (e.g.\ farms by district, students by class, pods by variety), then draw a simple random sample \emph{independently within each stratum}. With \cle{proportional allocation}, the number taken from stratum $i$ is $n_i = n \times N_i/N$, where $N_i$ is the stratum size. Stratification guarantees that every subgroup is represented and usually increases precision (reduces sampling variance) compared to SRS of the same total size.

\textbf{Cluster sampling.} Divide the population into heterogeneous groups called \cle{clusters} (often geographical, e.g.\ villages, schools, or enumeration areas), choose a random sample of \emph{whole clusters}, and observe \emph{every unit} inside the chosen clusters (one-stage cluster sampling). It is cheap (less travel) but less precise than stratified sampling, because units within a cluster tend to resemble each other (positive intra-cluster correlation), so the effective sample size is smaller than $n$.

\textbf{Multi-stage sampling.} Sampling is done in stages: e.g.\ first select divisions (primary sampling units), then schools within chosen divisions (secondary units), then students within chosen schools (ultimate units). Large national surveys in Cameroon (household surveys, school assessments, crop estimation) are typically multi-stage designs, often combining stratification at the first stage with clustering at later stages.

\begin{popfigure}
\begin{center}
\begin{tikzpicture}[scale=0.92, every node/.style={font=\small}]
% Population box
\draw[thick, popblue, rounded corners] (0,0) rectangle (6.4,4.4);
\node[popblue, anchor=west] at (0.1,4.15) {\textbf{Population} ($N$ pods)};
% Strata
\draw[fill=popblueL, draw=popblue] (0.3,2.9) rectangle (6.1,3.9);
\draw[fill=popgreenL, draw=popgreen] (0.3,1.6) rectangle (6.1,2.7);
\draw[fill=white, draw=poporange] (0.3,0.3) rectangle (6.1,1.4);
\node at (3.2,3.4) {Stratum 1: District A ($N_1$)};
\node at (3.2,2.15) {Stratum 2: District B ($N_2$)};
\node at (3.2,0.85) {Stratum 3: District C ($N_3$)};
% Sampled units
\fill[popdark] (1.0,3.15) circle (0.06); \fill[popdark] (2.2,3.7) circle (0.06); \fill[popdark] (4.5,3.2) circle (0.06);
\fill[popdark] (1.5,1.85) circle (0.06); \fill[popdark] (3.8,2.5) circle (0.06);
\fill[popdark] (2.6,0.6) circle (0.06); \fill[popdark] (5.0,1.1) circle (0.06);
% Arrows to sample
\draw[-latex, thick, poporange] (6.4,3.4) -- (7.6,3.0);
\draw[-latex, thick, poporange] (6.4,2.15) -- (7.6,2.15);
\draw[-latex, thick, poporange] (6.4,0.85) -- (7.6,1.3);
% Sample box
\draw[thick, poporange, rounded corners] (7.7,1.0) rectangle (10.6,3.3);
\node[poporange] at (9.15,3.0) {\textbf{Sample} ($n$)};
\node[align=center] at (9.15,2.45) {$n_1$ from A\\ $n_2$ from B\\ $n_3$ from C};
\node[align=center, popmuted] at (9.15,1.35) {random within\\ each stratum};
\end{tikzpicture}
\end{center}
\legende{Stratified random sampling: the population is split into homogeneous strata and a random sample is drawn independently within each stratum.}
\end{popfigure}

\begin{methode}[Drawing a simple random sample]
\begin{enumerate}
\item Obtain a \cle{sampling frame}: a numbered list of all $N$ units of the population, from $1$ (or $000$) to $N$. The frame must be complete, up-to-date, and free of duplicates.
\item Fix the sample size $n$.
\item Use a random number table, a calculator's \texttt{Rand} function, or statistical software to generate $n$ distinct random integers between $1$ and $N$. Read digits in groups matching the number of digits of $N$; ignore repeats and numbers exceeding $N$.
\item Select the units whose numbers were drawn. Measure them and record the data.
\end{enumerate}
\end{methode}

\begin{exoresolu}[Systematic and stratified sampling in practice]
A cooperative has $N = 2400$ pods registered, coming from three districts: District A ($1200$ pods), District B ($800$ pods), District C ($400$ pods). A sample of $n = 120$ pods is required.
\textbf{(a)} Describe a systematic sample. \textbf{(b)} Give the stratum sizes for a stratified sample with proportional allocation.
\tcblower
\textbf{(a) Systematic sampling.} The sampling interval is
\[
k = \frac{N}{n} = \frac{2400}{120} = 20.
\]
Choose a random start $r$ between $1$ and $20$, say $r = 7$. Then select pods numbered
\[
7,\ 27,\ 47,\ 67,\ \dots,\ 2387,
\]
which gives exactly $120$ pods spread evenly through the register.

\textbf{(b) Stratified sampling with proportional allocation.} Each stratum contributes a share proportional to its size:
\[
n_1 = 120 \times \frac{1200}{2400} = 60, \qquad
n_2 = 120 \times \frac{800}{2400} = 40, \qquad
n_3 = 120 \times \frac{400}{2400} = 20.
\]
So we draw simple random samples of $60$, $40$ and $20$ pods from Districts A, B and C respectively. Check: $60+40+20 = 120 = n$. \hfill$\blacksquare$
\end{exoresolu}

\courssub{Non-Probability Sampling Methods}

In \cle{non-probability sampling}, units are selected by judgement or convenience, so selection probabilities are unknown or zero for some units.

\begin{itemize}
\item \textbf{Convenience sampling:} take whoever is easiest to reach (e.g.\ interviewing shoppers at one market entrance). Fast and cheap, but rarely representative.
\item \textbf{Quota sampling:} the interviewer must fill fixed quotas (e.g.\ 30 men, 30 women, 20 youths) but chooses the individuals freely within each quota. Resembles stratified sampling but without random selection within quotas.
\item \textbf{Purposive (judgement) sampling:} the researcher hand-picks units believed to be typical or especially informative (e.g.\ choosing "model" farms for an agronomy pilot study).
\item \textbf{Snowball sampling:} respondents recruit other respondents; useful for hidden or hard-to-reach groups (e.g.\ informal-sector workers, drug users).
\end{itemize}

These methods are used when a sampling frame does not exist, when time or money is extremely limited, or in exploratory/pilot studies. Their weakness is that \emph{the accuracy of the estimates cannot be quantified}, and bias can go undetected.

\begin{attention}[The trap of convenience sampling]
Convenience sampling often leads to \cle{selection bias} and non-representative results. If the cocoa cooperative weighs only pods from trees near the road (easy to reach), and roadside trees happen to be better watered, the sample mean will overestimate the true mean bean weight — and no amount of calculation on that biased sample can repair the damage. \textbf{A large biased sample is worse than a small random one.}
\end{attention}

\courssub{Sampling Errors and Non-Sampling Errors}

Lesson 186 of the syllabus distinguishes two fundamentally different sources of error. Understanding this distinction is examinable.

\begin{definition}[Sampling Error vs Non-Sampling Error]
\begin{itemize}
\item \cle{Sampling error} is the difference between a sample statistic and the corresponding population parameter that arises \emph{purely by chance} because only a subset of the population is observed. It is an inherent feature of sampling, not a mistake. Its magnitude can be quantified (via the standard error) and reduced by increasing the sample size $n$.
\item \cle{Non-sampling errors} are all other errors that can occur in a survey, whether a sample or a census is used. They include:
      \begin{itemize}
      \item \cle{Selection bias} (coverage error): the sampling frame misses some units or includes units not in the target population.
      \item \cle{Non-response bias}: selected units cannot or will not respond, and non-respondents differ systematically from respondents.
      \item \cle{Measurement error}: poorly worded questions, interviewer effects, faulty instruments, recording mistakes.
      \item \cle{Processing error}: coding, data entry, or tabulation mistakes.
      \end{itemize}
      Non-sampling errors are \emph{not} reduced by increasing $n$; they require better design, training, and quality control.
\end{itemize}
\end{definition}

\begin{aretenir}[Key distinction]
\begin{center}
\begin{tabularx}{\linewidth}{|l|X|X|}
\hline
\rowcolor{popblueL} \textbf{Feature} & \textbf{Sampling Error} & \textbf{Non-Sampling Error} \\
\hline
Source & Random variation in which units are selected & Flaws in design, frame, measurement, processing \\
\hline
Present in census? & No & Yes \\
\hline
Quantifiable? & Yes (standard error, confidence intervals) & Difficult; often unknown magnitude \\
\hline
Reduced by larger $n$? & Yes (roughly $\propto 1/\sqrt{n}$) & No \\
\hline
Examples & Sample mean $\bar{x}$ differs from $\mu$ by chance & Biased frame, leading questions, refusals, data entry typos \\
\hline
\end{tabularx}
\end{center}
\end{aretenir}

\courssub{Advantages and Disadvantages of Sampling}

Why sample at all, instead of a full \cle{census} (observing the whole population)?

\begin{center}
\begin{tabularx}{\linewidth}{|l|X|X|}
\hline
\rowcolor{popblueL} \textbf{Aspect} & \textbf{Advantage of sampling} & \textbf{Disadvantage / risk} \\
\hline
Cost & Far cheaper: only $n$ units instead of $N$ & Poor design wastes the money spent \\
\hline
Speed & Results available quickly for decisions & Rushed fieldwork increases errors \\
\hline
Feasibility & Only option when testing is destructive (e.g.\ crushing pods, testing match lifetimes) or the population is huge & Cannot give information about every small subgroup \\
\hline
Accuracy & Resources concentrated on careful measurement of few units & \cle{Sampling error}: the sample differs from the population by chance \\
\hline
Representativeness & Random selection gives unbiased estimates & \cle{Sampling bias} and \cle{non-response} can distort results \\
\hline
\end{tabularx}
\end{center}

Two disadvantages deserve emphasis. \textbf{Sampling bias} arises when the selection procedure systematically favours some units (a bad frame, convenience selection). \textbf{Non-response} occurs when selected units cannot or will not respond (absent farmers, refusals); if non-respondents differ from respondents, the estimate is biased even with a random design. Good surveys plan call-backs and follow-ups to keep non-response low.

\courssub{Steps in Designing a Sample Survey}

A professional survey — whether by the cooperative, a school, or a national institute — follows a structured process:

\begin{enumerate}
\item \textbf{Define the objectives:} what exactly do we want to estimate (mean bean weight? proportion of export-grade pods?), and to what precision (margin of error, confidence level)?
\item \textbf{Define the target population} precisely (all pods of member farms this season, excluding rejected pods).
\item \textbf{Construct the sampling frame:} the list or map from which units will be drawn. A faulty frame (missing farms, outdated lists, duplicates) is a major source of coverage bias.
\item \textbf{Choose the sampling method} (simple random, systematic, stratified, cluster, multi-stage) suited to the population structure, budget, and required precision.
\item \textbf{Determine the sample size} $n$ (studied in Section 6) balancing precision against cost.
\item \textbf{Design the questionnaire or measurement protocol} and pilot-test it on a small trial sample to detect ambiguities and measurement errors.
\item \textbf{Carry out the fieldwork:} collect the data, monitor quality, manage non-response through call-backs and substitution protocols (if approved).
\item \textbf{Analyse the data:} compute statistics, estimate parameters, and report results with their margins of error and confidence intervals.
\end{enumerate}

\begin{popfigure}
\begin{center}
\begin{tikzpicture}[node distance=0.55cm, every node/.style={font=\small},
box/.style={draw=popblue, fill=popblueL, rounded corners, text width=3.4cm, align=center, minimum height=0.8cm},
arr/.style={-latex, thick, poporange}]
\node[box] (a) {1. Define objectives};
\node[box, below=of a] (b) {2. Define target population};
\node[box, below=of b] (c) {3. Build sampling frame};
\node[box, below=of c] (d) {4. Choose sampling method};
\node[box, right=1.6cm of d] (e) {5. Determine sample size};
\node[box, above=of e] (f) {6. Design \& pilot instrument};
\node[box, above=of f] (g) {7. Fieldwork \& non-response};
\node[box, above=of g] (h) {8. Analysis \& reporting};
\draw[arr] (a) -- (b); \draw[arr] (b) -- (c); \draw[arr] (c) -- (d);
\draw[arr] (d) -- (e); \draw[arr] (e) -- (f); \draw[arr] (f) -- (g); \draw[arr] (g) -- (h);
\end{tikzpicture}
\end{center}
\legende{Flowchart of the sampling design process, from objectives to reporting.}
\end{popfigure}

\begin{savaistu}[Sampling in Cameroon]
Large-scale sample surveys are routinely used in Cameroon: household consumption surveys (ECAM) to monitor living standards, crop estimation surveys before the cocoa and coffee seasons, and school-based assessments of learning (PASEC). Most of these use multi-stage stratified cluster designs — exactly the methods of this lesson — because a complete census of millions of households would be far too slow and costly.
\end{savaistu}

\begin{aretenir}[Key points of Section 1]
\begin{itemize}
\item A \cle{population} is described by \cle{parameters} ($\mu$, $\sigma^2$, $p$); a \cle{sample} yields \cle{statistics} ($\bar{x}$, $s^2$, $\hat{p}$) used to estimate them.
\item \cle{Probability sampling} (simple random, systematic, stratified, cluster, multi-stage) gives every unit a known non-zero selection probability — the basis of trustworthy estimation.
\item \cle{Non-probability sampling} (convenience, quota, purposive, snowball) is sometimes unavoidable but cannot quantify accuracy and risks bias.
\item \cle{Sampling error} is the chance variation due to observing only a sample; it is quantifiable and decreases with $n$. \cle{Non-sampling errors} (bias, non-response, measurement mistakes) are not reduced by larger samples and must be controlled by design.
\item Sampling saves \textbf{cost} and \textbf{time} and is often the only \textbf{feasible} option, but suffers from \textbf{sampling error}, \textbf{bias} and \textbf{non-response}.
\item A good survey follows a disciplined design process: objectives $\to$ population $\to$ frame $\to$ method $\to$ size $\to$ instrument $\to$ fieldwork $\to$ analysis.
\end{itemize}
\end{aretenir}

\begin{exercice}[Choosing a sampling method]
A Yaoundé education authority wants to estimate the proportion of Upper Sixth students in the city who own a scientific calculator. There are 60 schools, of very different sizes, spread across 7 subdivisions. Propose a suitable sampling design, naming the method(s) used at each stage and justifying your choices.
\end{exercice}

\begin{corrige}[Exercise 1]
A sensible design is a \textbf{stratified multi-stage cluster sample}. Stage 1: treat the 7 subdivisions as \emph{strata} (guaranteeing geographical coverage) and, within each stratum, select a random sample of \emph{schools} — schools act as \emph{clusters}, since it is cheap and practical to visit whole schools. Stage 2: within each selected school, draw a \emph{simple random sample} of Upper Sixth students from the class lists (the frame). If precision within each subdivision matters, use proportional allocation of schools/students to strata according to their Upper Sixth enrolment. Justification: stratification ensures all subdivisions are represented; clustering by school reduces travel cost; random selection at every stage keeps the design a probability sample, so the estimated proportion $\hat{p}$ can later be accompanied by a confidence interval. \hfill$\blacksquare$
\end{corrige}

\coursec{Sampling Errors and Point Estimation}

In the previous section we learned \emph{how} to select a sample from a population: simple random sampling, stratified sampling, systematic sampling, and so on. But selecting a sample is only the first step. Once the sample is in our hands, we use it to \emph{say something} about the whole population — and this is where two natural questions arise. First: \emph{how far wrong can we be?} Second: \emph{what single number should we report as our best guess?} These two questions lead us to the notions of \cle{sampling error} and \cle{point estimation}, which form the heart of this section.

\courssub{Sampling Error and Non-Sampling Error}

\textbf{A motivating situation.}\par
Suppose the Ministry of Public Health wishes to estimate the mean height $\mu$ of all Form Five students in Cameroon. The true value $\mu$ exists, but it is unknown — measuring every student in the country would be far too costly. An investigator selects a random sample of $n = 50$ students in Yaoundé, measures them, and computes the sample mean $\bar{x}$. Almost certainly $\bar{x} \neq \mu$. The difference is not a ``mistake'' in the everyday sense: even with a perfectly honest, perfectly executed survey, different samples give different means. This unavoidable variability is what we call \emph{sampling error}.

\begin{definition}[Sampling error]
The \cle{sampling error} of a statistic is the difference between the value of the statistic computed from a sample and the true value of the corresponding population parameter:
\[
\text{sampling error} = \text{statistic} - \text{parameter}.
\]
For the sample mean, the sampling error is $\bar{x} - \mu$. It is \textbf{not a mistake}: it arises purely because we observe a sample rather than the whole population. It \emph{varies from sample to sample}, and its expected size \emph{decreases as the sample size $n$ increases}.
\end{definition}

\begin{exemplebox}[Sampling error in action]
A bag contains the $5$ numbers $2, 4, 6, 8, 10$ (the population). The population mean is $\mu = \dfrac{2+4+6+8+10}{5} = 6$. If we draw a sample of size $2$ and obtain $\{4, 10\}$, then $\bar{x} = 7$ and the sampling error is $7 - 6 = +1$. If instead we draw $\{2, 6\}$, then $\bar{x} = 4$ and the sampling error is $4 - 6 = -2$. Same population, same sampling method, different errors: the error is a \emph{random} quantity.
\end{exemplebox}

By contrast, \cle{non-sampling errors} are genuine defects in the survey process. They would occur even if we attempted a complete census. The main types are:

\begin{itemize}
\item \textbf{Measurement error:} the recorded value differs from the true value (a faulty scale, a respondent who lies about his income, a badly worded question).
\item \textbf{Processing error:} mistakes made after collection — data entry errors, coding errors, calculation slips.
\item \textbf{Coverage error:} the sampling frame does not match the target population (a telephone survey excludes households without phones; some selected individuals refuse to answer — \emph{non-response}).
\end{itemize}

\begin{attention}[Sampling error vs non-sampling error]
Increasing the sample size reduces \emph{sampling} error, but it does \textbf{not} reduce \emph{non-sampling} errors. A badly designed questionnaire given to $10\,000$ people produces $10\,000$ bad answers. In many real surveys, non-sampling errors are larger and more dangerous than sampling error.
\end{attention}

\courssub{The Sampling Distribution of a Statistic}

Since the sampling error varies from sample to sample, the statistic itself is a \emph{random variable}. Before the sample is drawn, $\bar{X}$ is random; after the sample is drawn, $\bar{x}$ is one observed value of it.

\begin{definition}[Sampling distribution]
The \cle{sampling distribution} of a statistic is the probability distribution of the values the statistic takes over \emph{all possible samples} of a fixed size $n$ drawn from the population.
\end{definition}

\begin{propriete}[Mean and variance of the sample mean]
Let $X_1, X_2, \dots, X_n$ be a random sample from a population with mean $\mu$ and variance $\sigma^2$. Then
\[
E(\bar{X}) = \mu \qquad \text{and} \qquad \operatorname{Var}(\bar{X}) = \frac{\sigma^2}{n}.
\]
The quantity $\dfrac{\sigma}{\sqrt{n}}$ is called the \cle{standard error} of the mean.
\end{propriete}

\textbf{Proof (examinable and instructive).} Since the $X_i$ are independent with $E(X_i)=\mu$ and $\operatorname{Var}(X_i)=\sigma^2$,
\[
E(\bar{X}) = E\!\left(\frac{1}{n}\sum_{i=1}^{n} X_i\right) = \frac{1}{n}\sum_{i=1}^{n} E(X_i) = \frac{1}{n}(n\mu) = \mu,
\]
and, using independence so that variances add,
\[
\operatorname{Var}(\bar{X}) = \operatorname{Var}\!\left(\frac{1}{n}\sum_{i=1}^{n} X_i\right) = \frac{1}{n^2}\sum_{i=1}^{n}\operatorname{Var}(X_i) = \frac{n\sigma^2}{n^2} = \frac{\sigma^2}{n}. \qquad \blacksquare
\]

Notice what this says: on average the sample mean hits the target $\mu$ exactly, and its spread about $\mu$ shrinks like $1/\sqrt{n}$. To halve the standard error you must \emph{quadruple} the sample size.

\begin{propriete}[Central Limit Theorem (CLT)]
Let $X_1, \dots, X_n$ be a random sample from \emph{any} population with mean $\mu$ and finite variance $\sigma^2$. Then for large $n$ (in practice $n \geq 30$),
\[
\bar{X} \sim N\!\left(\mu,\ \frac{\sigma^2}{n}\right) \quad \text{approximately}, \qquad
Z = \frac{\bar{X} - \mu}{\sigma/\sqrt{n}} \sim N(0,1) \quad \text{approximately}.
\]
If the population itself is normal, the result is \emph{exact} for every $n$. (Proof not required at this level.)
\end{propriete}

\begin{savaistu}[A brief history of the CLT]
The Central Limit Theorem has a long history. Abraham de Moivre (1733) first proved it for the binomial distribution (coin tosses). Pierre-Simon Laplace extended it in 1810. The general form — for \emph{any} distribution with finite variance — was finally proved by Aleksandr Lyapunov in 1901. The term ``Central Limit Theorem'' was coined by George Pólya in 1920 because it serves as the ``centre'' of probability theory.
\end{savaistu}

The CLT is the most powerful result in this course: whatever the shape of the population — skewed, uniform, bimodal — the distribution of $\bar{X}$ becomes bell-shaped as $n$ grows. This is exactly why the normal distribution dominates all of estimation theory.

\begin{popfigure}
\begin{center}
\begin{tikzpicture}
\begin{axis}[
  width=12cm, height=6.5cm,
  axis lines=middle,
  xlabel={$\bar{x}$}, ylabel={density},
  xlabel style={at={(ticklabel* cs:1)}, anchor=west},
  xmin=8, xmax=22, ymin=0, ymax=0.55,
  xtick={15}, xticklabels={$\mu$},
  ytick=\empty,
  samples=200, domain=8:22,
  clip=false
]
\addplot[very thick, popblue] {0.45*exp(-((x-15)^2)/(2*1.8^2))};
\addplot[very thick, poporange, dashed] {0.18*exp(-((x-15)^2)/(2*4.2^2))};
\draw[dashed, popdark] (axis cs:15,0) -- (axis cs:15,0.45);
\node[popblue, align=center] at (axis cs:18.6,0.42) {small $\sigma/\sqrt{n}$\\ (large $n$)};
\node[poporange, align=center] at (axis cs:10.6,0.22) {large $\sigma/\sqrt{n}$\\ (small $n$)};
\end{axis}
\end{tikzpicture}
\end{center}
\legende{Sampling distribution of $\bar{X}$: normal curves centred on the population mean $\mu$. The spread (standard error $\sigma/\sqrt{n}$) decreases as $n$ increases.}
\end{popfigure}

\begin{experience}[Simulating the Central Limit Theorem]
Take a very non-normal population: the result of rolling one fair die (uniform on $\{1,2,3,4,5,6\}$, with $\mu = 3.5$ and $\sigma^2 = \frac{35}{12}$). Roll the die $n = 20$ times, compute the mean; repeat this experiment $200$ times and plot a histogram of the $200$ means. Although single rolls are flat, the histogram of means is unmistakably bell-shaped and centred near $3.5$ — the CLT made visible.
\end{experience}

\begin{popfigure}
\begin{center}
\begin{tikzpicture}
\begin{axis}[
  width=12cm, height=6.5cm,
  ybar, bar width=7pt,
  axis lines=left,
  xlabel={sample mean $\bar{x}$ of $20$ die rolls},
  ylabel={frequency},
  xmin=2.6, xmax=4.4, ymin=0, ymax=45,
  xtick={2.8,3.0,3.2,3.4,3.6,3.8,4.0,4.2},
  ytick={0,10,20,30,40},
]
\addplot+[fill=popblueL, draw=popblue] coordinates {
  (2.75,1) (2.85,2) (2.95,4) (3.05,7) (3.15,12) (3.25,18)
  (3.35,26) (3.45,34) (3.55,38) (3.65,30) (3.75,20)
  (3.85,12) (3.95,7) (4.05,4) (4.15,2) (4.25,1)
};
\addplot[very thick, poporange, domain=2.6:4.4, samples=200]
  {40*exp(-((x-3.5)^2)/(2*0.38^2))};
\node[poporange] at (axis cs:4.05,33) {CLT curve};
\end{axis}
\end{tikzpicture}
\end{center}
\legende{Histogram of $200$ simulated sample means (each the mean of $20$ die rolls), overlaid with the normal curve predicted by the Central Limit Theorem, centred at $\mu = 3.5$.}
\end{popfigure}

\courssub{Point Estimators and Point Estimates}

We now turn to the second question: what single number should we report?

\begin{definition}[Estimator and estimate]
A \cle{point estimator} is a statistic (a rule or formula, a random variable) used to estimate a population parameter. A \cle{point estimate} is the numerical value taken by the estimator for one particular observed sample.
\end{definition}

The distinction matters: $\bar{X}$ (capital letters, before sampling) is the \emph{estimator}; $\bar{x} = 165.4\ \mathrm{cm}$ (computed from the data) is the \emph{estimate}. The two most common point estimators are:

\begin{center}
\begin{tabularx}{\linewidth}{|l|X|X|}
\hline
\rowcolor{popblueL} \textbf{Parameter} & \textbf{Point estimator} & \textbf{Point estimate} \\
\hline
Population mean $\mu$ & $\bar{X} = \dfrac{1}{n}\displaystyle\sum_{i=1}^{n} X_i$ & $\bar{x} = \dfrac{\sum x_i}{n}$ \\
\hline
Population proportion $p$ & $\hat{P} = \dfrac{X}{n}$ ($X$ = number of ``successes'') & $\hat{p} = \dfrac{x}{n}$ \\
\hline
Population variance $\sigma^2$ & $S^2 = \dfrac{1}{n-1}\displaystyle\sum_{i=1}^{n}(X_i - \bar{X})^2$ & $s^2 = \dfrac{\sum (x_i-\bar{x})^2}{n-1}$ \\
\hline
\end{tabularx}
\end{center}

\begin{methode}[Computing a point estimate of $\mu$]
\begin{enumerate}
\item Collect the random sample $x_1, x_2, \dots, x_n$ and count the observations: $n$.
\item Add all the values: $\sum x_i$.
\item Divide by the sample size: $\bar{x} = \dfrac{\sum x_i}{n}$.
\item Report $\bar{x}$ as the point estimate of $\mu$, together with $n$ (the size matters for judging reliability).
\end{enumerate}
For a proportion: count the number $x$ of sampled units having the required characteristic and compute $\hat{p} = x/n$.
\end{methode}

\begin{exoresolu}[Point estimates from a market survey]
A trader at Mokolo market wants to estimate the mean daily expenditure $\mu$ of her customers and the proportion $p$ of customers who pay by mobile money. She observes a random sample of $8$ customers whose expenditures (in thousands of FCFA) are
\[
4.5,\quad 6.0,\quad 3.5,\quad 5.0,\quad 7.5,\quad 4.0,\quad 5.5,\quad 6.0,
\]
and notes that $3$ of the $8$ paid by mobile money. Find point estimates of $\mu$ and $p$.
\tcblower
\textbf{Solution.} For the mean:
\[
\sum x_i = 4.5 + 6.0 + 3.5 + 5.0 + 7.5 + 4.0 + 5.5 + 6.0 = 42.0,
\qquad
\bar{x} = \frac{42.0}{8} = 5.25.
\]
The point estimate of the mean daily expenditure is $\bar{x} = 5.25$ thousand FCFA, i.e. $5\,250\ \mathrm{FCFA}$. For the proportion:
\[
\hat{p} = \frac{x}{n} = \frac{3}{8} = 0.375.
\]
The point estimate of the proportion paying by mobile money is $0.375$, i.e. $37.5\%$.
\end{exoresolu}

\begin{attention}[A point estimate says nothing about its own precision]
Reporting only ``$\bar{x} = 5\,250\ \mathrm{FCFA}$'' hides the most important information: \emph{how trustworthy is this number?} An estimate from $n = 8$ customers and an estimate from $n = 800$ customers may give the same value of $\bar{x}$, yet their precisions are vastly different. A single point estimate carries no measure of its possible error — this is precisely why we shall need \emph{confidence intervals} later in the chapter.
\end{attention}

\courssub{Measuring Accuracy: The Mean Squared Error}

How do we judge whether an estimator $T$ of a parameter $\theta$ is ``good''? A natural idea is to look at the error $T - \theta$. Since this error is random and can be positive or negative, we square it and average:

\begin{definition}[Mean squared error]
The \cle{mean squared error} of an estimator $T$ of a parameter $\theta$ is
\[
\operatorname{MSE}(T) = E\!\left[(T - \theta)^2\right].
\]
The smaller the MSE, the more accurate the estimator.
\end{definition}

\begin{propriete}[Bias–variance decomposition]
For any estimator $T$ of $\theta$,
\[
\operatorname{MSE}(T) = \operatorname{Var}(T) + \bigl(\operatorname{bias}(T)\bigr)^2,
\qquad \text{where } \operatorname{bias}(T) = E(T) - \theta.
\]
In particular, if $T$ is \emph{unbiased} ($E(T) = \theta$), then $\operatorname{MSE}(T) = \operatorname{Var}(T)$.
\end{propriete}

\textbf{Proof.} Write $T - \theta = (T - E(T)) + (E(T) - \theta)$. Squaring and taking expectations,
\[
E[(T-\theta)^2] = E[(T - E(T))^2] + 2\,(E(T)-\theta)\,E[T - E(T)] + (E(T)-\theta)^2.
\]
The middle term vanishes because $E[T - E(T)] = 0$, leaving $\operatorname{Var}(T) + (\operatorname{bias}(T))^2$. $\blacksquare$

This decomposition is deeply intuitive: inaccuracy comes from two sources — the \emph{variability} of the estimator (it jumps around) and its \emph{systematic deviation} from the target (it aims at the wrong place). A good estimator controls both.

\begin{exoresolu}[Comparing two estimators with the MSE]
A population has mean $\mu$ and variance $\sigma^2 = 36$. A sample of size $n = 4$ is taken. Two estimators of $\mu$ are proposed:
\[
T_1 = \bar{X} \qquad \text{and} \qquad T_2 = \frac{X_1 + X_2}{2} + 1.
\]
Compute the MSE of each estimator and decide which is preferable.
\tcblower
\textbf{Solution.} For $T_1 = \bar{X}$: $E(T_1) = \mu$, so $T_1$ is unbiased and
\[
\operatorname{MSE}(T_1) = \operatorname{Var}(\bar{X}) = \frac{\sigma^2}{n} = \frac{36}{4} = 9.
\]
For $T_2$: $E(T_2) = \dfrac{\mu + \mu}{2} + 1 = \mu + 1$, so $\operatorname{bias}(T_2) = 1$. Also
\[
\operatorname{Var}(T_2) = \operatorname{Var}\!\left(\frac{X_1 + X_2}{2}\right) = \frac{1}{4}\bigl(\sigma^2 + \sigma^2\bigr) = \frac{2 \times 36}{4} = 18.
\]
Hence
\[
\operatorname{MSE}(T_2) = \operatorname{Var}(T_2) + \bigl(\operatorname{bias}(T_2)\bigr)^2 = 18 + 1^2 = 19.
\]
Since $\operatorname{MSE}(T_1) = 9 < 19 = \operatorname{MSE}(T_2)$, the sample mean $T_1$ is the more accurate estimator — it is both unbiased \emph{and} less variable.
\end{exoresolu}

\begin{attention}[Common mistakes with sampling error and the MSE]
\begin{itemize}
\item Confusing the \emph{standard deviation} $\sigma$ (spread of the population) with the \emph{standard error} $\sigma/\sqrt{n}$ (spread of $\bar{X}$). Forgetting the $\sqrt{n}$ is the most frequent error in examinations.
\item Believing the CLT requires a normal population: it does not — that is its whole power. Normality of the population only guarantees normality of $\bar{X}$ for \emph{small} $n$.
\item Writing $\operatorname{MSE}(T) = \operatorname{Var}(T)$ when the estimator is biased: the bias term $(E(T)-\theta)^2$ must be added.
\item Treating the sampling error $\bar{x} - \mu$ as a fixed number: it is random and changes from sample to sample.
\end{itemize}
\end{attention}

\begin{exercice}[Die-roll sampling distribution]
A population consists of the numbers $\{1, 3, 5\}$ with equal probability. Samples of size $2$ are drawn with replacement.
\begin{enumerate}
\item Compute $\mu$ and $\sigma^2$.
\item List all $9$ equally likely samples and the corresponding values of $\bar{X}$; hence give the sampling distribution of $\bar{X}$.
\item Verify directly that $E(\bar{X}) = \mu$ and $\operatorname{Var}(\bar{X}) = \sigma^2/2$.
\end{enumerate}
\end{exercice}

\begin{corrige}[Exercise 1]
\textbf{1.} $\mu = \dfrac{1+3+5}{3} = 3$; \quad $E(X^2) = \dfrac{1+9+25}{3} = \dfrac{35}{3}$, so $\sigma^2 = \dfrac{35}{3} - 9 = \dfrac{8}{3}$.

\textbf{2.} The $9$ samples give means: $(1,1)\to1$; $(1,3),(3,1)\to2$; $(1,5),(3,3),(5,1)\to3$; $(3,5),(5,3)\to4$; $(5,5)\to5$. Hence

\begin{center}
\begin{tabular}{|c|ccccc|}
\hline
\rowcolor{popblueL} $\bar{x}$ & $1$ & $2$ & $3$ & $4$ & $5$ \\
\hline
$P(\bar{X}=\bar{x})$ & $\tfrac{1}{9}$ & $\tfrac{2}{9}$ & $\tfrac{3}{9}$ & $\tfrac{2}{9}$ & $\tfrac{1}{9}$ \\
\hline
\end{tabular}
\end{center}

\textbf{3.} $E(\bar{X}) = \dfrac{1(1)+2(2)+3(3)+4(2)+5(1)}{9} = \dfrac{27}{9} = 3 = \mu$. \checkmark

$E(\bar{X}^2) = \dfrac{1(1)+4(2)+9(3)+16(2)+25(1)}{9} = \dfrac{93}{9} = \dfrac{31}{3}$, so $\operatorname{Var}(\bar{X}) = \dfrac{31}{3} - 9 = \dfrac{4}{3} = \dfrac{\sigma^2}{2}$. \checkmark \quad Even for $n=2$, the distribution of $\bar{X}$ is already triangular — a first step towards the bell shape promised by the CLT.
\end{corrige}

\begin{aretenir}[Key points of Section 2]
\begin{itemize}
\item \textbf{Sampling error} $=$ statistic $-$ parameter; it is random, unavoidable, and shrinks as $n$ grows. \textbf{Non-sampling errors} (measurement, processing, coverage) are defects of the survey itself and are not cured by larger samples.
\item The \textbf{sampling distribution} of $\bar{X}$ has $E(\bar{X}) = \mu$ and $\operatorname{Var}(\bar{X}) = \sigma^2/n$; the \textbf{standard error} is $\sigma/\sqrt{n}$.
\item \textbf{Central Limit Theorem:} for large $n$, $\bar{X} \sim N(\mu, \sigma^2/n)$ approximately, whatever the population.
\item A \textbf{point estimator} is a rule (random variable); a \textbf{point estimate} is its observed value. Standard estimators: $\bar{X}$ for $\mu$, $\hat{P} = X/n$ for $p$.
\item Accuracy is measured by $\operatorname{MSE}(T) = \operatorname{Var}(T) + (\operatorname{bias}(T))^2$; for unbiased estimators, $\operatorname{MSE} = \operatorname{Var}$.
\item A point estimate alone gives \textbf{no information about its precision} — this motivates the confidence intervals of the next sections.
\end{itemize}
\end{aretenir}

Having learned how to \emph{measure} the quality of an estimator through its bias, variance and mean squared error, we are ready for the natural next step: Section 3 will turn these ideas into formal \emph{criteria} — the properties that define a genuinely \emph{good} estimator.

\coursec{Properties of Good Estimators}

In the previous section we met \cle{point estimation}: from a random sample $X_1, X_2, \dots, X_n$ we compute a single number $\hat{\theta}$ intended to guess an unknown population parameter $\theta$. We saw, for example, that the sample mean $\bar{X}$ is a natural estimator of the population mean $\mu$, and the sample proportion $\hat{p}$ estimates a population proportion $p$.

But here is a natural question a careful student should ask: \emph{there are infinitely many possible estimators of the same parameter — how do we choose a good one?} For instance, to estimate the mean height of Upper Sixth students in Cameroon, should we use the sample mean, the sample median, the average of the largest and smallest values, or simply the first value observed? All are legitimate estimators. Statisticians have developed four quality criteria to compare them: \cle{unbiasedness}, \cle{consistency}, \cle{efficiency} and \cle{sufficiency}. This section studies each of them in turn, following the official GCE Advanced Level syllabus (Lesson 188).

\courssub{Unbiasedness: aiming at the right target}

\textbf{Intuition.} Imagine a marksman shooting at a target. If his shots are scattered but \emph{centred} on the bullseye, we say his aim is good: on average he hits the right spot. An estimator is like the marksman: each sample gives a different value of $\hat{\theta}$ (this is sampling variability), but we would like the \emph{average} of all possible values of $\hat{\theta}$ to be exactly $\theta$.

\begin{definition}[Unbiased estimator]
An estimator $\hat{\theta}$ of a parameter $\theta$ is said to be \cle{unbiased} if
\[ \mathbb{E}(\hat{\theta}) = \theta. \]
The quantity $\operatorname{bias}(\hat{\theta}) = \mathbb{E}(\hat{\theta}) - \theta$ is called the \cle{bias} of the estimator. If $\operatorname{bias}(\hat{\theta}) \neq 0$, the estimator is \emph{biased}.
\end{definition}

\begin{methode}[Checking unbiasedness]
To show that an estimator $\hat{\theta}$ is unbiased:
\begin{enumerate}
\item Write $\hat{\theta}$ explicitly as a function of the sample $X_1, \dots, X_n$.
\item Compute $\mathbb{E}(\hat{\theta})$ using linearity of expectation and the fact that $\mathbb{E}(X_i) = \mu$ for every $i$ (or the appropriate population moment).
\item Compare the result with $\theta$: if $\mathbb{E}(\hat{\theta}) = \theta$, the estimator is unbiased; otherwise compute the bias.
\end{enumerate}
\end{methode}

\begin{exemplebox}[The sample mean is unbiased for $\mu$]
Let $X_1, \dots, X_n$ be a random sample from a population with mean $\mu$. Then
\[ \mathbb{E}(\bar{X}) = \mathbb{E}\!\left(\frac{1}{n}\sum_{i=1}^{n} X_i\right) = \frac{1}{n}\sum_{i=1}^{n}\mathbb{E}(X_i) = \frac{1}{n}(n\mu) = \mu. \]
Hence $\bar{X}$ is an unbiased estimator of $\mu$.
\end{exemplebox}

\begin{exemplebox}[The sample proportion is unbiased for $p$]
Let $X_1, \dots, X_n$ be a random sample from a Bernoulli population with parameter $p$ (success probability). The sample proportion is $\hat{p} = \frac{1}{n}\sum_{i=1}^n X_i$. Since $\mathbb{E}(X_i)=p$,
\[ \mathbb{E}(\hat{p}) = \frac{1}{n}\sum_{i=1}^n \mathbb{E}(X_i) = \frac{1}{n}(np) = p. \]
Thus $\hat{p}$ is an unbiased estimator of $p$.
\end{exemplebox}

\begin{attention}[The trap of the "dividing by $n$" sample variance]
The quantity $S_n^2 = \dfrac{1}{n}\sum_{i=1}^n (X_i - \bar{X})^2$ is \emph{not} an unbiased estimator of $\sigma^2$: one can show that $\mathbb{E}(S_n^2) = \dfrac{n-1}{n}\sigma^2$. This is why we define the \cle{sample variance} with divisor $n-1$:
\[ S^2 = \frac{1}{n-1}\sum_{i=1}^n (X_i - \bar{X})^2, \qquad \mathbb{E}(S^2) = \sigma^2. \]
Dividing by $n$ instead of $n-1$ is one of the most frequent errors at the GCE. The correction $n-1$ is called \emph{Bessel's correction} and accounts for the loss of one degree of freedom when estimating $\mu$ by $\bar{X}$.
\end{attention}

\courssub{Consistency: improving with more data}

\textbf{Intuition.} A good estimator should get better and better as the sample size grows. If you interview 30 voters in Bamenda, your estimate of the proportion favouring a candidate is rough; if you interview 30\,000, it should be very close to the truth. An estimator with this property is called \emph{consistent}.

\begin{definition}[Consistent estimator]
An estimator $\hat{\theta}_n$ (based on a sample of size $n$) is \cle{consistent} for $\theta$ if it converges to $\theta$ in probability as $n \to \infty$:
\[ \text{for every } \varepsilon > 0, \quad P\big(|\hat{\theta}_n - \theta| > \varepsilon\big) \longrightarrow 0 \quad \text{as } n \to \infty. \]
This is also called \emph{weak consistency}. (Strong consistency, i.e. almost sure convergence, is not required at this level.)
\end{definition}

A practical sufficient condition, easily checked in examinations, is the following.

\begin{propriete}[Sufficient condition for consistency (Chebyshev)]
If $\hat{\theta}_n$ is unbiased (or asymptotically unbiased, i.e. $\mathbb{E}(\hat{\theta}_n) \to \theta$) and $\operatorname{Var}(\hat{\theta}_n) \to 0$ as $n \to \infty$, then $\hat{\theta}_n$ is consistent. This follows from Chebyshev's inequality:
\[ P\big(|\hat{\theta}_n - \theta| > \varepsilon\big) \le \frac{\operatorname{Var}(\hat{\theta}_n) + [\mathbb{E}(\hat{\theta}_n)-\theta]^2}{\varepsilon^2} \longrightarrow 0. \]
\end{propriete}

\begin{exemplebox}[Consistency of $\bar{X}$ and $\hat{p}$]
Since $\mathbb{E}(\bar{X}) = \mu$ and $\operatorname{Var}(\bar{X}) = \dfrac{\sigma^2}{n} \to 0$, the sample mean is a consistent estimator of $\mu$. Similarly, the sample proportion $\hat{p}$ satisfies $\mathbb{E}(\hat{p})=p$ and $\operatorname{Var}(\hat{p}) = \dfrac{p(1-p)}{n} \to 0$, so $\hat{p}$ is consistent for $p$.
\end{exemplebox}

\begin{exemplebox}[An unbiased but inconsistent estimator]
Consider $\hat{\mu} = X_1$ (using only the first observation). It is unbiased because $\mathbb{E}(X_1)=\mu$, but $\operatorname{Var}(\hat{\mu}) = \sigma^2$ does \emph{not} tend to $0$ as $n$ increases. Hence $\hat{\mu}$ is \emph{not} consistent. This shows that unbiasedness alone does not guarantee consistency.
\end{exemplebox}

\courssub{Efficiency: among unbiased estimators, choose the least variable}

\textbf{Intuition.} Two marksmen may both aim correctly (both unbiased), but the one whose shots are \emph{tightly grouped} is better. Between two unbiased estimators, we prefer the one with the \emph{smaller variance}, because its values fluctuate less from sample to sample.

\begin{definition}[Efficient estimator]
Let $\hat{\theta}_1$ and $\hat{\theta}_2$ be two unbiased estimators of $\theta$. We say that $\hat{\theta}_1$ is \cle{more efficient} than $\hat{\theta}_2$ if
\[ \operatorname{Var}(\hat{\theta}_1) < \operatorname{Var}(\hat{\theta}_2). \]
An unbiased estimator whose variance attains the smallest possible value among all unbiased estimators is called \emph{efficient} (or a \emph{minimum variance unbiased estimator}, MVUE).
\end{definition}

How small can the variance be? A celebrated result gives a floor that no unbiased estimator can go below.

\begin{propriete}[Cramér–Rao lower bound]
Under regularity conditions (the support of the distribution does not depend on $\theta$, and the density is differentiable with respect to $\theta$), for any unbiased estimator $\hat{\theta}$ of $\theta$ based on a sample of size $n$,
\[ \operatorname{Var}(\hat{\theta}) \ge \frac{1}{n\,\mathcal{I}(\theta)}, \qquad \text{where } \mathcal{I}(\theta) = \mathbb{E}\!\left[\left(\frac{\partial}{\partial \theta}\ln f(X;\theta)\right)^2\right] \]
is the \cle{Fisher information} of a single observation. An unbiased estimator attaining this bound is efficient. (The proof is \emph{not} examinable; the statement and its use are.)
\end{propriete}

\begin{propriete}[The sample mean is BLUE and, for normal data, MVUE]
Under i.i.d. sampling from a population with finite variance $\sigma^2$ (normality \emph{not} required), the sample mean $\bar{X}$ is the \cle{BLUE} — \emph{Best Linear Unbiased Estimator} — of $\mu$: among all unbiased estimators that are linear combinations $\sum a_i X_i$ of the observations, $\bar{X}$ has the smallest variance, namely $\operatorname{Var}(\bar{X}) = \sigma^2/n$. This is the Gauss–Markov theorem.

If, in addition, the population is normal $N(\mu, \sigma^2)$, then $\bar{X}$ attains the Cramér–Rao bound, so it is efficient among \emph{all} unbiased estimators (not just linear ones). In this case $\bar{X}$ is the MVUE for $\mu$. (Proofs are not examinable.)
\end{propriete}

\begin{attention}[Unbiased but useless]
An unbiased estimator may have a \emph{huge} variance and therefore be practically worthless. For example, the estimator $\hat{\mu} = X_1$ (use only the first observation) is unbiased since $\mathbb{E}(X_1) = \mu$, but $\operatorname{Var}(X_1) = \sigma^2$, which does not decrease with $n$: it is unbiased but neither efficient nor consistent. \textbf{Unbiasedness alone is not enough — efficiency matters.}
\end{attention}

\courssub{Sufficiency: keeping all the information}

\textbf{Intuition.} Suppose a census officer collects the incomes of $n$ households and wants to estimate the mean income $\mu$. Must he keep the full list of $n$ figures, or is the sample mean enough? A statistic is \emph{sufficient} if, once you know it, the individual data values tell you \emph{nothing more} about the parameter. It compresses the data without losing any relevant information.

\begin{definition}[Sufficient statistic]
A statistic $T = T(X_1, \dots, X_n)$ is \cle{sufficient} for a parameter $\theta$ if the conditional distribution of the sample given $T$ does not depend on $\theta$. Intuitively: $T$ captures \emph{all} the information the sample contains about $\theta$.
\end{definition}

Checking this definition directly is difficult; the following theorem gives a practical recipe.

\begin{propriete}[Fisher–Neyman factorisation theorem]
$T$ is sufficient for $\theta$ if and only if the joint density (or probability mass function) of the sample factorises as
\[ f(x_1, \dots, x_n; \theta) = g\big(T(x_1, \dots, x_n); \theta\big) \times h(x_1, \dots, x_n), \]
where $g$ depends on the data \emph{only through} $T$ and on $\theta$, while $h$ does \emph{not} involve $\theta$.
\end{propriete}

\begin{popfigure}
\begin{center}
\begin{tikzpicture}[>=stealth, node distance=1.1cm]
\node[draw=popblue, fill=popblueL, rounded corners, text width=3.1cm, align=center, font=\small] (data) {\textbf{Raw sample}\\ $x_1, x_2, \dots, x_n$};
\node[draw=popdark, fill=popgreenL, rounded corners, text width=3.4cm, align=center, font=\small, right=2.6cm of data] (T) {\textbf{Sufficient statistic}\\ $T = g\text{-part, e.g. } \sum x_i$};
\node[draw=poppurple, fill=white, rounded corners, text width=3.2cm, align=center, font=\small, right=2.6cm of T] (theta) {\textbf{Parameter $\theta$}\\ all information flows through $T$};
\draw[->, very thick, popblue] (data) -- node[above, font=\small]{summarise} (T);
\draw[->, very thick, popdark] (T) -- node[above, font=\small]{estimate} (theta);
\node[below=0.5cm of T, text width=5.6cm, align=center, font=\small\itshape, text=popmuted] {$f(\mathbf{x};\theta) = \underbrace{g(T;\theta)}_{\text{keeps } \theta} \times \underbrace{h(\mathbf{x})}_{\text{no } \theta}$};
\end{tikzpicture}
\end{center}
\legende{The factorisation theorem: the likelihood splits into a part carrying $\theta$ only through $T$ and a part free of $\theta$.}
\end{popfigure}

\begin{exoresolu}[Sufficiency of $\sum X_i$ for a Poisson mean]
Let $X_1, \dots, X_n$ be i.i.d. Poisson with mean $\lambda$. Show that $T = \sum_{i=1}^n X_i$ is sufficient for $\lambda$.
\tcblower
The joint probability mass function is
\[ f(\mathbf{x};\lambda) = \prod_{i=1}^n \frac{e^{-\lambda}\lambda^{x_i}}{x_i!} = e^{-n\lambda}\lambda^{\sum x_i} \times \frac{1}{\prod x_i!}. \]
Take $g(T;\lambda) = e^{-n\lambda}\lambda^{T}$ (depends on the data only through $T = \sum x_i$) and $h(\mathbf{x}) = 1/\prod x_i!$ (free of $\lambda$). By the factorisation theorem, $T = \sum X_i$ — equivalently $\bar{X} = T/n$ — is sufficient for $\lambda$.
\end{exoresolu}

\begin{exemplebox}[Sufficiency of $\bar{X}$ for $\mu$ in $N(\mu,\sigma^2)$ with known $\sigma^2$]
The joint density is
\[ f(\mathbf{x};\mu) = (2\pi\sigma^2)^{-n/2} \exp\!\left(-\frac{1}{2\sigma^2}\sum(x_i-\mu)^2\right). \]
Expanding $\sum(x_i-\mu)^2 = \sum(x_i-\bar{x})^2 + n(\bar{x}-\mu)^2$, the density factorises as
\[ \underbrace{(2\pi\sigma^2)^{-n/2} \exp\!\left(-\frac{n(\bar{x}-\mu)^2}{2\sigma^2}\right)}_{g(\bar{x};\mu)} \times \underbrace{\exp\!\left(-\frac{\sum(x_i-\bar{x})^2}{2\sigma^2}\right)}_{h(\mathbf{x})}. \]
Hence $\bar{X}$ is sufficient for $\mu$.
\end{exemplebox}

\courssub{Sample mean versus sample median: efficiency in action}

A classical comparison illustrates all four criteria at once. For a \emph{normal} population $N(\mu, \sigma^2)$, both $\bar{X}$ and the sample median $\tilde{X}$ are unbiased and consistent for $\mu$. But their variances differ:
\[ \operatorname{Var}(\bar{X}) = \frac{\sigma^2}{n}, \qquad \operatorname{Var}(\tilde{X}) \approx \frac{\pi\,\sigma^2}{2n} \approx 1.571\,\frac{\sigma^2}{n} \quad (\text{asymptotic variance}). \]
So for normal data the sample mean is about $36\%$ more efficient: to match the precision of $\bar{X}$ from $100$ observations, the median would need about $157$.

However, for \emph{heavy-tailed} distributions (data with frequent extreme outliers, such as incomes in a town where a few traders earn one hundred times more than most farmers), the situation reverses: a single extreme value pulls $\bar{X}$ violently, while the median, which depends only on the \emph{middle order}, is barely affected. The median is then the more efficient (more "robust") estimator.

\begin{popfigure}
\begin{center}
\begin{tikzpicture}
\begin{axis}[
  width=12cm, height=6.6cm,
  xlabel={Sample size $n$}, ylabel={Variance of the estimator ($\sigma^2=1$)},
  xmin=0, xmax=100, ymin=0, ymax=1.05,
  xtick={20,40,60,80,100},
  legend style={at={(0.97,0.97)}, anchor=north east, font=\small},
  grid=both, grid style={popline, very thin},
]
\addplot[domain=5:100, samples=60, very thick, popblue] {1/x};
\addlegendentry{$\operatorname{Var}(\overline{X}) = 1/n$ (normal data)}
\addplot[domain=5:100, samples=60, very thick, poporange, dashed] {1.571/x};
\addlegendentry{$\operatorname{Var}(\tilde{X}) \approx 1.571/n$ (normal data)}
\addplot[domain=5:100, samples=60, very thick, popgreen] {0.55/x};
\addlegendentry{Median{,} heavy-tailed (illustrative)}
\end{axis}
\end{tikzpicture}
\end{center}
\legende{Variance of the sample mean and sample median (with $\sigma^2 = 1$). For normal data the mean (blue) beats the median (orange); for heavy-tailed data the median (green) can win. The green curve is illustrative only.}
\end{popfigure}

\begin{exoresolu}[Comparing two unbiased estimators]
A sample $X_1, \dots, X_n$ is drawn from $N(\mu, \sigma^2)$. Consider $\hat{\mu}_1 = \bar{X}$ and $\hat{\mu}_2 = \dfrac{X_1 + X_n}{2}$. Show both are unbiased and determine which is more efficient.
\tcblower
\textbf{Unbiasedness.} $\mathbb{E}(\hat{\mu}_1) = \mu$ (shown above). Also
\[ \mathbb{E}(\hat{\mu}_2) = \frac{\mathbb{E}(X_1) + \mathbb{E}(X_n)}{2} = \frac{\mu + \mu}{2} = \mu. \]
Both are unbiased. \textbf{Efficiency.} By independence,
\[ \operatorname{Var}(\hat{\mu}_1) = \frac{\sigma^2}{n}, \qquad \operatorname{Var}(\hat{\mu}_2) = \frac{\operatorname{Var}(X_1) + \operatorname{Var}(X_n)}{4} = \frac{2\sigma^2}{4} = \frac{\sigma^2}{2}. \]
For every $n > 2$, $\operatorname{Var}(\hat{\mu}_1) < \operatorname{Var}(\hat{\mu}_2)$, so $\bar{X}$ is more efficient. Note also that $\hat{\mu}_2$ is \emph{not} consistent, since its variance does not tend to $0$.
\end{exoresolu}

\begin{aretenir}[The four quality criteria — summary for GCE]
\begin{itemize}
\item \textbf{Unbiased:} $\mathbb{E}(\hat{\theta}) = \theta$ — right on average.
\item \textbf{Consistent:} $\hat{\theta} \to \theta$ in probability as $n \to \infty$ — improves with more data. Sufficient condition: unbiased (or asymptotically unbiased) and $\operatorname{Var}(\hat{\theta}_n) \to 0$.
\item \textbf{Efficient:} smallest variance among unbiased estimators; bounded below by the Cramér–Rao bound $1/(n\mathcal{I}(\theta))$. For normal data, $\bar{X}$ is MVUE (attains the bound).
\item \textbf{Sufficient:} uses all the information in the sample; detected by the factorisation $f(\mathbf{x};\theta) = g(T;\theta)\,h(\mathbf{x})$.
\item For normal data, $\bar{X}$ is unbiased, consistent, sufficient and BLUE for $\mu$.
\end{itemize}
\end{aretenir}

\begin{exercice}[Practice]
Let $X_1, \dots, X_n$ be i.i.d. with mean $\mu$ and variance $\sigma^2$. Consider $\hat{\mu} = \dfrac{2X_1 + X_2 + \dots + X_n}{n+1}$.
\begin{enumerate}
\item Show that $\hat{\mu}$ is unbiased for $\mu$.
\item Compute $\operatorname{Var}(\hat{\mu})$ and deduce that $\hat{\mu}$ is consistent.
\item Is $\hat{\mu}$ more or less efficient than $\bar{X}$? Justify.
\end{enumerate}
\end{exercice}

\begin{corrige}[Exercise 1]
\textbf{1.} $\mathbb{E}(\hat{\mu}) = \dfrac{2\mu + (n-1)\mu}{n+1} = \dfrac{(n+1)\mu}{n+1} = \mu$, so $\hat{\mu}$ is unbiased. \quad
\textbf{2.} By independence, $\operatorname{Var}(\hat{\mu}) = \dfrac{4\sigma^2 + (n-1)\sigma^2}{(n+1)^2} = \dfrac{(n+3)\sigma^2}{(n+1)^2} \to 0$ as $n\to\infty$, so $\hat{\mu}$ is consistent. \quad
\textbf{3.} $\operatorname{Var}(\bar{X}) = \dfrac{\sigma^2}{n}$. Compare:
\[ \frac{\operatorname{Var}(\bar{X})}{\operatorname{Var}(\hat{\mu})} = \frac{\sigma^2/n}{(n+3)\sigma^2/(n+1)^2} = \frac{(n+1)^2}{n(n+3)} = \frac{n^2+2n+1}{n^2+3n} < 1 \quad \text{for all } n>1. \]
Hence $\operatorname{Var}(\bar{X}) < \operatorname{Var}(\hat{\mu})$: $\bar{X}$ is more efficient, as expected since it is BLUE.
\end{corrige}

\begin{savaistu}[Why $n-1$?]
The divisor $n-1$ in the sample variance is called \emph{Bessel's correction}, after the German astronomer Friedrich Bessel. Intuitively, the deviations $(X_i - \bar{X})$ are computed from an \emph{estimated} centre $\bar{X}$ rather than the true $\mu$, so they are slightly too small; dividing by $n-1$ instead of $n$ exactly compensates for this loss of one "degree of freedom". This correction is used every day by laboratories, banks and polling institutes all over the world — including those analysing agricultural yields in the North West and South West regions of Cameroon.
\end{savaistu}

\coursec{Confidence Intervals for the Mean (Known Variance)}

In the previous section we studied what makes a point estimator \emph{good}: unbiasedness, consistency and efficiency. But even the best point estimator has an obvious weakness. If we estimate the mean mass of bags of rice produced in a Ndop mill to be $50.2$ kg, this single number tells us nothing about how \emph{close} to the true mean $\mu$ it is likely to be. A point estimate is almost certainly wrong; what we need is a statement of \emph{how wrong} it might be. That is exactly what a \cle{confidence interval} provides: instead of one number, we give a whole \emph{range} of plausible values for $\mu$, together with a measure of how confident we are in that range. In this section we build this idea rigorously in the simplest case, where the population variance $\sigma^{2}$ is known.

\courssub{From a Point Estimate to an Interval Estimate}

Suppose a machine fills packets of sugar and the mass $X$ of a packet is normally distributed with unknown mean $\mu$ and \emph{known} standard deviation $\sigma = 5$ g. We weigh a random sample of $n = 25$ packets and obtain $\bar{x} = 498$ g. Is $\mu = 498$ g? Almost certainly not exactly: another sample of 25 packets would give a slightly different value of $\bar{x}$. However, we know from the sampling distribution of the mean that
\[
\bar{X} \sim N\!\left(\mu,\ \frac{\sigma^{2}}{n}\right),
\]
so most sample means fall within about two standard errors $\sigma/\sqrt{n}$ of $\mu$. Here $\sigma/\sqrt{n} = 5/\sqrt{25} = 1$ g, so $\bar{x}$ is very likely within about $2$ g of $\mu$. It is therefore much more informative to say ``$\mu$ is probably between $496$ g and $500$ g'' than to say ``$\mu = 498$ g''. This interval idea is what we now formalise.

\begin{definition}[Confidence interval, confidence level, margin of error]
Let $\theta$ be an unknown population parameter. A \cle{confidence interval} for $\theta$ is a random interval $(L, U)$, whose endpoints are statistics computed from the sample, constructed so that
\[
P(L \le \theta \le U) = 1 - \alpha,
\]
where the probability is taken over repeated random sampling. The number $1-\alpha$ is called the \cle{confidence level} (often $0.90$, $0.95$ or $0.99$), and $\alpha$ is the \cle{significance level}. For a symmetric interval of the form $\bar{x} \pm E$, the quantity
\[
E = z_{\alpha/2}\,\frac{\sigma}{\sqrt{n}}
\]
is called the \cle{margin of error}. The \cle{width} of the interval is $2E$.
\end{definition}

\begin{attention}[The probability is about the method, not about one interval]
Once a sample has been collected and the interval computed, say $(496, 500)$, the true mean $\mu$ is a fixed (unknown) number: it either lies in $(496, 500)$ or it does not. It is therefore \textbf{wrong} to write ``$P(496 \le \mu \le 500) = 0.95$''. The correct statement is: ``the \emph{procedure} that produced this interval captures the true $\mu$ in $95\%$ of all possible samples''. Examiners regularly penalise candidates who attach the probability to a specific computed interval.
\end{attention}

\courssub{Derivation of the Z-Interval for $\mu$ ($\sigma^{2}$ Known)}

\begin{propriete}[Confidence interval for $\mu$ when $\sigma^{2}$ is known]
Let $X_{1}, \dots, X_{n}$ be a random sample from a population with mean $\mu$ and \emph{known} variance $\sigma^{2}$, where either the population is normal or $n$ is large (so that the Central Limit Theorem applies). Then a $100(1-\alpha)\%$ confidence interval for $\mu$ is
\[
\bar{x} \pm z_{\alpha/2}\,\frac{\sigma}{\sqrt{n}}
\qquad\text{i.e.}\qquad
\left(\bar{x} - z_{\alpha/2}\frac{\sigma}{\sqrt{n}},\ \bar{x} + z_{\alpha/2}\frac{\sigma}{\sqrt{n}}\right),
\]
where $z_{\alpha/2}$ is the critical value satisfying $P(Z > z_{\alpha/2}) = \alpha/2$ for $Z \sim N(0,1)$. \emph{The derivation below is instructive and its steps are examinable.}
\end{propriete}

\textbf{Derivation.} Since $\bar{X} \sim N(\mu, \sigma^{2}/n)$ (exactly for a normal population, approximately for large $n$ by the Central Limit Theorem), standardising gives
\[
Z = \frac{\bar{X} - \mu}{\sigma/\sqrt{n}} \sim N(0,1).
\]
Choose $z_{\alpha/2}$ so that the central area under the standard normal curve is $1-\alpha$:
\[
P\!\left(-z_{\alpha/2} \le \frac{\bar{X}-\mu}{\sigma/\sqrt{n}} \le z_{\alpha/2}\right) = 1-\alpha.
\]
Now we \emph{rearrange the inequality to isolate $\mu$}. Multiplying through by $\sigma/\sqrt{n}$ (positive, so inequalities unchanged):
\[
P\!\left(-z_{\alpha/2}\frac{\sigma}{\sqrt{n}} \le \bar{X}-\mu \le z_{\alpha/2}\frac{\sigma}{\sqrt{n}}\right) = 1-\alpha.
\]
Subtracting $\bar{X}$ from all three parts:
\[
P\!\left(-\bar{X} - z_{\alpha/2}\frac{\sigma}{\sqrt{n}} \le -\mu \le -\bar{X} + z_{\alpha/2}\frac{\sigma}{\sqrt{n}}\right) = 1-\alpha.
\]
Multiplying by $-1$ (which reverses the inequalities):
\[
P\!\left(\bar{X} - z_{\alpha/2}\frac{\sigma}{\sqrt{n}} \le \mu \le \bar{X} + z_{\alpha/2}\frac{\sigma}{\sqrt{n}}\right) = 1-\alpha. \qquad \blacksquare
\]

\begin{popfigure}
\begin{center}
\begin{tikzpicture}
\begin{axis}[
  width=12cm, height=6cm,
  axis lines=middle,
  xlabel={$z$}, ylabel={},
  xmin=-4, xmax=4, ymin=0, ymax=0.45,
  xtick={-1.96,0,1.96},
  xticklabels={$-z_{\alpha/2}$,$0$,$z_{\alpha/2}$},
  ytick=\empty,
  enlargelimits=false,
  clip=false
]
\addplot[domain=-4:4, samples=200, thick, popblue] {exp(-x^2/2)/sqrt(max(0,2*pi))};
\addplot[domain=-1.96:1.96, samples=100, fill=popblueL, draw=none] {exp(-x^2/2)/sqrt(max(0,2*pi))} \closedcycle;
\addplot[domain=-4:-1.96, samples=50, fill=poporange!50, draw=none] {exp(-x^2/2)/sqrt(max(0,2*pi))} \closedcycle;
\addplot[domain=1.96:4, samples=50, fill=poporange!50, draw=none] {exp(-x^2/2)/sqrt(max(0,2*pi))} \closedcycle;
\draw[dashed, popdark] (axis cs:-1.96,0) -- (axis cs:-1.96,0.242);
\draw[dashed, popdark] (axis cs:1.96,0) -- (axis cs:1.96,0.242);
\node[popdark] at (axis cs:0,0.2) {$1-\alpha$};
\node[popdark] at (axis cs:-2.9,0.03) {$\alpha/2$};
\node[popdark] at (axis cs:2.9,0.03) {$\alpha/2$};
\end{axis}
\end{tikzpicture}
\legende{Standard normal curve: the central region of area $1-\alpha$ between the critical values $\pm z_{\alpha/2}$, with tails of area $\alpha/2$ each.}
\end{center}
\end{popfigure}

\begin{aretenir}[Standard critical values]
\begin{center}
\begin{tabular}{|c|c|c|c|}
\hline
\rowcolor{popblueL} Confidence level $1-\alpha$ & $\alpha$ & $\alpha/2$ & $z_{\alpha/2}$ \\
\hline
$90\%$ & $0.10$ & $0.05$ & $1.645$ \\
\hline
$95\%$ & $0.05$ & $0.025$ & $1.960$ \\
\hline
$99\%$ & $0.01$ & $0.005$ & $2.576$ \\
\hline
\end{tabular}
\end{center}
These three values must be known by heart for the GCE examination.
\end{aretenir}

\courssub{Constructing and Interpreting a Z-Interval}

\begin{methode}[Steps to construct a Z-interval for $\mu$ ($\sigma$ known)]
\begin{enumerate}
\item \textbf{Identify the parameter:} the population mean $\mu$.
\item \textbf{Check the conditions:} the sample is random; $\sigma$ is known; the population is normal \emph{or} $n$ is large (rule of thumb $n \ge 30$ for the Central Limit Theorem).
\item \textbf{Compute} the sample mean $\bar{x}$.
\item \textbf{Choose} the confidence level $1-\alpha$ and read off $z_{\alpha/2}$ from normal tables.
\item \textbf{Compute the margin of error} $E = z_{\alpha/2}\,\sigma/\sqrt{n}$.
\item \textbf{Write the interval} $\bar{x} \pm E$ and \textbf{interpret it in context}: ``We are $95\%$ confident that the true mean \dots lies between \dots and \dots''.
\end{enumerate}
\end{methode}

\begin{exoresolu}[Masses of cement bags]
A factory in Douala produces bags of cement whose mass is known from long experience to have standard deviation $\sigma = 0.4$ kg. A random sample of $64$ bags gives a mean mass of $\bar{x} = 49.7$ kg. Construct a $95\%$ confidence interval for the true mean mass $\mu$ and interpret it.
\tcblower
\textbf{Solution.} The sample is random, $\sigma = 0.4$ is known and $n = 64 \ge 30$, so the Z-interval applies (Central Limit Theorem).
\begin{align*}
E &= z_{0.025}\,\frac{\sigma}{\sqrt{n}} = 1.96 \times \frac{0.4}{\sqrt{64}} = 1.96 \times \frac{0.4}{8} = 1.96 \times 0.05 = 0.098.
\end{align*}
Hence the $95\%$ confidence interval is
\[
49.7 \pm 0.098 = (49.602,\ 49.798)\ \text{kg}.
\]
\textbf{Interpretation:} we are $95\%$ confident that the true mean mass of all bags produced by the factory lies between $49.602$ kg and $49.798$ kg. Equivalently, the method used would capture $\mu$ in $95\%$ of all possible samples of $64$ bags.
\end{exoresolu}

\courssub{The Long-Run Frequency Interpretation}

What does ``$95\%$ confident'' really mean? Imagine repeating the whole experiment many times: each time we draw a new random sample of size $n$, compute $\bar{x}$, and build the interval $\bar{x} \pm E$. Because $\bar{x}$ varies from sample to sample, the interval moves around; but the true $\mu$ stays fixed. In the long run, about $95\%$ of these intervals will contain $\mu$, and about $5\%$ will miss it. The confidence level is a property of the \emph{method}, not of any single interval.

\begin{experience}[Simulating the meaning of 95\%]
Suppose $\mu = 50$ and $\sigma/\sqrt{n} = 1$. Four samples give $\bar{x} = 49.3,\ 50.6,\ 51.4,\ 49.8$. With $E = 1.96$, the intervals are $(47.34, 51.26)$, $(48.64, 52.56)$, $(49.44, 53.36)$ and $(47.84, 51.76)$. All four contain $\mu = 50$. Had we drawn a sample with $\bar{x} = 52.1$, the interval $(50.14, 54.06)$ would have \emph{missed} $\mu$ — this happens about once in twenty samples at the $95\%$ level, and we never know which sample we have.
\end{experience}

\courssub{What Controls the Width of the Interval?}

The width of the interval is
\[
2E = 2\,z_{\alpha/2}\,\frac{\sigma}{\sqrt{n}}.
\]
Three factors act on it:
\begin{itemize}
\item \textbf{Sample size $n$:} the width is proportional to $1/\sqrt{n}$. To \emph{halve} the width you must \emph{quadruple} the sample size. Larger samples give sharper estimates.
\item \textbf{Confidence level $1-\alpha$:} a higher level uses a larger $z_{\alpha/2}$, so the interval is \emph{wider}. Greater confidence costs precision.
\item \textbf{Population variability $\sigma$:} the more dispersed the population, the wider the interval. We cannot change $\sigma$; we can only compensate through $n$.
\end{itemize}

\begin{popfigure}
\begin{center}
\begin{tikzpicture}[scale=1]
\draw[->, popdark] (0,0) -- (12.5,0) node[right] {$\mu$};
\draw[thick, popdark] (6.25,-0.12) -- (6.25,0.12);
\node[below, popdark] at (6.25,-0.1) {$\bar{x}$};
% 90%
\draw[very thick, popgreen] (4.43,0.7) -- (8.07,0.7);
\fill[popgreen] (4.43,0.7) circle (1.5pt); \fill[popgreen] (8.07,0.7) circle (1.5pt);
\node[left, popgreen] at (4.3,0.7) {$90\%$};
% 95%
\draw[very thick, popblue] (4.09,1.5) -- (8.41,1.5);
\fill[popblue] (4.09,1.5) circle (1.5pt); \fill[popblue] (8.41,1.5) circle (1.5pt);
\node[left, popblue] at (3.95,1.5) {$95\%$};
% 99%
\draw[very thick, poporange] (3.41,2.3) -- (9.09,2.3);
\fill[poporange] (3.41,2.3) circle (1.5pt); \fill[poporange] (9.09,2.3) circle (1.5pt);
\node[left, poporange] at (3.25,2.3) {$99\%$};
\end{tikzpicture}
\legende{Nested confidence intervals from the same data ($\bar{x} = 49.7$, $\sigma/\sqrt{n} = 0.05$): higher confidence levels give wider intervals, all centred on $\bar{x}$.}
\end{center}
\end{popfigure}

\begin{exemplebox}[Effect of the confidence level]
Using the cement data ($\bar{x} = 49.7$, $\sigma/\sqrt{n} = 0.05$):
\begin{align*}
90\%:&\quad 49.7 \pm 1.645(0.05) = (49.618,\ 49.782),\\
95\%:&\quad 49.7 \pm 1.960(0.05) = (49.602,\ 49.798),\\
99\%:&\quad 49.7 \pm 2.576(0.05) = (49.571,\ 49.829).
\end{align*}
The $99\%$ interval is about $57\%$ wider than the $90\%$ interval: extra confidence is paid for with extra width.
\end{exemplebox}

\courssub{One-Sided Confidence Bounds}

Sometimes only one direction matters. A consumer association wants to be sure the mean mass is \emph{at least} some value; an engineer wants an \emph{upper} bound for the mean level of a pollutant in the Wouri river. In such cases we use a \cle{one-sided} confidence bound, placing the whole of $\alpha$ in a single tail.

\begin{propriete}[One-sided confidence bounds for $\mu$ ($\sigma$ known)]
A $100(1-\alpha)\%$ \textbf{lower} confidence bound for $\mu$ is
\[
\mu \ge \bar{x} - z_{\alpha}\,\frac{\sigma}{\sqrt{n}},
\]
and a $100(1-\alpha)\%$ \textbf{upper} confidence bound is
\[
\mu \le \bar{x} + z_{\alpha}\,\frac{\sigma}{\sqrt{n}}.
\]
Note carefully: one-sided bounds use $z_{\alpha}$ (e.g. $1.645$ for $95\%$), \emph{not} $z_{\alpha/2}$.
\end{propriete}

\begin{exoresolu}[A lower bound for the mean mass]
With the cement data ($\bar{x} = 49.7$ kg, $\sigma = 0.4$ kg, $n = 64$), find a $95\%$ lower confidence bound for $\mu$.
\tcblower
\textbf{Solution.} For a one-sided $95\%$ bound, $z_{0.05} = 1.645$:
\[
\mu \ge 49.7 - 1.645 \times \frac{0.4}{8} = 49.7 - 0.0823 = 49.618\ \text{kg}.
\]
We are $95\%$ confident that the true mean mass is at least $49.618$ kg. Note that this is exactly the lower endpoint of the $90\%$ two-sided interval — a useful check.
\end{exoresolu}

\begin{attention}[Common mistakes with Z-intervals]
\begin{itemize}
\item Using $z_{\alpha/2}$ for a one-sided bound, or $z_{\alpha}$ for a two-sided interval. Decide first whether the question is one- or two-sided.
\item Forgetting the $\sqrt{n}$: the standard error is $\sigma/\sqrt{n}$, not $\sigma/n$ or $\sigma$.
\item Writing the interval as $\mu \pm z_{\alpha/2}\sigma/\sqrt{n}$: the interval is centred on the \emph{statistic} $\bar{x}$, since $\mu$ is unknown.
\item Claiming ``there is a $95\%$ probability that $\mu$ lies in this interval'' after computing it: the $95\%$ refers to the long-run success rate of the method.
\item Believing a wider interval means a worse study: at fixed $n$, a wider interval simply reflects a higher confidence level.
\end{itemize}
\end{attention}

\begin{exercice}[Practice]
A random sample of $100$ pupils in a large school has a mean height of $162.4$ cm. Assume heights have known standard deviation $\sigma = 8$ cm.
\begin{enumerate}
\item Construct a $95\%$ confidence interval for the mean height $\mu$ of all pupils.
\item Construct a $99\%$ confidence interval and comment on the change in width.
\item Find a $95\%$ upper confidence bound for $\mu$.
\item How large a sample would halve the width of the $95\%$ interval?
\end{enumerate}
\end{exercice}

\begin{corrige}[Exercise 1]
\begin{enumerate}
\item $E = 1.96 \times 8/\sqrt{100} = 1.96 \times 0.8 = 1.568$, so the interval is $162.4 \pm 1.568 = (160.832,\ 163.968)$ cm.
\item $E = 2.576 \times 0.8 = 2.0608$, giving $(160.339,\ 164.461)$ cm. The interval is wider because greater confidence requires a larger critical value.
\item $\mu \le 162.4 + 1.645 \times 0.8 = 162.4 + 1.316 = 163.716$ cm.
\item Width is proportional to $1/\sqrt{n}$; to halve it we need $4n = 400$ pupils.
\end{enumerate}
\end{corrige}

\begin{savaistu}[Why 95\%?]
The choice of $95\%$ is a convention, not a law of nature: it became standard in the twentieth century through the influence of statisticians such as Jerzy Neyman, who developed the modern theory of confidence intervals in the 1930s. In medical research $99\%$ or even higher levels are common when lives are at stake, while quick market surveys in Cameroon often work at $90\%$ to keep sample sizes — and costs — manageable.
\end{savaistu}

\begin{aretenir}[Key points of this section]
\begin{itemize}
\item A confidence interval gives a \emph{range} of plausible values for $\mu$ plus a confidence level $1-\alpha$.
\item For $\sigma$ known: $\bar{x} \pm z_{\alpha/2}\,\sigma/\sqrt{n}$; memorise $z_{0.05} = 1.645$, $z_{0.025} = 1.96$, $z_{0.005} = 2.576$.
\item Width decreases with $\sqrt{n}$, increases with the confidence level and with $\sigma$.
\item One-sided bounds use $z_{\alpha}$ and keep all of $\alpha$ in one tail.
\item The confidence level describes the long-run behaviour of the \emph{method}, never the probability content of one computed interval.
\end{itemize}
\end{aretenir}

In the next section we remove the unrealistic assumption that $\sigma$ is known. This leads naturally to Student's $t$-distribution for the mean, and to a parallel theory of confidence intervals for a population \emph{proportion}.

\coursec{Confidence Intervals for the Mean (Unknown Variance) and for a Proportion}

In the previous section we constructed confidence intervals for the population mean $\mu$ assuming the population variance $\sigma^2$ was known. In practice, $\sigma^2$ is almost never known. We must estimate it from the sample using the sample variance $s^2$. This extra layer of estimation changes the sampling distribution of the standardised mean: instead of a standard normal distribution, we obtain \textbf{Student's $t$-distribution}. This section develops the $t$-interval for a mean, then turns to estimating a population proportion $p$, where the Wald interval is standard but the Wilson (score) interval often performs better, especially for small samples or proportions near $0$ or $1$.

\courssub{From Known to Unknown Variance: The Need for a New Distribution}

Recall that for a random sample $X_1,\dots,X_n$ from a normal population $N(\mu,\sigma^2)$, the statistic
\[
Z = \frac{\bar{X} - \mu}{\sigma/\sqrt{n}} \sim N(0,1).
\]
When $\sigma$ is unknown we replace it by the sample standard deviation
\[
S = \sqrt{\frac{1}{n-1}\sum_{i=1}^n (X_i - \bar{X})^2}.
\]
The resulting statistic
\[
T = \frac{\bar{X} - \mu}{S/\sqrt{n}}
\]
no longer follows a standard normal distribution. Its distribution has heavier tails because $S$ varies from sample to sample, adding extra uncertainty. This distribution was derived by William Sealy Gosset (publishing under the pseudonym ``Student'') in 1908 while working for Guinness Breweries in Dublin; he needed a method to monitor the quality of stout using very small samples. It is called the \textbf{Student's $t$-distribution}.

\begin{savaistu}[William Gosset and the Guinness Connection]
Gosset was a chemist and statistician hired by Arthur Guinness \& Son to apply scientific methods to brewing. He discovered the $t$-distribution while working with small samples of barley and hops (often $n=4$). Guinness considered the work a trade secret but allowed him to publish under the pen name ``Student'' in 1908. The $t$-test revolutionised industrial quality control and remains the cornerstone of small-sample inference.
\end{savaistu}

\courssub{The Student's $t$-Distribution}

\begin{definition}[Student's $t$-statistic and degrees of freedom]
Let $X_1,\dots,X_n$ be a random sample from a normal distribution $N(\mu,\sigma^2)$. The statistic
\[
T = \frac{\bar{X} - \mu}{S/\sqrt{n}}
\]
follows a \textbf{Student's $t$-distribution with $\nu = n-1$ degrees of freedom}, denoted $t_{\nu}$ or $t_{n-1}$.
The number of degrees of freedom $\nu$ equals the number of independent pieces of information used to estimate $\sigma$ (the $n$ observations minus the one constraint $\sum_{i=1}^n (X_i-\bar{X})=0$).
\end{definition}

\begin{propriete}[Properties of the $t$-distribution]
\begin{enumerate}
\item It is symmetric about $0$, bell-shaped, but with \textbf{heavier tails} than the standard normal.
\item As $\nu \to \infty$, the $t_{\nu}$ distribution converges to $N(0,1)$.
\item The variance of $t_{\nu}$ is $\frac{\nu}{\nu-2}$ for $\nu > 2$ (undefined for $\nu \le 2$), which is always $>1$, confirming the heavier tails.
\item Critical values $t_{\alpha/2,\nu}$ are such that $P(T > t_{\alpha/2,\nu}) = \alpha/2$. They are larger in absolute value than the corresponding $z_{\alpha/2}$.
\end{enumerate}
\end{propriete}

\begin{popfigure}
\centering
\begin{tikzpicture}
\begin{axis}[
    width=0.9\linewidth, height=7cm,
    axis lines=middle,
    xlabel={$t$}, ylabel={Density},
    xmin=-4, xmax=4, ymin=0, ymax=0.42,
    xtick={-3,-2,-1,0,1,2,3},
    ytick={0,0.1,0.2,0.3,0.4},
    legend style={at={(0.98,0.98)}, anchor=north east, font=\small},
    samples=200, smooth,
    domain=-4:4,
    clip=false
]
% Standard normal
\addplot[popdark, thick] {1/sqrt(max(0,2*pi))*exp(-x^2/2)};
\addlegendentry{$N(0{,}1)$}
% t with df=2
\addplot[poporange, thick, dashed] {1/(2*sqrt(max(0,2)))*(1+x^2/2)^(-1.5)};
\addlegendentry{$t_2$}
% t with df=5
\addplot[popgreen, thick, dotted] {8/(3*sqrt(max(0,5))*pi)*(1+x^2/5)^(-3)};
\addlegendentry{$t_5$}
% t with df=30
\addplot[poppurple, thick, dashdotted] {0.3956*(1+x^2/30)^(-15.5)};
\addlegendentry{$t_{30}$}
\end{axis}
\end{tikzpicture}
\legende{Probability density functions of $t$-distributions with $\nu=2,5,30$ and the standard normal $N(0,1)$. As $\nu$ increases, the $t$-distribution approaches the normal.}
\end{popfigure}

\courssub{Confidence Interval for $\mu$ when $\sigma$ is Unknown}

If the population is normal, or if $n$ is large enough for the Central Limit Theorem to apply, a $(1-\alpha)100\%$ confidence interval for $\mu$ is
\[
\bar{x} \pm t_{\alpha/2,\,n-1} \cdot \frac{s}{\sqrt{n}}.
\]

\begin{aretenir}[Confidence interval for $\mu$ ($\sigma$ unknown)]
\[
\boxed{\bar{x} - t_{\alpha/2,\,n-1}\,\frac{s}{\sqrt{n}} \;<\; \mu \;<\; \bar{x} + t_{\alpha/2,\,n-1}\,\frac{s}{\sqrt{n}}}
\]
where
\begin{itemize}
\item $\bar{x}$ = sample mean,
\item $s$ = sample standard deviation,
\item $n$ = sample size,
\item $t_{\alpha/2,\,n-1}$ = upper $\alpha/2$ critical value of the $t$-distribution with $n-1$ degrees of freedom.
\end{itemize}
\end{aretenir}

\begin{definition}[Standard error of the mean (estimated)]
The quantity $\displaystyle \widehat{\sigma}_{\bar{X}} = \frac{s}{\sqrt{n}}$ is called the \textbf{estimated standard error of the mean}. It estimates the standard deviation of the sampling distribution of $\bar{X}$.
\end{definition}

\begin{methode}[Constructing a $t$-confidence interval for $\mu$]
\begin{enumerate}
\item Verify the conditions: random sample; normality (if $n<30$) or large $n$ ($\ge 30$).
\item Compute the sample mean $\bar{x}$ and sample standard deviation $s$.
\item Determine the degrees of freedom $\nu = n-1$ and the confidence level $1-\alpha$.
\item Find the critical value $t_{\alpha/2,\nu}$ from the $t$-table (or software).
\item Calculate the margin of error $E = t_{\alpha/2,\nu} \cdot \dfrac{s}{\sqrt{n}}$.
\item Write the interval: $\bar{x} \pm E$ or $(\bar{x}-E,\; \bar{x}+E)$.
\item Interpret the interval in the context of the problem.
\end{enumerate}
\end{methode}

\courssub{Conditions for Validity of the $t$-Interval}

The exact $t$-interval relies on the assumption that the sample comes from a \textbf{normal population}. In practice we check this assumption.

\begin{methode}[Checking the normality assumption for the $t$-interval]
\begin{enumerate}
\item \textbf{Small sample ($n < 30$):} The population must be approximately normal. Check with:
      \begin{itemize}
      \item A \textbf{Q--Q plot} (quantile-quantile plot): points should lie close to a straight line.
      \item A formal test such as the \textbf{Shapiro--Wilk test} (if $p$-value $> 0.05$, do not reject normality).
      \item A histogram or boxplot for gross non-normality (skewness, outliers).
      \end{itemize}
\item \textbf{Large sample ($n \ge 30$):} By the Central Limit Theorem, $\bar{X}$ is approximately normal even if the population is not. The $t$-interval is robust to moderate departures from normality.
\item \textbf{If normality fails for small $n$:} Consider non-parametric methods (e.g., bootstrap confidence interval, sign test) or transform the data (e.g., log transformation).
\end{enumerate}
\end{methode}

\courssub{Confidence Interval for a Population Proportion $p$}

Let $X$ be the number of ``successes'' in a random sample of size $n$ from a large population (or with replacement). The sample proportion is $\hat{p} = X/n$. For large $n$, by the Central Limit Theorem,
\[
\frac{\hat{p} - p}{\sqrt{p(1-p)/n}} \approx N(0,1).
\]
Replacing the unknown $p$ in the standard error by its estimate $\hat{p}$ gives the \textbf{Wald interval} (also called the normal approximation interval).

\begin{propriete}[Wald confidence interval for $p$]
\[
\boxed{\hat{p} \pm z_{\alpha/2}\,\sqrt{\frac{\hat{p}(1-\hat{p})}{n}}}
\]
provided $n\hat{p} \ge 5$ and $n(1-\hat{p}) \ge 5$ (some textbooks use $10$ as a safer threshold).
\end{propriete}

\begin{attention}[Common mistake: Wald interval with small $n$ or extreme $\hat{p}$]
The Wald interval can have actual coverage probability far below the nominal level when $n$ is small or $\hat{p}$ is close to $0$ or $1$. It can even produce limits outside $[0,1]$. \textbf{Always check $n\hat{p} \ge 5$ and $n(1-\hat{p}) \ge 5$}. If not satisfied, use the Wilson (score) interval or the exact Clopper--Pearson interval.
\end{attention}

\courssub{The Wilson (Score) Interval}

The Wilson interval is derived by inverting the score test. It solves for $p$ in
\[
\frac{\hat{p} - p}{\sqrt{p(1-p)/n}} = \pm z_{\alpha/2}.
\]
The resulting interval has better coverage properties, especially for small $n$ or extreme $p$.

\begin{propriete}[Wilson (score) confidence interval for $p$]
\[
\boxed{
\frac{\hat{p} + \frac{z^2}{2n} \;\pm\; z\sqrt{\frac{\hat{p}(1-\hat{p})}{n} + \frac{z^2}{4n^2}}}{1 + \frac{z^2}{n}}
}
\]
where $z = z_{\alpha/2}$. The interval is always contained in $[0,1]$.
\end{propriete}

\begin{popfigure}
\centering
\begin{tikzpicture}
\begin{axis}[
    width=0.95\linewidth, height=7cm,
    axis lines=middle,
    xlabel={$\hat{p}$}, ylabel={Interval half-width},
    xmin=0, xmax=1, ymin=0, ymax=0.55,
    xtick={0,0.2,0.4,0.5,0.6,0.8,1},
    ytick={0,0.1,0.2,0.3,0.4,0.5},
    legend style={at={(0.5,1.02)}, anchor=south, font=\small, legend columns=2},
    samples=101, smooth,
    domain=0:1,
    clip=false
]
\def\n{20}
\def\z{1.96}
% Wald half-width
\addplot[poporange, thick] {\z*sqrt(max(0,x*(1-x)/\n))};
\addlegendentry{Wald half-width}
% Wilson half-width (upper minus centre)
\addplot[popgreen, thick, dashed] {
    ( (x + \z^2/(2*\n) + \z*sqrt(x*(1-x)/\n + \z^2/(4*\n^2)) )/(1+\z^2/\n)
    - (x + \z^2/(2*\n))/(1+\z^2/\n) )
};
\addlegendentry{Wilson half-width}
% Centre of Wilson (for reference)
\addplot[popblue, thin, dotted] { (x + \z^2/(2*\n))/(1+\z^2/\n) };
\addlegendentry{Wilson centre}
\end{axis}
\end{tikzpicture}
\legende{Comparison of half-widths of the 95\% Wald and Wilson intervals for $n=20$. The Wilson interval is asymmetric and never exceeds $[0,1]$.}
\end{popfigure}

\courssub{Worked Examples}

\begin{exoresolu}[Example 1: $t$-interval for a mean]
A food safety laboratory in Douala measures the lead content (in mg/kg) in a random sample of $10$ cans of imported sardines. The data are:
\[
0.12,\; 0.15,\; 0.11,\; 0.14,\; 0.13,\; 0.16,\; 0.10,\; 0.14,\; 0.13,\; 0.15.
\]
Construct a $95\%$ confidence interval for the true mean lead content $\mu$. Assume normality.
\tcblower
\textbf{Solution.}
\begin{enumerate}
\item Compute sample statistics:
      $n = 10$,
      $\bar{x} = \frac{1.33}{10} = 0.133$,
      $s^2 = \frac{1}{9}\sum (x_i - 0.133)^2 = \frac{0.00321}{9} = 0.0003567$,
      $s = \sqrt{0.0003567} \approx 0.01889$.
\item Degrees of freedom: $\nu = n-1 = 9$.
\item For $95\%$ confidence, $\alpha = 0.05$, $\alpha/2 = 0.025$. From $t$-tables: $t_{0.025,9} = 2.262$.
\item Standard error: $\displaystyle \frac{s}{\sqrt{n}} = \frac{0.01889}{\sqrt{10}} \approx 0.00597$.
\item Margin of error: $E = 2.262 \times 0.00597 \approx 0.0135$.
\item Confidence interval:
      \[
      0.133 \pm 0.0135 \quad\Rightarrow\quad (0.1195,\; 0.1465).
      \]
\end{enumerate}
We are $95\%$ confident that the true mean lead content lies between $0.120$ and $0.147$ mg/kg.
\end{exoresolu}

\begin{exoresolu}[Example 2: Choosing between $Z$ and $t$]
A researcher wants a $99\%$ confidence interval for the mean height of a certain variety of plant.
\begin{enumerate}
\item In a pilot study of $50$ plants, the sample mean is $45.2$ cm and the sample standard deviation is $3.8$ cm. The population variance is unknown. Which distribution should be used?
\item If the population standard deviation were known to be $3.8$ cm, which distribution would be used?
\end{enumerate}
\tcblower
\textbf{Solution.}
\begin{enumerate}
\item \textbf{Unknown $\sigma$:} Use the $t$-distribution with $\nu = 49$. For $99\%$ confidence, $\alpha/2 = 0.005$. $t_{0.005,49} \approx 2.680$ (from tables or software). The interval is $45.2 \pm 2.680 \times \frac{3.8}{\sqrt{50}}$.
\item \textbf{Known $\sigma$:} Use the standard normal distribution. $z_{0.005} = 2.576$. The interval would be $45.2 \pm 2.576 \times \frac{3.8}{\sqrt{50}}$.
\end{enumerate}
\textbf{Note:} Even for $n=50$, the $t$ critical value ($2.680$) is slightly larger than the $z$ critical value ($2.576$), reflecting the extra uncertainty from estimating $\sigma$. As $n$ grows, the difference vanishes.
\end{exoresolu}

\begin{exoresolu}[Example 3: Proportion intervals -- Wald vs. Wilson]
In a survey of $25$ randomly selected voters in a constituency, $18$ say they support Candidate A. Compute a $95\%$ confidence interval for the true proportion $p$ of supporters using both the Wald and Wilson methods.
\tcblower
\textbf{Solution.}
$n = 25$, $X = 18$, $\hat{p} = 18/25 = 0.72$. $z_{0.025} = 1.96$.

\textbf{Wald interval:}
Check conditions: $n\hat{p} = 18 \ge 5$, $n(1-\hat{p}) = 7 \ge 5$ (barely).
\[
\text{SE} = \sqrt{\frac{0.72 \times 0.28}{25}} = \sqrt{0.008064} \approx 0.0898.
\]
Margin of error: $E = 1.96 \times 0.0898 \approx 0.176$.
Interval: $0.72 \pm 0.176 \quad\Rightarrow\quad (0.544,\; 0.896)$.

\textbf{Wilson interval:}
$z^2 = 1.96^2 = 3.8416$.
Centre: $\displaystyle \frac{0.72 + \frac{3.8416}{50}}{1 + \frac{3.8416}{25}} = \frac{0.72 + 0.07683}{1 + 0.15366} = \frac{0.79683}{1.15366} \approx 0.6907$.
Half-width: $\displaystyle \frac{1.96 \sqrt{\frac{0.72\times0.28}{25} + \frac{3.8416}{2500}}}{1.15366}
= \frac{1.96 \sqrt{0.008064 + 0.0015366}}{1.15366}
= \frac{1.96 \times 0.0980}{1.15366} \approx 0.1664$.
Interval: $0.6907 \pm 0.1664 \quad\Rightarrow\quad (0.5243,\; 0.8571)$.

\textbf{Comparison:} The Wilson interval is shifted slightly toward $0.5$ and is narrower on the upper side, staying well within $[0,1]$. For this small sample, the Wilson interval is preferred.
\end{exoresolu}

\courssub{Summary of Interval Selection}

\begin{center}
\begin{tabularx}{\linewidth}{|l|X|X|}\hline
\rowcolor{popblueL}
\textbf{Parameter} & \textbf{Condition} & \textbf{Interval Formula} \\ \hline
$\mu$ & $\sigma$ known, normal pop. or large $n$ & $\bar{x} \pm z_{\alpha/2}\,\dfrac{\sigma}{\sqrt{n}}$ \\ \hline
$\mu$ & $\sigma$ unknown, normal pop. & $\bar{x} \pm t_{\alpha/2,\,n-1}\,\dfrac{s}{\sqrt{n}}$ \\ \hline
$\mu$ & $\sigma$ unknown, non-normal, large $n$ ($\ge 30$) & $\bar{x} \pm t_{\alpha/2,\,n-1}\,\dfrac{s}{\sqrt{n}}$ (approx.) \\ \hline
$p$ & Large $n$, $n\hat{p}\ge5$, $n(1-\hat{p})\ge5$ & $\hat{p} \pm z_{\alpha/2}\sqrt{\dfrac{\hat{p}(1-\hat{p})}{n}}$ (Wald) \\ \hline
$p$ & Small $n$ or extreme $\hat{p}$ & Wilson (score) interval (preferred) \\ \hline
\end{tabularx}
\end{center}

\begin{aretenir}[Key points to remember]
\begin{itemize}
\item When $\sigma$ is unknown, use the $t$-distribution with $n-1$ degrees of freedom. The $t$-interval is exact for normal populations and approximate for large $n$.
\item The $t$-distribution has heavier tails than the normal; critical values $t_{\alpha/2,\nu} > z_{\alpha/2}$.
\item For a proportion, the Wald interval is simple but can be inaccurate for small $n$ or $\hat{p}$ near $0$ or $1$. Check $n\hat{p}\ge5$ and $n(1-\hat{p})\ge5$.
\item The Wilson (score) interval has better coverage properties and is always within $[0,1]$. Use it when Wald conditions are not met.
\item Always verify assumptions (normality for small-sample $t$, success/failure counts for proportions) before computing an interval.
\end{itemize}
\end{aretenir}

\begin{exercice}[Practice exercises]
\begin{itemize}
\item A random sample of $16$ bags of cement from a factory in Figuil has a mean weight of $49.8$ kg and a standard deviation of $0.6$ kg. Assuming normality, find a $99\%$ confidence interval for the true mean weight.
\item In a study of malaria prevalence, $42$ out of $120$ children tested positive. Compute a $95\%$ confidence interval for the true prevalence using (a) the Wald method, (b) the Wilson method. Comment on any difference.
\item A quality control engineer measures the diameter of $8$ ball bearings (in mm): $10.02, 9.98, 10.01, 10.03, 9.99, 10.00, 10.02, 9.97$. Construct a $90\%$ confidence interval for the mean diameter. Check the normality assumption with a Q--Q plot (describe what you would look for).
\item Explain why the $t$-interval is wider than the $Z$-interval for the same data and confidence level.
\item A survey of $30$ households finds that $3$ have access to fibre internet. Can you use the Wald interval to estimate the proportion? If not, compute the Wilson $95\%$ interval.
\end{itemize}
\end{exercice}

\begin{corrige}[Exercise 1]
$n=16$, $\bar{x}=49.8$, $s=0.6$, $\nu=15$. $t_{0.005,15}=2.947$. $\text{SE}=0.6/\sqrt{16}=0.15$. $E=2.947\times0.15=0.442$. Interval: $49.8\pm0.442 = (49.358,\;50.242)$ kg.
\end{corrige}

\begin{corrige}[Exercise 2]
$\hat{p}=42/120=0.35$. 
(a) Wald: $\text{SE}=\sqrt{0.35\times0.65/120}=0.0435$. $E=1.96\times0.0435=0.085$. Interval: $(0.265,\;0.435)$. 
(b) Wilson: $z^2=3.8416$. Centre $=\frac{0.35+3.8416/240}{1+3.8416/120}=\frac{0.3660}{1.0320}=0.3547$. Half-width $=\frac{1.96\sqrt{0.35\times0.65/120+3.8416/57600}}{1.0320}=\frac{1.96\times0.0443}{1.0320}=0.0841$. Interval: $(0.2706,\;0.4388)$. 
Comment: The intervals are very similar because $n$ is moderately large and $\hat{p}$ is not extreme. The Wilson interval is slightly shifted and asymmetric.
\end{corrige}

\begin{corrige}[Exercise 3]
$n=8$, $\bar{x}=10.0025$, $s=0.0212$, $\nu=7$. $t_{0.05,7}=1.895$. $\text{SE}=0.0212/\sqrt{8}=0.0075$. $E=1.895\times0.0075=0.0142$. Interval: $(9.988,\;10.017)$ mm. Q--Q plot: plot ordered sample quantiles vs. theoretical normal quantiles; look for points lying approximately on a straight line.
\end{corrige}

\begin{corrige}[Exercise 4]
The $t$-interval uses $s$ instead of $\sigma$, adding estimation variability. The $t$-distribution has heavier tails, so $t_{\alpha/2,\nu} > z_{\alpha/2}$, giving a larger margin of error.
\end{corrige}

\begin{corrige}[Exercise 5]
$n=30$, $\hat{p}=0.1$. $n\hat{p}=3 < 5$, Wald not recommended. Wilson: $z^2=3.8416$. Centre $=\frac{0.1+3.8416/60}{1+3.8416/30}=\frac{0.1640}{1.1281}=0.1454$. Half-width $=\frac{1.96\sqrt{0.1\times0.9/30+3.8416/3600}}{1.1281}=\frac{1.96\times0.0638}{1.1281}=0.1108$. Interval: $(0.0346,\;0.2562)$.
\end{corrige}

\coursec{Margin of Error, Sample Size Determination and Applications}

In the previous section we learnt how to build confidence intervals for a population mean when the variance is unknown (using the $t$-distribution) and for a population proportion. In every case, the interval had the form
\[
\text{point estimate} \;\pm\; \text{(critical value)}\times\text{(standard error)}.
\]
The quantity added and subtracted on each side of the estimate is called the \cle{margin of error}. In this final section we reverse the question: instead of computing an interval from a given sample, we ask \emph{how large must the sample be} so that the margin of error does not exceed a prescribed value? This is one of the most practical questions in statistics, and it is exactly the question asked by a cocoa exporter in Douala, a polling organisation before an election, or a quality-control engineer in a brewery.

\courssub{The Margin of Error}

\textbf{Intuition.} When a newspaper announces ``$52\%$ of voters intend to vote for candidate A, with a margin of error of $3$ points'', it means that the true proportion probably lies in $[0.49,\,0.55]$. The margin of error measures the \emph{precision} of the estimate: the smaller it is, the more precise the estimate.

\begin{definition}[Margin of error]
The \cle{margin of error} $E$ of a confidence interval is the half-width of the interval, that is, the distance between the point estimate and either endpoint of the interval:
\[
E = (\text{critical value})\times(\text{standard error of the estimator}).
\]
For a population mean with known standard deviation $\sigma$, at confidence level $(1-\alpha)$:
\[
E = z_{\alpha/2}\,\dfrac{\sigma}{\sqrt{n}}.
\]
For a population mean with unknown standard deviation (estimated by $s$, large $n$):
\[
E = t_{\alpha/2,\,n-1}\,\dfrac{s}{\sqrt{n}}.
\]
For a population proportion:
\[
E = z_{\alpha/2}\,\sqrt{\dfrac{\hat{p}(1-\hat{p})}{n}}.
\]
\end{definition}

From the formula $E = z_{\alpha/2}\,\sigma/\sqrt{n}$ we can read off the three factors that control precision:

\begin{itemize}
\item \textbf{Confidence level.} Increasing the confidence level increases $z_{\alpha/2}$ (e.g.\ $1.645$ at $90\%$, $1.96$ at $95\%$, $2.576$ at $99\%$), hence \emph{increases} $E$. More confidence costs precision.
\item \textbf{Variability of the population.} A larger $\sigma$ gives a larger $E$: a heterogeneous population is harder to estimate.
\item \textbf{Sample size.} $E$ is proportional to $1/\sqrt{n}$. To \emph{halve} the margin of error, one must \emph{quadruple} the sample size. This square-root law is the fundamental economic fact of sampling.
\end{itemize}

\begin{attention}[The square-root law]
A very common error is to believe that doubling the sample size halves the margin of error. In fact $E\propto 1/\sqrt{n}$, so doubling $n$ only divides $E$ by $\sqrt{2}\approx 1.41$. To halve $E$ you need $4n$ observations; to divide $E$ by $10$ you need $100n$.
\end{attention}

\courssub{Determining the Sample Size for Estimating a Mean ($\sigma$ known)}

\textbf{Derivation.} Suppose we want to estimate the population mean $\mu$ with a margin of error at most $E$, at confidence level $(1-\alpha)$, and $\sigma$ is known (from previous studies or a pilot survey). We require
\[
z_{\alpha/2}\,\dfrac{\sigma}{\sqrt{n}} \le E
\quad\Longleftrightarrow\quad
\sqrt{n}\ \ge\ \dfrac{z_{\alpha/2}\,\sigma}{E}
\quad\Longleftrightarrow\quad
n \ge \left(\dfrac{z_{\alpha/2}\,\sigma}{E}\right)^{2}.
\]

\begin{propriete}[Sample size for estimating a mean]
To estimate a population mean $\mu$ with margin of error at most $E$ at confidence level $(1-\alpha)$, when $\sigma$ is known, the minimum sample size is
\[
\boxed{\,n = \left(\dfrac{z_{\alpha/2}\,\sigma}{E}\right)^{2}\,}
\]
always rounded \textbf{up} to the next integer. (This derivation is examinable.)
\end{propriete}

\begin{attention}[Always round up]
If the formula gives $n = 153.2$, we must take $n = 154$, never $153$. Rounding down would make the actual margin of error slightly \emph{larger} than the target $E$.
\end{attention}

\begin{attention}[What if $\sigma$ is unknown?]
The formula above requires $\sigma$. In practice, when $\sigma$ is unknown, one uses an estimate from a \emph{pilot study} or from previous data, and it is prudent to take a slightly \emph{larger} (conservative) value of $\sigma$. Strictly, with unknown variance the critical value should come from the $t$-distribution, which itself depends on $n$; the standard examination practice is to use the $z$-formula with the estimated $\sigma$, which is accurate for the moderate-to-large samples involved in planning.
\end{attention}

\begin{exoresolu}[Estimating the average weight of cocoa beans]
A cocoa cooperative in the Centre Region wants to estimate the mean weight $\mu$ of dried cocoa beans in a consignment. From previous seasons, the standard deviation of bean weights is known to be $\sigma = 0.30\ \mathrm{g}$. How many beans must be weighed so that the $95\%$ confidence interval for $\mu$ has a margin of error of at most $0.05\ \mathrm{g}$?
\tcblower
\textbf{Solution.} At the $95\%$ level, $z_{\alpha/2} = z_{0.025} = 1.96$. Then
\[
n = \left(\dfrac{1.96 \times 0.30}{0.05}\right)^{2} = (11.76)^{2} = 138.30.
\]
Rounding up, the cooperative must weigh $\boxed{n = 139}$ beans.

If a $99\%$ level were required instead, $z_{0.005} = 2.576$ and
\[
n = \left(\dfrac{2.576 \times 0.30}{0.05}\right)^{2} = (15.456)^{2} = 238.9 \;\Rightarrow\; n = 239.
\]
Greater confidence demands a substantially larger sample.
\end{exoresolu}

\courssub{Determining the Sample Size for Estimating a Proportion}

For a proportion, the margin of error is $E = z_{\alpha/2}\sqrt{\hat{p}(1-\hat{p})/n}$. Solving for $n$ exactly as before:
\[
z_{\alpha/2}\sqrt{\dfrac{\hat{p}(1-\hat{p})}{n}} \le E
\quad\Longleftrightarrow\quad
n \ge \dfrac{z_{\alpha/2}^{2}\,\hat{p}(1-\hat{p})}{E^{2}}.
\]

\begin{propriete}[Sample size for estimating a proportion]
To estimate a population proportion $p$ with margin of error at most $E$ at confidence level $(1-\alpha)$:
\[
\boxed{\,n = \dfrac{z_{\alpha/2}^{2}\,\hat{p}(1-\hat{p})}{E^{2}}\,}
\]
where $\hat{p}$ is a \emph{prior estimate} of $p$ (from a pilot study or previous survey). If no prior estimate is available, take $\hat{p} = 0.5$, which maximises $\hat{p}(1-\hat{p})$ and therefore gives the \emph{largest} (safest) sample size:
\[
n_{\max} = \dfrac{z_{\alpha/2}^{2}}{4E^{2}}.
\]
\end{propriete}

\textbf{Why $\hat{p}=0.5$ is the worst case.} The function $g(p) = p(1-p)$ is a downward parabola on $[0,1]$, with maximum $g(0.5) = 0.25$. Using $\hat{p}=0.5$ therefore guarantees that the sample size is sufficient whatever the true value of $p$.

\begin{exoresolu}[Voter intention in a constituency]
A polling organisation wants to estimate the proportion $p$ of registered voters in a Yaound\'e constituency who intend to vote for a given candidate, with a margin of error of $3$ percentage points at the $95\%$ level.
\begin{itemize}
\item[(a)] No prior information is available.
\item[(b)] A previous poll suggested $\hat{p} = 0.60$.
\end{itemize}
Find the required sample size in each case.
\tcblower
\textbf{Solution.} Here $E = 0.03$ and $z_{0.025} = 1.96$.

(a) With $\hat{p} = 0.5$:
\[
n = \dfrac{(1.96)^{2}(0.5)(0.5)}{(0.03)^{2}} = \dfrac{3.8416 \times 0.25}{0.0009} = \dfrac{0.9604}{0.0009} = 1067.1 \;\Rightarrow\; n = 1068.
\]

(b) With $\hat{p} = 0.60$:
\[
n = \dfrac{(1.96)^{2}(0.6)(0.4)}{(0.03)^{2}} = \dfrac{3.8416 \times 0.24}{0.0009} = 1024.4 \;\Rightarrow\; n = 1025.
\]
Using the prior estimate saves $43$ interviews. Note the famous rule of thumb visible here: for $E = 0.03$ at $95\%$, about $1068$ respondents are needed, and for $E = 0.02$, about $2401$.
\end{exoresolu}

\courssub{Finite Population Correction and Non-response}

The formulae above assume the population is infinite (or very large compared with the sample). When the population has a finite size $N$ and the sample is an appreciable fraction of it, the standard error is reduced.

\begin{definition}[Finite population correction factor]
If a sample of size $n$ is drawn \emph{without replacement} from a population of size $N$, the standard error of the mean is multiplied by the \cle{finite population correction factor} (FPC):
\[
\sqrt{\dfrac{N-n}{N-1}} \;\approx\; \sqrt{1-\dfrac{n}{N}}.
\]
The corrected sample size formula becomes
\[
n = \dfrac{n_0}{1 + \dfrac{n_0 - 1}{N}} \;\approx\; \dfrac{n_0}{1 + \dfrac{n_0}{N}},
\qquad
n_0 = \left(\dfrac{z_{\alpha/2}\,\sigma}{E}\right)^{2}.
\]
The correction is negligible when $n/N < 5\%$ and is usually ignored in that case.
\end{definition}

\textbf{Anticipated non-response.} In real surveys, some selected individuals refuse to answer or cannot be contacted. If a response rate $r$ (as a decimal) is anticipated, the initial sample size must be inflated:
\[
n_{\text{adjusted}} = \dfrac{n}{r}.
\]
For example, if the calculation gives $n = 1068$ and a response rate of $80\%$ is expected, one should contact $1068/0.80 = 1335$ persons.

\begin{methode}[Step-by-step sample size calculation]
\begin{enumerate}
\item \textbf{Identify the target}: mean or proportion? State the confidence level and read $z_{\alpha/2}$ ($1.645$, $1.96$ or $2.576$ for $90\%$, $95\%$, $99\%$).
\item \textbf{State the desired margin of error} $E$ in the same units as $\sigma$ (for a mean) or as a decimal (for a proportion).
\item \textbf{Gather the ingredients}: $\sigma$ (from a pilot study or past data), or $\hat{p}$ (take $0.5$ if unknown).
\item \textbf{Apply the formula}: $n = (z_{\alpha/2}\sigma/E)^2$ or $n = z_{\alpha/2}^2\hat{p}(1-\hat{p})/E^2$.
\item \textbf{Round up} to the next integer.
\item \textbf{Adjust} if necessary: finite population correction if $n/N \ge 5\%$; divide by the expected response rate to allow for non-response.
\end{enumerate}
\end{methode}

\begin{attention}[Underpowered studies]
Two mistakes ruin survey planning. First, \emph{ignoring non-response}: computing $n$ exactly and then losing $20\%$ of the respondents leaves the study with a margin of error larger than promised. Second, \emph{using an unrealistically small $\sigma$}: if the true variability is larger than assumed, the achieved precision is worse than planned. When in doubt, take a conservative (larger) value of $\sigma$, or use $\hat{p}=0.5$ for proportions. An \emph{underpowered} study wastes money because its conclusions are too imprecise to be useful.
\end{attention}

\courssub{Graphical View: Sample Size versus Margin of Error}

The figure below shows how the required sample size explodes as the demanded margin of error shrinks, for a mean with $\sigma = 0.30$ (in the units of the cocoa example), at three confidence levels.

\begin{popfigure}
\begin{tikzpicture}
\begin{axis}[
    width=0.85\linewidth, height=7cm,
    xlabel={Margin of error $E$},
    ylabel={Required sample size $n$},
    xmin=0.02, xmax=0.10, ymin=0, ymax=1500,
    legend pos=north east,
    legend style={font=\small},
    grid=major, grid style={popline, very thin},
    axis lines=left, thick
]
\addplot[domain=0.02:0.10, samples=80, very thick, popblue] {(1.645*0.30/x)^2};
\addlegendentry{$90\%$ ($z=1.645$)}
\addplot[domain=0.02:0.10, samples=80, very thick, poporange] {(1.96*0.30/x)^2};
\addlegendentry{$95\%$ ($z=1.96$)}
\addplot[domain=0.02:0.10, samples=80, very thick, poppurple] {(2.576*0.30/x)^2};
\addlegendentry{$99\%$ ($z=2.576$)}
\end{axis}
\end{tikzpicture}
\legende{Required sample size $n=(z_{\alpha/2}\sigma/E)^2$ as a function of the margin of error $E$, for $\sigma=0.30$. Halving $E$ multiplies $n$ by four; higher confidence levels require larger samples.}
\end{popfigure}

The next chart compares, at the $95\%$ level with $E=0.03$, the sample sizes required for a mean (with $\sigma=0.15$, so that $n=(1.96\times0.15/0.03)^2\approx 97$) and for a proportion (worst case $\hat p=0.5$, $n\approx 1068$).

\begin{popfigure}
\begin{tikzpicture}
\begin{axis}[
    width=0.75\linewidth, height=6.5cm,
    ybar, bar width=1.6cm,
    symbolic x coords={Mean ($\sigma=0.15$), Proportion ($\hat{p}=0.5$)},
    xtick=data,
    ylabel={Required sample size $n$},
    ymin=0, ymax=1250,
    nodes near coords,
    every node near coord/.append style={font=\small},
    axis lines=left, thick,
    x tick label style={font=\small, align=center}
]
\addplot[fill=popblue, draw=popdark] coordinates {(Mean ($\sigma=0.15$),97) (Proportion ($\hat{p}=0.5$),1068)};
\end{axis}
\end{tikzpicture}
\legende{Sample sizes needed at $95\%$ confidence with $E=0.03$: estimating a mean with small $\sigma$ is cheap; estimating a proportion in the worst case is expensive.}
\end{popfigure}

\courssub{Case Studies and Applications}

\begin{experience}[Case study 1 --- Quality control in a brewery]
A brewery in Douala bottles beer in nominal \SI{65}{cl} bottles. The filling machine is known to have a standard deviation $\sigma = \SI{4}{ml}$. The quality-control manager wants to estimate the mean fill volume to within $E = \SI{1}{ml}$ at $99\%$ confidence.
\tcblower
With $z_{0.005} = 2.576$:
\[
n = \left(\dfrac{2.576 \times 4}{1}\right)^{2} = (10.304)^{2} = 106.2 \;\Rightarrow\; n = 107 \text{ bottles}.
\]
If the manager later decides that $E = \SI{0.5}{ml}$ is needed, the sample size quadruples: $n = (2.576\times4/0.5)^2 = 424.8 \Rightarrow 425$ bottles. This illustrates the square-root law in practice: precision is expensive.
\end{experience}

\begin{experience}[Case study 2 --- Voter intention with finite population and non-response]
In a small municipality of $N = 12\,000$ registered voters, a candidate's team wants to estimate its vote share to within $4$ points at $95\%$ confidence, expecting a $75\%$ response rate.
\tcblower
Step 1: worst-case sample size ($\hat p = 0.5$, $E=0.04$):
\[
n_0 = \dfrac{(1.96)^2(0.25)}{(0.04)^2} = \dfrac{0.9604}{0.0016} = 600.25 \Rightarrow 601.
\]
Step 2: finite population correction (since $601/12000 \approx 5\%$):
\[
n = \dfrac{601}{1 + \frac{601}{12000}} = \dfrac{601}{1.0501} = 572.3 \Rightarrow 573.
\]
Step 3: adjust for non-response:
\[
n_{\text{contact}} = \dfrac{573}{0.75} = 764.
\]
The team should plan to contact about $764$ voters.
\end{experience}

\begin{savaistu}[Why polls of 1000 people work]
Whether a country has $5$ million or $50$ million voters, a well-conducted random sample of about $1068$ people estimates a proportion to within $\pm 3$ points at $95\%$ confidence. The required sample size depends on the \emph{precision} demanded, not on the population size (as long as the population is large). This is why international polling organisations routinely interview samples of only $1000$ to $2000$ people.
\end{savaistu}

\begin{exercice}[Practice]
\begin{enumerate}
\item A nutritionist wants to estimate the mean daily calorie intake of students at a university in Buea to within $50$ kcal at $95\%$ confidence. A pilot study gives $\sigma = 320$ kcal. Find the minimum sample size.
\item An NGO wants to estimate the proportion of households in a town that own a mosquito net, to within $2$ points at $90\%$ confidence, with no prior estimate. Find $n$.
\item In question 2, the town has $N = 8000$ households and a response rate of $85\%$ is expected. Adjust the sample size.
\item A factory manager claims: ``To double the precision of my estimate, I only need to double my sample.'' Comment.
\end{enumerate}
\end{exercice}

\begin{corrige}[Exercise 6.1]
\textbf{1.} $n = (1.96\times 320/50)^2 = (12.544)^2 = 157.35 \Rightarrow n = 158$ students.

\textbf{2.} $n = (1.645)^2(0.25)/(0.02)^2 = 0.6765/0.0004 = 1691.3 \Rightarrow n = 1692$.

\textbf{3.} FPC: $n = 1692/(1+1692/8000) = 1692/1.2115 = 1396.6 \Rightarrow 1397$. Non-response: $1397/0.85 = 1643.5 \Rightarrow 1644$ households to contact.

\textbf{4.} False. Since $E \propto 1/\sqrt{n}$, doubling the precision (halving $E$) requires \emph{four times} the sample size, not two.
\end{corrige}

\begin{aretenir}[Key points of Section 6]
\begin{itemize}
\item The \cle{margin of error} is $E = z_{\alpha/2}\sigma/\sqrt{n}$ (mean, $\sigma$ known) or $E = z_{\alpha/2}\sqrt{\hat p(1-\hat p)/n}$ (proportion).
\item Sample size for a mean: $n = \left(z_{\alpha/2}\sigma/E\right)^2$; for a proportion: $n = z_{\alpha/2}^2\hat p(1-\hat p)/E^2$, with $\hat p = 0.5$ if unknown. \textbf{Always round up.}
\item Halving $E$ multiplies $n$ by $4$; raising the confidence level increases $n$.
\item Apply the finite population correction when $n/N \ge 5\%$, and inflate $n$ by dividing by the expected response rate.
\item Never use an optimistic $\sigma$: an underpowered study produces useless, imprecise conclusions.
\end{itemize}
\end{aretenir}

\newpage

\coursec{Integration activity}

\courssub{Aims of the activity}
\begin{itemize}
    \item Integrate the complete chain of statistical inference studied in this chapter: sampling design, point estimation, confidence interval construction, margin of error control, and sample size determination.
    \item Apply theoretical concepts to a realistic Cameroonian policy scenario --- estimating youth unemployment for the \cle{National Youth Employment Programme (NYEP)}.
    \item Develop critical thinking about \cle{non-sampling errors} and practical mitigation strategies.
    \item Practise writing a concise, evidence-based statistical report suitable for decision-makers at the Ministry of Youth Affairs.
\end{itemize}

\begin{attention}[Before you start]
This case study uses \emph{every} lesson of the chapter. Before attempting the tasks, make sure you can:
\begin{itemize}
    \item distinguish a \cle{population} from a \cle{sample}, and a \cle{parameter} from a \cle{statistic};
    \item name the main sampling methods (simple random, systematic, stratified, cluster, quota) with their advantages and disadvantages;
    \item use the confidence interval formula for a proportion, $\hat{p} \pm z_{\alpha/2}\sqrt{\hat{p}(1-\hat{p})/n}$;
    \item invert the margin of error formula to obtain a required sample size.
\end{itemize}
\end{attention}

\courssub{Situation}
The Ministry of Youth Affairs and Civic Education (MINJEC) is preparing the next phase of the \cle{National Youth Employment Programme (NYEP)}. To allocate resources efficiently, the Ministry needs a reliable estimate of the \textbf{proportion of unemployed youths (aged 15--35)} in the \cle{Centre Region}. According to the latest available figures from the National Institute of Statistics (INS), the Centre Region contains approximately \textbf{1\,200\,000} youths in this age bracket.

The Minister has requested a survey with the following specifications:
\begin{itemize}
    \item \textbf{Confidence level:} 95\% (so $z_{0.025}=1.96$).
    \item \textbf{Margin of error:} $\pm 3\%$ (i.e. $E = 0.03$).
    \item \textbf{No prior estimate of the proportion} is available; hence the conservative value $\hat{p}=0.5$ must be used for the sample size calculation (this maximises $\hat{p}(1-\hat{p})$ and therefore gives the safest, largest sample size).
\end{itemize}

A sampling frame exists: the INS master sample of enumeration areas (EAs), covering all divisions of the Centre Region. Each EA lists households with age-eligible members.

After the sample size is determined, a pilot data collection is carried out. The field teams obtain responses from \textbf{1068} youths (the realised sample size after non-response adjustments). Among them, \textbf{542} report being currently unemployed and actively seeking work.

Your group is tasked with designing the survey, performing the calculations, and reporting the findings.

\courssub{Tasks}
\begin{enumerate}
    \item \textbf{Sampling design:} Propose a suitable \cle{sampling frame} and \cle{sampling method} for this survey. Justify your choice with respect to representativeness, cost, and feasibility in the Cameroonian context.
    \item \textbf{Sample size determination:} Compute the minimum required sample size $n$ to achieve a 95\% confidence interval with a margin of error of 3\%, using $\hat{p}=0.5$. Apply the finite population correction (FPC) since the population size $N=1\,200\,000$ is known. Show all steps.
    \item \textbf{Point estimation and confidence intervals:} Using the pilot data ($n=1068$, $x=542$), calculate:
        \begin{itemize}
            \item the point estimate $\hat{p}$;
            \item the 95\% Wald (normal approximation) confidence interval;
            \item the 95\% Wilson score confidence interval.
        \end{itemize}
        Compare the two intervals and comment on their suitability.
    \item \textbf{Non-sampling errors:} Identify at least three potential sources of non-sampling error in this survey (e.g. non-response, measurement error, coverage error). For each, propose a concrete mitigation strategy.
    \item \textbf{Report writing:} Draft a brief report (maximum 500 words) addressed to the Minister of Youth Affairs, summarising the methodology, key results (point estimate and confidence interval), limitations, and clear recommendations for the NYEP.
\end{enumerate}

\begin{methode}[Strategy for solving an estimation problem]
\begin{enumerate}
    \item \textbf{Identify the target parameter.} Here it is a \emph{proportion} $p$ (unemployment rate), not a mean --- this decides which formulae apply.
    \item \textbf{Choose the design before computing.} The sampling method determines whether the simple random sampling formulae are (approximately) valid.
    \item \textbf{Compute the sample size} from $n_0 = z_{\alpha/2}^2\,\hat{p}(1-\hat{p})/E^2$, then apply the FPC if $N$ is known and $n_0/N$ is not negligible.
    \item \textbf{Estimate from the data:} compute $\hat{p}=x/n$, then the standard error, then the interval.
    \item \textbf{Interpret in context:} always conclude with a sentence of the form ``we are 95\% confident that the true proportion lies between \dots and \dots''.
\end{enumerate}
\end{methode}

\courssub{Model answer}

\textbf{Task 1: Sampling design}
\begin{itemize}
    \item \textbf{Sampling frame:} the INS master sample of enumeration areas (EAs), the most up-to-date and geographically exhaustive frame available. It covers urban, peri-urban and rural areas across all divisions of the Centre Region (Mfoundi, Mefou-et-Afamba, Mefou-et-Akono, Mbam-et-Inoubou, Mbam-et-Kim, Nyong-et-K\'ell\'e, Nyong-et-Mfoumou, Nyong-et-So'o, Leki\'e, Nfoundi-area divisions).
    \item \textbf{Sampling method:} \cle{stratified multi-stage cluster sampling}.
        \begin{enumerate}
            \item \emph{Stratification:} divide the region into urban (Yaound\'e) and rural strata to guarantee proportional representation and reduce the variance of the estimator.
            \item \emph{First stage:} randomly select a fixed number of EAs within each stratum, with probability proportional to size (PPS).
            \item \emph{Second stage:} within each selected EA, systematically sample a fixed number of households (e.g. 20) from the household listing.
            \item \emph{Third stage:} in each sampled household, list all eligible youths (15--35) and randomly select one (e.g. using a Kish grid) to avoid intra-household selection bias.
        \end{enumerate}
    \item \textbf{Justification:} this design respects the geographical spread of the population, controls costs by clustering interviews within EAs, allows stratum-specific estimates, and is standard practice for INS household surveys. Simple random sampling over the whole region would be prohibitively expensive because of travel distances; systematic sampling alone risks periodicity bias; quota sampling would not allow valid probability-based confidence intervals.
\end{itemize}

\textbf{Task 2: Sample size determination}

The basic formula for estimating a proportion, using the conservative value $\hat{p}=0.5$ (which maximises $\hat{p}(1-\hat{p})=0.25$), is
\[
n_0 = \frac{z_{\alpha/2}^2 \cdot \hat{p}(1-\hat{p})}{E^2}.
\]
With $z_{0.025}=1.96$, $\hat{p}=0.5$ and $E=0.03$:
\[
n_0 = \frac{1.96^2 \times 0.5 \times 0.5}{0.03^2}
    = \frac{3.8416 \times 0.25}{0.0009}
    = \frac{0.9604}{0.0009}
    = 1067.11\ldots
\]
We always round \emph{up} (rounding down would give a margin of error slightly larger than requested), so $n_0 = 1068$.

Since the population size $N = 1\,200\,000$ is finite, apply the \cle{finite population correction (FPC)}:
\[
n = \frac{n_0}{1 + \dfrac{n_0-1}{N}}
  = \frac{1068}{1 + \dfrac{1067}{1\,200\,000}}
  = \frac{1068}{1.000889}
  \approx 1067.05.
\]
Rounding up again, the required sample size is \textbf{1068 youths}.

\begin{attention}[The FPC is often negligible]
The FPC changes the sample size only when the sample is a substantial fraction of the population (a common rule of thumb: apply it when $n_0/N > 5\%$). Here $n_0/N \approx 0.09\%$, so the correction has virtually no effect. In the GCE examination, state the correction, compute it, and comment --- do not simply ignore it.
\end{attention}

\textbf{Task 3: Point estimation and confidence intervals}

Given $n=1068$ and $x=542$, the point estimate is
\[
\hat{p} = \frac{x}{n} = \frac{542}{1068} = 0.50749\ldots \approx 0.5075.
\]

\emph{Wald (normal approximation) interval:}
\[
\hat{p} \pm z_{0.025}\sqrt{\frac{\hat{p}(1-\hat{p})}{n}}.
\]
Standard error:
\[
SE = \sqrt{\frac{0.5075 \times 0.4925}{1068}}
    = \sqrt{\frac{0.24994}{1068}}
    = \sqrt{0.0002340}
    \approx 0.01530.
\]
Margin of error:
\[
ME = 1.96 \times 0.01530 \approx 0.0300.
\]
Wald 95\% CI:
\[
(0.5075 - 0.0300,\; 0.5075 + 0.0300) = (0.4775,\; 0.5375).
\]

\emph{Wilson score interval} (more accurate coverage, especially for small samples or proportions near 0 or 1):
\[
\frac{\hat{p} + \dfrac{z^2}{2n} \pm z \sqrt{\dfrac{\hat{p}(1-\hat{p})}{n} + \dfrac{z^2}{4n^2}}}{1 + \dfrac{z^2}{n}}.
\]
Compute the components step by step, with $z^2 = 1.96^2 = 3.8416$:
\[
\frac{z^2}{2n} = \frac{3.8416}{2136} = 0.001798, \qquad
\frac{z^2}{n} = \frac{3.8416}{1068} = 0.003597, \qquad
\frac{z^2}{4n^2} = \frac{3.8416}{4 \times 1\,140\,624} \approx 8.42 \times 10^{-7}.
\]
Adjusted centre:
\[
\tilde{p} = \hat{p} + \frac{z^2}{2n} = 0.5075 + 0.0018 = 0.5093.
\]
Square-root term:
\[
\sqrt{\frac{\hat{p}(1-\hat{p})}{n} + \frac{z^2}{4n^2}}
= \sqrt{0.0002340 + 0.0000008}
= \sqrt{0.0002348}
\approx 0.01532.
\]
Margin in the numerator:
\[
z \times 0.01532 = 1.96 \times 0.01532 \approx 0.03003.
\]
Denominator:
\[
1 + \frac{z^2}{n} = 1 + 0.003597 = 1.003597.
\]
Lower and upper limits:
\[
L = \frac{0.5093 - 0.03003}{1.003597} = \frac{0.47927}{1.003597} \approx 0.4776,
\qquad
U = \frac{0.5093 + 0.03003}{1.003597} = \frac{0.53933}{1.003597} \approx 0.5374.
\]
Wilson 95\% CI: \textbf{(0.4776,\; 0.5374)}.

\begin{center}
\begin{tabularx}{\linewidth}{|l|X|X|X|}
\hline
\rowcolor{popblueL} \textbf{Method} & \textbf{Point estimate} & \textbf{95\% CI} & \textbf{Width} \\
\hline
Wald (normal approx.) & 0.5075 & $(0.4775,\;0.5375)$ & 0.0600 \\
\hline
Wilson score & 0.5075 & $(0.4776,\;0.5374)$ & 0.0598 \\
\hline
\end{tabularx}
\end{center}

\emph{Comparison:} the two intervals are nearly identical because $\hat{p}$ is close to $0.5$ and $n$ is large --- exactly the conditions under which the normal approximation works best. The Wilson interval is slightly narrower and is not symmetric about $\hat{p}$; it has better coverage properties for small samples or extreme proportions. For this report either is acceptable, but the Wilson interval is statistically preferable.

\textbf{Task 4: Non-sampling errors and mitigation}
\begin{enumerate}
    \item \textbf{Non-response bias:} unemployed youths may be more likely to be at home and to respond, while employed youths are at work during visits. \emph{Mitigation:} make at least three callbacks at different times (including evenings and weekends); keep the questionnaire short; offer a small incentive such as airtime credit.
    \item \textbf{Measurement error (misclassification):} respondents may misunderstand ``unemployed'' (e.g. students not seeking work, or underemployed workers in the informal sector). \emph{Mitigation:} use an explicit, standard definition in the questionnaire (``Did you actively look for work in the last four weeks?''); train interviewers with role-plays; pilot-test the questions.
    \item \textbf{Coverage error:} the EA frame may miss homeless youths, those in informal settlements, or recent migrants. \emph{Mitigation:} supplement the frame with a targeted module for out-of-household youths; collaborate with local councils and youth associations to locate hidden populations.
    \item \textbf{Processing error:} data entry mistakes or coding errors. \emph{Mitigation:} use Computer-Assisted Personal Interviewing (CAPI) with built-in consistency checks; apply double-entry verification for any paper questionnaires.
\end{enumerate}

\begin{attention}[Sampling error vs non-sampling error]
Do not confuse the two in the examination. \cle{Sampling error} decreases when $n$ increases and is quantified by the standard error. \cle{Non-sampling errors} (non-response, measurement, coverage, processing) do \emph{not} disappear with a larger sample --- a large biased sample is still biased. They are controlled by good design and fieldwork discipline, not by size.
\end{attention}

\textbf{Task 5: Draft report (about 450 words)}

\begin{quote}
\textbf{REPORT TO THE MINISTER OF YOUTH AFFAIRS AND CIVIC EDUCATION}\\
\textbf{Subject:} Estimate of Youth Unemployment in the Centre Region --- NYEP Baseline Survey\\
\textbf{Prepared by:} Statistics Unit, MINJEC / INS Collaboration Team
\end{quote}

\textbf{1. Objective}\\
To provide a statistically sound estimate of the proportion of unemployed youths (15--35 years) in the Centre Region, to inform the next phase of the National Youth Employment Programme (NYEP).

\textbf{2. Methodology}\\
A stratified multi-stage cluster sample was drawn from the INS master frame of enumeration areas. Strata: urban (Yaound\'e) and rural divisions. Within the selected EAs, 20 households each were systematically sampled, and one eligible youth per household was randomly selected. The target sample size of 1068 was calculated for a 95\% confidence level and a $\pm 3\%$ margin of error, using the conservative proportion $0.5$; the finite population correction was applied and found negligible. Fieldwork produced 1068 completed interviews after up to three callbacks per household.

\textbf{3. Key results}
\begin{itemize}
    \item \textbf{Point estimate:} $\hat{p} = 50.7\%$ (542 of 1068 youths unemployed and actively seeking work).
    \item \textbf{95\% confidence interval (Wilson score):} \textbf{47.8\% to 53.7\%}.
    \item The Wald interval (47.8\%--53.8\%) is virtually identical, as expected for a large sample with a central proportion.
\end{itemize}
We are therefore 95\% confident that the true youth unemployment rate in the Centre Region lies between \textbf{47.8\% and 53.7\%}.

\textbf{4. Limitations}
\begin{itemize}
    \item Non-response was slightly higher in urban areas, which could bias the estimate if non-respondents differ systematically from respondents.
    \item The sampling frame excludes homeless and institutionalised youths.
    \item A cross-sectional survey cannot capture seasonal variations in employment.
\end{itemize}

\textbf{5. Recommendations}
\begin{enumerate}
    \item Target NYEP resources to the Centre Region, planning for an unemployment band of roughly 48\%--54\%.
    \item Commission a follow-up survey to estimate unemployment by division, gender and education level for finer targeting.
    \item Improve frame coverage by updating EA maps in fast-growing peri-urban zones with the municipal councils.
    \item Repeat the survey annually with the same methodology to monitor the impact of the NYEP.
\end{enumerate}

\textbf{6. Conclusion}\\
Approximately one in two young people in the Centre Region is unemployed and seeking work. The NYEP should scale up its interventions urgently, using this estimate as a rigorous baseline for future impact evaluation.

\emph{End of report.}

\begin{aretenir}[What this case study teaches]
\begin{itemize}
    \item With no prior estimate, use $\hat{p}=0.5$ in $n_0 = z_{\alpha/2}^2\hat{p}(1-\hat{p})/E^2$; always round the sample size \emph{up}.
    \item The FPC, $n = n_0/\bigl(1+(n_0-1)/N\bigr)$, matters only when the sample is a non-negligible fraction of the population.
    \item For large $n$ and $\hat{p}\approx 0.5$, Wald and Wilson intervals coincide; Wilson is safer in general.
    \item A confidence interval quantifies \emph{sampling} error only; \emph{non-sampling} errors require design and fieldwork solutions.
    \item A statistical report always ends with an interpretation in context and actionable recommendations.
\end{itemize}
\end{aretenir}

\newpage

\coursec{Methods and worked examples}

\courssub{Methods and Worked Examples}

This section turns the theory of sampling and estimation into examination skills. Each method card gives you the reflexes to adopt; the worked example that follows shows, step by step, how a GCE Advanced Level answer should be presented.

\begin{methode}[Choosing and describing a sampling method]
To answer a question asking you to \emph{select} or \emph{describe} a sample:
\begin{enumerate}
\item \textbf{Identify the population} and state whether a \cle{sampling frame} (a complete list of members) exists.
\item \textbf{Decide the sample size} $n$ required.
\item \textbf{Choose the method} according to the situation:
\begin{itemize}
\item \cle{Simple random sampling}: every member has an equal chance; use random number tables or a calculator's random function on a numbered frame.
\item \cle{Systematic sampling}: compute the sampling interval $k=\dfrac{N}{n}$, pick a random start $r$ between $1$ and $k$, then select $r, r+k, r+2k, \dots$
\item \cle{Stratified sampling}: split the population into strata, then take from each stratum a number proportional to its size: $n_h = n\times\dfrac{N_h}{N}$.
\item \cle{Quota / convenience sampling}: non-random; quick but biased — mention this if asked for a disadvantage.
\end{itemize}
\item \textbf{Justify} your choice in one sentence (representativeness, cost, existence of a frame).
\end{enumerate}
\textbf{Reflex:} if the population has identifiable subgroups that differ (streams, sexes, classes), think \emph{stratified}; if there is only a list, think \emph{random} or \emph{systematic}.
\end{methode}

\begin{exoresolu}[Selecting a stratified sample]
A secondary school in Bamenda has $600$ students: $240$ in Lower Sixth Science, $160$ in Lower Sixth Arts, $120$ in Upper Sixth Science and $80$ in Upper Sixth Arts. The principal wants a sample of $60$ students to survey study habits.
\begin{enumerate}
\item Explain why a stratified sample is preferable to a simple random sample here.
\item Calculate the number of students to be selected from each stratum.
\item Describe how the students within one stratum should be chosen.
\end{enumerate}
\tcblower
\textbf{Solution.}
\begin{enumerate}
\item The four classes form natural subgroups whose study habits may differ (examination classes versus non-examination classes, Science versus Arts). A simple random sample could, by chance, over-represent one class and under-represent another. \textbf{Stratified sampling guarantees that each class is represented in proportion to its size}, so the sample reflects the structure of the school.
\item The sampling fraction is $\dfrac{60}{600}=\dfrac{1}{10}$. Applying $n_h = n\times\dfrac{N_h}{N}$:
\begin{align*}
\text{L6 Science: } & 60\times\frac{240}{600}=24, &
\text{L6 Arts: } & 60\times\frac{160}{600}=16,\\
\text{U6 Science: } & 60\times\frac{120}{600}=12, &
\text{U6 Arts: } & 60\times\frac{80}{600}=8.
\end{align*}
Check: $24+16+12+8=60$. \checkmark
\item Within each stratum, use \textbf{simple random sampling}: number the students of the class from the class list (the sampling frame), then generate random numbers (calculator or random number table), ignoring repeats and numbers outside the range, until the required number of students is obtained.
\end{enumerate}
\end{exoresolu}

\begin{methode}[Distinguishing parameter, statistic and sampling error]
\begin{enumerate}
\item A \cle{parameter} is a numerical characteristic of the \textbf{whole population} (e.g. $\mu$, $\sigma^2$, $p$). It is fixed but usually unknown.
\item A \cle{statistic} is a numerical characteristic computed from the \textbf{sample} (e.g. $\bar{x}$, $s^2$, $\hat{p}$). It varies from sample to sample.
\item The \cle{sampling error} is the difference $|\text{statistic}-\text{parameter}|$ due to chance; it decreases roughly like $\dfrac{1}{\sqrt{n}}$.
\item \textbf{Reflex:} in an exam, always write one sentence stating which symbol is the parameter and which is the statistic before computing anything.
\end{enumerate}
\end{methode}

\begin{exoresolu}[Parameter versus statistic]
A farmer in the NWR wants to know the mean mass of all the $10\,000$ yams in his store. He weighs a random sample of $50$ yams and obtains a mean mass of $2.35\ \mathrm{kg}$. The true (unknown) mean mass of all the yams is $2.40\ \mathrm{kg}$.
\begin{enumerate}
\item Identify the population, the sample, the parameter and the statistic.
\item Calculate the sampling error and state one way of reducing it.
\end{enumerate}
\tcblower
\textbf{Solution.}
\begin{enumerate}
\item \textbf{Population:} all $10\,000$ yams in the store. \textbf{Sample:} the $50$ yams weighed. \textbf{Parameter:} the true mean mass $\mu = 2.40\ \mathrm{kg}$ (fixed, characteristic of the whole population). \textbf{Statistic:} the sample mean $\bar{x}=2.35\ \mathrm{kg}$ (computed from the sample; another sample of $50$ yams would give a different value).
\item The sampling error is
\[
|\bar{x}-\mu| = |2.35-2.40| = 0.05\ \mathrm{kg}.
\]
To reduce it, \textbf{increase the sample size} $n$: since the standard error of the mean is $\dfrac{\sigma}{\sqrt{n}}$, quadrupling $n$ (from $50$ to $200$) would halve the typical sampling error. Using a properly random selection also removes systematic (non-sampling) bias.
\end{enumerate}
\end{exoresolu}

\begin{methode}[Computing point estimates of $\mu$ and $\sigma^2$ from raw data]
Given a sample $x_1,\dots,x_n$:
\begin{enumerate}
\item The \cle{point estimate} of the population mean is the sample mean
\[
\bar{x}=\frac{\sum x}{n}.
\]
\item The \textbf{unbiased} estimate of the population variance uses $n-1$:
\[
s^2=\frac{1}{n-1}\left(\sum x^2-\frac{(\sum x)^2}{n}\right)=\frac{\sum(x-\bar x)^2}{n-1}.
\]
\item The estimate of the standard deviation is $s=\sqrt{s^2}$.
\item \textbf{Reflex:} compute $\sum x$ and $\sum x^2$ first, in a small table; then apply the computing formula — it is far faster and less error-prone than summing deviations.
\end{enumerate}
\textbf{Warning:} dividing by $n$ instead of $n-1$ gives a \emph{biased} estimate of $\sigma^2$; the GCE expects the unbiased version unless stated otherwise.
\end{methode}

\begin{exoresolu}[Point estimates from a small sample]
The masses (in kg) of a random sample of $8$ bags of rice from a consignment at the Douala port are:
\[
49.2,\ 50.1,\ 49.8,\ 50.4,\ 49.5,\ 50.0,\ 49.7,\ 50.3.
\]
Calculate unbiased point estimates of the mean and of the variance of the masses of all bags in the consignment.
\tcblower
\textbf{Solution.} Build the sums:
\[
\sum x = 49.2+50.1+49.8+50.4+49.5+50.0+49.7+50.3 = 399.0,
\]
\[
\sum x^2 = 49.2^2+50.1^2+49.8^2+50.4^2+49.5^2+50.0^2+49.7^2+50.3^2 = 19\,901.28.
\]
\textbf{Estimate of the mean:}
\[
\bar{x}=\frac{\sum x}{n}=\frac{399.0}{8}=49.875\ \mathrm{kg}.
\]
\textbf{Estimate of the variance} (unbiased, divisor $n-1=7$):
\begin{align*}
s^2 &= \frac{1}{n-1}\left(\sum x^2 - \frac{(\sum x)^2}{n}\right)
= \frac{1}{7}\left(19\,901.28-\frac{399.0^2}{8}\right)\\
&= \frac{1}{7}\left(19\,901.28-19\,900.125\right)
= \frac{1.155}{7} = 0.165\ \mathrm{kg^2}.
\end{align*}
Hence $s=\sqrt{0.165}\approx 0.41\ \mathrm{kg}$.
\textbf{Conclusion:} $\hat{\mu}=49.875\ \mathrm{kg}$, $\hat{\sigma}^2 = 0.165\ \mathrm{kg}^2$.
\end{exoresolu}

\begin{attention}[A frequent arithmetic trap]
In the computing formula $\sum x^2 - \dfrac{(\sum x)^2}{n}$, the two terms are large and nearly equal, so a small slip in $\sum x^2$ badly distorts $s^2$. \textbf{Always re-add the squares carefully}, and check plausibility: here the data range over about $1.2\ \mathrm{kg}$, so a standard deviation near $0.4\ \mathrm{kg}$ is reasonable, whereas $s^2$ coming out negative or larger than $1$ would signal an error.
\end{attention}

\begin{methode}[Verifying that an estimator is unbiased]
\begin{enumerate}
\item Recall: an estimator $T$ is an \cle{unbiased estimator} of a parameter $\theta$ if $E(T)=\theta$.
\item Use the linearity of expectation: $E(aX+bY)=aE(X)+bE(Y)$, and for a random sample $E(X_i)=\mu$ for every $i$.
\item Compute $E(T)$ symbolically and compare with $\theta$.
\item \textbf{Key results to quote:} $E(\bar{X})=\mu$ (so $\bar X$ is unbiased for $\mu$); $E(S^2)=\sigma^2$ with divisor $n-1$.
\item \textbf{Reflex:} if $E(T)=c\,\theta$ with $c\neq 1$, then $\dfrac{T}{c}$ is an unbiased estimator of $\theta$.
\end{enumerate}
\end{methode}

\begin{exoresolu}[Testing unbiasedness of a proposed estimator]
Let $X_1, X_2, X_3$ be a random sample from a population with mean $\mu$. A student proposes the estimator
\[
T=\frac{X_1+2X_2+X_3}{3}.
\]
\begin{enumerate}
\item Show that $T$ is a biased estimator of $\mu$ and find its bias.
\item Write down an unbiased estimator of $\mu$ based on the same sample.
\end{enumerate}
\tcblower
\textbf{Solution.}
\begin{enumerate}
\item Since the sample is random, $E(X_1)=E(X_2)=E(X_3)=\mu$. By linearity of expectation:
\begin{align*}
E(T) &= \frac{1}{3}\left[E(X_1)+2E(X_2)+E(X_3)\right]
= \frac{1}{3}\left(\mu+2\mu+\mu\right)
= \frac{4\mu}{3}.
\end{align*}
Since $E(T)=\dfrac{4\mu}{3}\neq\mu$, $T$ is \textbf{biased}. The bias is
\[
E(T)-\mu=\frac{4\mu}{3}-\mu=\frac{\mu}{3}.
\]
\item Two natural answers:
\begin{itemize}
\item Correcting $T$: $T'=\dfrac{3T}{4}=\dfrac{X_1+2X_2+X_3}{4}$, since $E(T')=\dfrac{3}{4}\cdot\dfrac{4\mu}{3}=\mu$.
\item The sample mean: $\bar{X}=\dfrac{X_1+X_2+X_3}{3}$, with $E(\bar{X})=\dfrac{3\mu}{3}=\mu$.
\end{itemize}
Both are unbiased; $\bar X$ is preferred because it also has the \textbf{smallest variance} among linear unbiased estimators (it is efficient).
\end{enumerate}
\end{exoresolu}

\begin{methode}[Confidence interval for $\mu$ when $\sigma$ is known]
\begin{enumerate}
\item \textbf{Check the conditions:} $\sigma$ known; either the population is normal, or $n$ is large ($n\ge 30$, Central Limit Theorem).
\item Write the formula:
\[
\bar{x}\pm z_{\alpha/2}\,\frac{\sigma}{\sqrt{n}}
\quad\text{i.e.}\quad
\left[\bar{x}-z_{\alpha/2}\frac{\sigma}{\sqrt{n}}\;;\;\bar{x}+z_{\alpha/2}\frac{\sigma}{\sqrt{n}}\right].
\]
\item Read the critical value: $z_{0.025}=1.96$ (95\%), $z_{0.05}=1.645$ (90\%), $z_{0.005}=2.576$ (99\%).
\item Compute the standard error $\dfrac{\sigma}{\sqrt{n}}$, then the margin of error, then the bounds.
\item \textbf{Interpret in context:} "We are $95\%$ confident that the true mean ... lies between ... and ...".
\end{enumerate}
\end{methode}

\begin{exoresolu}[CI for a mean, known variance — GCE style]
A machine fills packets of sugar. The masses are normally distributed with known standard deviation $\sigma = 4\ \mathrm{g}$. A random sample of $25$ packets gives a mean mass of $\bar{x}=502.3\ \mathrm{g}$.
\begin{enumerate}
\item Construct a $95\%$ confidence interval for the true mean mass $\mu$.
\item The label claims $500\ \mathrm{g}$. Comment.
\end{enumerate}
\tcblower
\textbf{Solution.}
\begin{enumerate}
\item Population normal, $\sigma$ known: use $\bar{x}\pm z_{0.025}\dfrac{\sigma}{\sqrt{n}}$ with $z_{0.025}=1.96$.
\[
\text{Standard error: } \frac{\sigma}{\sqrt{n}}=\frac{4}{\sqrt{25}}=\frac{4}{5}=0.8\ \mathrm{g}.
\]
\[
\text{Margin of error: } 1.96\times 0.8 = 1.568\ \mathrm{g}.
\]
\[
\text{CI: } 502.3\pm 1.568 = [\,500.732\ ;\ 503.868\,]\ \mathrm{g}.
\]
We are $95\%$ confident that the true mean mass of all packets lies between approximately $500.73\ \mathrm{g}$ and $503.87\ \mathrm{g}$.
\item The claimed value $500\ \mathrm{g}$ lies \textbf{below} the lower bound $500.73\ \mathrm{g}$ of the interval. The data therefore suggest that the machine is \emph{overfilling}: the true mean is very likely greater than $500\ \mathrm{g}$. The producer may wish to adjust the machine to avoid giving away sugar (while still respecting the legal minimum).
\end{enumerate}
\end{exoresolu}

\begin{popfigure}
\begin{center}
\begin{tikzpicture}[scale=1]
\draw[->, popmuted] (-0.3,0) -- (10.3,0) node[below left] {\small $\mu$};
\draw[popdark, thick] (5,0.1) -- (5,-0.1) node[below] {\small $\mu$};
\foreach \y/\a/\b/\ok in {0.4/3.6/6.2/1, 0.8/4.1/6.9/1, 1.2/3.2/5.8/1, 1.6/4.4/7.0/1, 2.0/5.4/8.0/0, 2.4/3.8/6.4/1, 2.8/4.0/6.6/1, 3.2/3.4/6.0/1, 3.6/4.2/6.8/1, 4.0/3.7/6.3/1, 4.4/4.5/7.1/1, 4.8/1.9/4.5/0, 5.2/3.9/6.5/1, 5.6/4.3/6.9/1, 6.0/3.5/6.1/1, 6.4/4.1/6.7/1, 6.8/3.3/5.9/1, 7.2/4.6/7.2/1, 7.6/3.8/6.4/1, 8.0/4.0/6.6/1}{
\ifnum\ok=1
\draw[popblue, thick] (\a,\y) -- (\b,\y);
\fill[popblue] ({0.5*(\a+\b)},\y) circle (0.05);
\else
\draw[poporange, thick] (\a,\y) -- (\b,\y);
\fill[poporange] ({0.5*(\a+\b)},\y) circle (0.05);
\fi
}
\draw[popdark, dashed] (5,0) -- (5,8.3);
\node[popdark, align=center] at (5,8.7) {\small true mean};
\node[popblue, align=center] at (8.6,7.6) {\small intervals\\ \small containing $\mu$};
\node[poporange, align=center] at (8.6,5.0) {\small intervals\\ \small missing $\mu$};
\end{tikzpicture}
\end{center}
\legende{Twenty 95\% confidence intervals from twenty different samples of the same population. In the long run about 95\% of such intervals capture the fixed (unknown) mean $\mu$; here 18 out of 20 do.}
\end{popfigure}

\begin{methode}[CI for $\mu$ with unknown variance, and CI for a proportion]
\textbf{Mean, $\sigma$ unknown, small sample ($n<30$, population assumed normal):}
\begin{enumerate}
\item Estimate $\sigma$ by $s$ (divisor $n-1$).
\item Use Student's $t$-distribution with $\nu=n-1$ degrees of freedom:
\[
\bar{x}\pm t_{n-1,\,\alpha/2}\,\frac{s}{\sqrt{n}}.
\]
Read $t$ from the table (e.g. $t_{9,\,0.025}=2.262$).
\end{enumerate}
\textbf{Proportion (large sample):}
\begin{enumerate}
\item Compute $\hat{p}=\dfrac{x}{n}$.
\item Use
\[
\hat{p}\pm z_{\alpha/2}\sqrt{\frac{\hat{p}(1-\hat{p})}{n}}.
\]
\item \textbf{Reflex:} for a proportion, never forget the $\hat p(1-\hat p)$ under the square root; for a mean, the standard error is $\dfrac{\sigma\text{ or }s}{\sqrt n}$.
\end{enumerate}
\end{methode}

\begin{exoresolu}[Two intervals in one problem]
\textbf{(a)} The heights (cm) of a random sample of $10$ students give $\bar{x}=168.2$ and $s=6.4$. Heights are assumed normal. Find a $95\%$ confidence interval for the mean height of all students. (Use $t_{9,\,0.025}=2.262$.)

\textbf{(b)} In a random sample of $200$ voters in a municipality, $86$ say they will vote for candidate A. Find a $99\%$ confidence interval for the proportion $p$ of all voters supporting A. (Use $z_{0.005}=2.576$.)
\tcblower
\textbf{Solution.}

\textbf{(a)} $\sigma$ unknown, $n=10<30$, population normal: use the $t$-interval with $\nu=9$:
\[
\bar{x}\pm t_{9,\,0.025}\frac{s}{\sqrt{n}}
=168.2\pm 2.262\times\frac{6.4}{\sqrt{10}}
=168.2\pm 2.262\times 2.0239
=168.2\pm 4.578.
\]
\[
\text{CI} = [\,163.62\ ;\ 172.78\,]\ \mathrm{cm}.
\]
We are $95\%$ confident that the mean height of all students lies between about $163.6$ cm and $172.8$ cm.

\textbf{(b)} $\hat{p}=\dfrac{86}{200}=0.43$, $n=200$ (large). Then
\[
\sqrt{\frac{\hat{p}(1-\hat{p})}{n}}=\sqrt{\frac{0.43\times0.57}{200}}=\sqrt{0.0012255}=0.03501.
\]
\[
\text{Margin of error: } 2.576\times0.03501=0.0902.
\]
\[
\text{CI} = 0.43\pm0.0902 = [\,0.3398\ ;\ 0.5202\,].
\]
We are $99\%$ confident that between about $34.0\%$ and $52.0\%$ of all voters support candidate A. Since the interval contains values below $50\%$, the sample does not allow us to conclude, at this confidence level, that A has majority support.
\end{exoresolu}

\begin{methode}[Determining the required sample size]
\begin{enumerate}
\item \textbf{For a mean} (target margin of error $E$, confidence level giving $z$):
\[
E=z_{\alpha/2}\frac{\sigma}{\sqrt{n}}
\quad\Longrightarrow\quad
n=\left(\frac{z_{\alpha/2}\,\sigma}{E}\right)^{2}.
\]
\item \textbf{For a proportion:}
\[
n=\frac{z_{\alpha/2}^{2}\,\hat{p}(1-\hat{p})}{E^{2}};
\qquad\text{if no estimate of }p\text{ exists, take } \hat p = 0.5
\]
(worst case, since $p(1-p)$ is maximal at $p=0.5$).
\item \textbf{Always round UP} to the next integer (rounding down would make the error larger than allowed).
\item \textbf{Reflex:} halving the margin of error multiplies $n$ by $4$; raising confidence from $95\%$ to $99\%$ also increases $n$.
\end{enumerate}
\end{methode}

\begin{exoresolu}[Sample size for a pre-election poll]
A polling organisation in Yaoundé wants to estimate the proportion $p$ of voters favouring a candidate, with a margin of error of at most $2\%$ at the $95\%$ confidence level.
\begin{enumerate}
\item A pilot survey suggests $p\approx 0.55$. What sample size is needed?
\item What size is needed if no prior estimate of $p$ is available?
\item State one advantage and one disadvantage of increasing the confidence level to $99\%$.
\end{enumerate}
\tcblower
\textbf{Solution.}
\begin{enumerate}
\item With $z_{0.025}=1.96$, $E=0.02$, $\hat p = 0.55$:
\[
n=\frac{z^{2}\hat p(1-\hat p)}{E^{2}}
=\frac{1.96^{2}\times0.55\times0.45}{0.02^{2}}
=\frac{3.8416\times0.2475}{0.0004}
=\frac{0.950796}{0.0004}=2376.99.
\]
Round up: $\boxed{n=2377}$ voters.
\item With no prior estimate, use the worst case $\hat p = 0.5$:
\[
n=\frac{1.96^{2}\times0.25}{0.0004}=\frac{0.9604}{0.0004}=2401.
\]
So $\boxed{n=2401}$ voters guarantee the required precision whatever the true $p$.
\item \textbf{Advantage:} greater confidence that the interval captures the true proportion. \textbf{Disadvantage:} $z$ rises from $1.96$ to $2.576$, so $n$ is multiplied by $\left(\dfrac{2.576}{1.96}\right)^2\approx 1.73$ — the survey becomes considerably more expensive and slower to conduct.
\end{enumerate}
\end{exoresolu}

\begin{methode}[Avoiding the classic mistakes in estimation]
\begin{enumerate}
\item Never use $z$ when $\sigma$ is unknown and $n$ is small — use $t_{n-1}$.
\item Never divide by $n$ when estimating $\sigma^2$ from a sample — divide by $n-1$.
\item Never interpret a $95\%$ CI as "a $95\%$ probability that $\mu$ is in this interval" after it is computed: $\mu$ is fixed; the $95\%$ refers to the \emph{method} capturing $\mu$ in $95\%$ of repeated samples.
\item Never forget to round the sample size \emph{up}.
\item Never apply the proportion formula to a mean, or vice versa.
\end{enumerate}
\end{methode}

\begin{exoresolu}[Spotting and correcting errors — synoptic exercise]
A student analyses the lifetime (hours) of $n=16$ phone batteries, assumed normally distributed, and obtains $\bar{x}=41.2$, $s=3.6$. He writes: "A $95\%$ CI for the mean lifetime is $41.2\pm1.96\times\frac{3.6}{\sqrt{16}}=[39.44\,;\,42.96]$, so there is a $95\%$ probability that the true mean lies in this interval." (Given $t_{15,\,0.025}=2.131$.)
\begin{enumerate}
\item Identify and correct the two errors.
\item Give the correct interval and a correct interpretation.
\end{enumerate}
\tcblower
\textbf{Solution.}
\begin{enumerate}
\item \textbf{Error 1 (wrong distribution):} $\sigma$ is unknown and $n=16<30$ with a normal population, so Student's $t$ with $\nu=15$ must be used, not $z=1.96$. \textbf{Error 2 (wrong interpretation):} once the interval is computed, the true mean $\mu$ either is or is not inside it; the $95\%$ describes the long-run success rate of the method, not a probability about this particular interval.
\item Correct computation:
\[
\bar{x}\pm t_{15,\,0.025}\frac{s}{\sqrt{n}}
=41.2\pm2.131\times\frac{3.6}{4}
=41.2\pm2.131\times0.9
=41.2\pm1.918.
\]
\[
\text{CI}=[\,39.28\ ;\ 43.12\,]\ \text{hours}.
\]
\textbf{Correct interpretation:} "We are $95\%$ confident that the true mean lifetime of all such batteries lies between about $39.3$ and $43.1$ hours" — meaning that if we repeated the sampling many times, about $95\%$ of the intervals built this way would contain $\mu$. Note the correct interval is \emph{wider} than the student's: using $t$ instead of $z$ properly reflects the extra uncertainty from estimating $\sigma$.
\end{enumerate}
\end{exoresolu}

\begin{aretenir}[Formula summary for the examination]
\begin{center}
\begin{tabularx}{\linewidth}{|l|X|X|}
\hline
\rowcolor{popblueL} \textbf{Quantity estimated} & \textbf{Point estimate} & \textbf{Confidence interval} \\
\hline
Mean $\mu$, $\sigma$ known & $\bar{x}=\dfrac{\sum x}{n}$ & $\bar{x}\pm z_{\alpha/2}\dfrac{\sigma}{\sqrt{n}}$ \\
\hline
Mean $\mu$, $\sigma$ unknown, $n<30$ (normal pop.) & $\bar{x}$,\quad $s^2=\dfrac{\sum(x-\bar x)^2}{n-1}$ & $\bar{x}\pm t_{n-1,\,\alpha/2}\dfrac{s}{\sqrt{n}}$ \\
\hline
Proportion $p$, $n$ large & $\hat{p}=\dfrac{x}{n}$ & $\hat{p}\pm z_{\alpha/2}\sqrt{\dfrac{\hat{p}(1-\hat{p})}{n}}$ \\
\hline
\end{tabularx}
\end{center}
Critical values: $1.645$ (90\%), $1.96$ (95\%), $2.576$ (99\%). Sample sizes: $n=\left(\dfrac{z\sigma}{E}\right)^2$ for a mean; $n=\dfrac{z^2\hat p(1-\hat p)}{E^2}$ for a proportion — always round \textbf{up}.
\end{aretenir}

\newpage

\coursec{Exercises}

\courssub{I check my knowledge}

\begin{exercice}[MCQ: Sampling Terminology]
A researcher selects every 10th name from an alphabetical list of 5000 registered voters in Yaoundé to form a sample of 500. This sampling method is best described as:
\begin{enumerate}
\item Simple random sampling
\item Stratified sampling
\item Systematic sampling
\item Cluster sampling
\end{enumerate}
\end{exercice}

\begin{exercice}[True/False with Justification]
State whether each statement is True or False. Justify your answer briefly.
\begin{enumerate}
\item The sample mean $\bar{X}$ is always equal to the population mean $\mu$.
\item A 95\% confidence interval for a population mean implies there is a 0.95 probability that the true mean lies within the calculated interval.
\item The standard error of the mean decreases as the sample size increases.
\item The $t$-distribution approaches the standard normal distribution as the degrees of freedom increase.
\end{enumerate}
\end{exercice}

\begin{exercice}[Definitions and Concepts]
Define the following terms in the context of statistical inference:
\begin{enumerate}
\item Sampling frame
\item Parameter vs. Statistic
\item Sampling error vs. Non-sampling error
\item Unbiased estimator
\end{enumerate}
\end{exercice}

\begin{exercice}[Direct Calculation: Standard Error]
A population has a standard deviation $\sigma = 12$. Calculate the standard error of the mean for samples of size:
\begin{enumerate}
\item $n = 36$
\item $n = 144$
\item What happens to the standard error when the sample size is quadrupled?
\end{enumerate}
\end{exercice}

\courssub{I practise}

\begin{exercice}[Confidence Interval for Mean (Known Variance)]
A random sample of 50 cocoa pods from a plantation in the South-West Region yields a mean weight of $\SI{320}{g}$ with a known population standard deviation of $\SI{45}{g}$. Construct a 99\% confidence interval for the true mean pod weight. Interpret your result.
\end{exercice}

\begin{exercice}[Confidence Interval for Mean (Unknown Variance)]
The daily milk yields (in litres) of 15 dairy cows in the West Region are recorded as follows:
\[ 12.4,\ 13.1,\ 11.8,\ 14.2,\ 12.9,\ 13.5,\ 11.6,\ 12.7,\ 13.8,\ 12.2,\ 14.0,\ 13.3,\ 12.5,\ 13.0,\ 12.8 \]
Assuming the yields are normally distributed, construct a 95\% confidence interval for the true mean daily milk yield.
\end{exercice}

\begin{exercice}[Confidence Interval for a Proportion]
In a survey of 400 randomly selected voters in Douala, 210 indicate they support Candidate A.
\begin{enumerate}
\item Compute the 95\% Wald (standard) confidence interval for the true proportion of support.
\item Compute the 95\% Wilson score confidence interval for the same proportion.
\item Compare the two intervals and comment on which is more appropriate.
\end{enumerate}
\end{exercice}

\begin{exercice}[Determining Sample Size]
A researcher wants to estimate the average daily milk yield of dairy cows in the West Region with a margin of error of $\SI{0.5}{L}$ at 95\% confidence. Pilot data from a small study give a sample standard deviation $s = \SI{2.1}{L}$.
\begin{enumerate}
\item Determine the minimum sample size required using the pilot standard deviation as an estimate for $\sigma$.
\item If the researcher wants to halve the margin of error to $\SI{0.25}{L}$, by what factor must the sample size increase?
\end{enumerate}
\end{exercice}

\courssub{I prepare for the GCE}

\begin{exercice}[Properties of Estimators and Simulation] \hfill (Total: 12 marks)
\begin{enumerate}
\item Define what it means for an estimator $\hat{\theta}$ to be \textbf{unbiased}, \textbf{consistent}, and \textbf{efficient}. \hfill (3 marks)
\item Let $X_1, X_2, \dots, X_n$ be a random sample from a normal distribution $N(\mu, \sigma^2)$. Show that the sample mean $\bar{X}$ is an unbiased estimator of $\mu$ and find its variance. \hfill (3 marks)
\item The sample median $\tilde{X}$ is also an unbiased estimator of $\mu$ for a normal distribution. Explain why $\bar{X}$ is preferred over $\tilde{X}$ as an estimator of $\mu$. \hfill (2 marks)
\item A student simulates 10\,000 samples of size $n=20$ from a $N(100, 25)$ distribution and computes both the sample mean and sample median for each sample. The standard deviation of the 10\,000 sample means is found to be approximately 1.12, while the standard deviation of the 10\,000 sample medians is approximately 1.40. Use these results to comment on the relative efficiency of the two estimators. \hfill (4 marks)
\end{enumerate}
\end{exercice}

\begin{exercice}[Confidence Intervals: Interpretation and Common Mistakes] \hfill (Total: 13 marks)
A quality control manager at a bottling plant in Douala takes a random sample of 40 bottles. The sample mean fill volume is $\SI{752}{ml}$ with a sample standard deviation of $\SI{8}{ml}$. The machine is set to fill $\SI{750}{ml}$.
\begin{enumerate}
\item Construct a 95\% confidence interval for the true mean fill volume $\mu$. State any assumptions made. \hfill (5 marks)
\item Based on your interval, does the machine appear to be correctly calibrated? Justify your answer. \hfill (2 marks)
\item Identify and describe \textbf{three} common mistakes students make when interpreting a 95\% confidence interval. For each mistake, provide the correct interpretation. \hfill (6 marks)
\end{enumerate}
\end{exercice}

\begin{exercice}[Sample Size Determination and Practical Application] \hfill (Total: 15 marks)
The Ministry of Public Health wishes to estimate the proportion $p$ of households in a certain health district that use insecticide-treated mosquito nets (ITNs).
\begin{enumerate}
\item A pilot survey of 100 households finds 68 using ITNs. Calculate a 95\% confidence interval for $p$. \hfill (3 marks)
\item The Ministry wants a new survey with a margin of error not exceeding 3\% at 95\% confidence. Determine the minimum sample size required:
      \begin{itemize}
      \item[a)] Using the pilot study estimate $\hat{p} = 0.68$.
      \item[b)] Using a conservative approach (no prior estimate).
      \end{itemize}
      \hfill (5 marks)
\item Suppose the actual sample size used is 1200 households, and 810 report using ITNs. Test at the 5\% significance level whether the true proportion of ITN usage differs from 0.68. State your hypotheses, test statistic, critical region, and conclusion in context. \hfill (7 marks)
\end{enumerate}
\end{exercice}

\newpage

\coursec{Solutions to the exercises}

\coursec{Sampling and Estimation}

\courssub{I check my knowledge}

\begin{corrige}[Exercise 1 — MCQ: Sampling Terminology]
The correct answer is \textbf{3. Systematic sampling}.

\textbf{Justification.} In systematic sampling, the population is arranged in a list (here, the alphabetical list of 5000 registered voters), a sampling interval $k$ is computed as
\[ k = \frac{N}{n} = \frac{5000}{500} = 10, \]
a random starting point is chosen among the first $k$ names, and then every $k$-th name (here every 10th) is selected. This matches exactly the procedure described.

The other options do not fit:
\begin{itemize}
\item \emph{Simple random sampling} would require each of the 500 voters to be chosen purely at random, with no fixed interval.
\item \emph{Stratified sampling} would require dividing the voters into homogeneous subgroups (strata), e.g. by quarter or by age group, and sampling within each stratum.
\item \emph{Cluster sampling} would require dividing the population into clusters (e.g. neighbourhoods), randomly selecting whole clusters, and surveying everyone (or a subsample) within the chosen clusters.
\end{itemize}
\end{corrige}

\begin{corrige}[Exercise 2 — True/False with Justification]
\begin{enumerate}
\item \textbf{False.} The sample mean $\bar{X}$ is a random variable whose value depends on the particular sample drawn. It is an \emph{unbiased estimator} of $\mu$, meaning $E(\bar{X}) = \mu$: it equals $\mu$ \emph{on average} over all possible samples, but for any single sample $\bar{X}$ generally differs from $\mu$ by a sampling error.
\item \textbf{False.} Once the interval is computed, it either contains $\mu$ or it does not; $\mu$ is a fixed (unknown) constant, not a random variable. The correct interpretation is: the \emph{procedure} used produces intervals that capture the true mean in $95\%$ of repeated samples. We say we are $95\%$ confident, not that the probability is $0.95$ for this particular interval.
\item \textbf{True.} The standard error is $\mathrm{SE}(\bar{X}) = \dfrac{\sigma}{\sqrt{n}}$. Since $\sqrt{n}$ increases with $n$, the standard error decreases (it is halved when $n$ is multiplied by 4).
\item \textbf{True.} As the number of degrees of freedom $\nu$ increases, the $t$-distribution becomes less heavy-tailed and converges to the standard normal distribution $N(0,1)$. For $\nu \geq 30$, the two distributions are already very close.
\end{enumerate}
\end{corrige}

\begin{corrige}[Exercise 3 — Definitions and Concepts]
\begin{enumerate}
\item \textbf{Sampling frame:} the complete list (or operational representation) of all members of the target population from which the sample is actually drawn, e.g. the electoral register of Yaoundé voters. Ideally it should be exhaustive, up to date and free of duplicates.
\item \textbf{Parameter vs.\ Statistic:} a \emph{parameter} is a fixed numerical characteristic of the whole population (e.g. the population mean $\mu$, the population variance $\sigma^2$, the population proportion $p$); it is usually unknown. A \emph{statistic} is a numerical characteristic computed from a sample (e.g. $\bar{X}$, $S^2$, $\hat{p}$); it varies from sample to sample and is used to estimate the corresponding parameter.
\item \textbf{Sampling error vs.\ Non-sampling error:} \emph{sampling error} is the discrepancy between a sample statistic and the true parameter caused purely by the fact that we observe only a part of the population; it decreases as $n$ increases. \emph{Non-sampling errors} are all other errors: non-response, poorly worded questionnaires, measurement errors, faulty sampling frame, data-entry mistakes; they can occur even in a census and do not necessarily decrease with sample size.
\item \textbf{Unbiased estimator:} an estimator $\hat{\theta}$ of a parameter $\theta$ is unbiased if its expected value equals the parameter, i.e. $E(\hat{\theta}) = \theta$. On average, over repeated sampling, it hits the true value; its bias $E(\hat{\theta}) - \theta$ is zero.
\end{enumerate}
\end{corrige}

\begin{corrige}[Exercise 4 — Direct Calculation: Standard Error]
The standard error of the mean is $\mathrm{SE}(\bar{X}) = \dfrac{\sigma}{\sqrt{n}}$ with $\sigma = 12$.
\begin{enumerate}
\item For $n = 36$: $\mathrm{SE} = \dfrac{12}{\sqrt{36}} = \dfrac{12}{6} = 2$.
\item For $n = 144$: $\mathrm{SE} = \dfrac{12}{\sqrt{144}} = \dfrac{12}{12} = 1$.
\item When the sample size is quadrupled (from 36 to 144), the standard error is \textbf{divided by 2} (from 2 to 1). In general, since $\mathrm{SE} \propto \dfrac{1}{\sqrt{n}}$, multiplying $n$ by 4 divides the standard error by $\sqrt{4} = 2$. This illustrates the diminishing returns of increasing sample size: to halve the error, one must quadruple the sample.
\end{enumerate}
\end{corrige}

\courssub{I practise}

\begin{corrige}[Exercise 5 — Confidence Interval for Mean (Known Variance)]
\textbf{Given:} $n = 50$, $\bar{x} = 320\ \mathrm{g}$, $\sigma = 45\ \mathrm{g}$ (known), confidence level $99\%$.

Since $\sigma$ is known and $n = 50$ is large, we use the normal distribution:
\[ \bar{x} \pm z_{\alpha/2}\,\frac{\sigma}{\sqrt{n}}. \]
For a $99\%$ confidence level, $z_{0.005} = 2.576$.

Standard error:
\[ \frac{\sigma}{\sqrt{n}} = \frac{45}{\sqrt{50}} = \frac{45}{7.0711} \approx 6.364. \]

Margin of error:
\[ E = 2.576 \times 6.364 \approx 16.4\ \mathrm{g}. \]

Confidence interval:
\[ 320 \pm 16.4 = [303.6\,;\,336.4]\ \mathrm{g}. \]

\textbf{Interpretation.} We are $99\%$ confident that the true mean weight of cocoa pods on this plantation lies between approximately $303.6\ \mathrm{g}$ and $336.4\ \mathrm{g}$. Equivalently, if we repeated this sampling procedure many times, about $99\%$ of the intervals constructed would contain the true mean pod weight.
\end{corrige}

\begin{corrige}[Exercise 6 — Confidence Interval for Mean (Unknown Variance)]
\textbf{Step 1: Compute the sample mean.}
\[ \sum x_i = 12.4 + 13.1 + 11.8 + 14.2 + 12.9 + 13.5 + 11.6 + 12.7 + 13.8 + 12.2 + 14.0 + 13.3 + 12.5 + 13.0 + 12.8 = 191.8. \]
\[ \bar{x} = \frac{191.8}{15} \approx 12.787\ \mathrm{L}. \]

\textbf{Step 2: Compute the sample standard deviation} (divisor $n-1 = 14$).
\begin{center}
\begin{tabular}{|c|c|c|}
\hline
\rowcolor{popblueL} $x_i$ & $x_i - \bar{x}$ & $(x_i - \bar{x})^2$ \\
\hline
12.4 & $-0.387$ & 0.1495 \\
13.1 & 0.313 & 0.0982 \\
11.8 & $-0.987$ & 0.9735 \\
14.2 & 1.413 & 1.9975 \\
12.9 & 0.113 & 0.0128 \\
13.5 & 0.713 & 0.5088 \\
11.6 & $-1.187$ & 1.4082 \\
12.7 & $-0.087$ & 0.0075 \\
13.8 & 1.013 & 1.0268 \\
12.2 & $-0.587$ & 0.3442 \\
14.0 & 1.213 & 1.4722 \\
13.3 & 0.513 & 0.2635 \\
12.5 & $-0.287$ & 0.0822 \\
13.0 & 0.213 & 0.0455 \\
12.8 & 0.013 & 0.0002 \\
\hline
\textbf{Sum} & & $\approx 8.3907$ \\
\hline
\end{tabular}
\end{center}
\[ s^2 = \frac{8.3907}{14} \approx 0.5993, \qquad s \approx 0.774\ \mathrm{L}. \]

\textbf{Step 3: Choose the distribution.} $\sigma$ is unknown, $n = 15$ is small, and the yields are assumed normal, so we use the $t$-distribution with $\nu = n - 1 = 14$ degrees of freedom. From tables, $t_{0.025,\,14} = 2.145$.

\textbf{Step 4: Compute the interval.}
\[ E = t_{0.025,\,14}\,\frac{s}{\sqrt{n}} = 2.145 \times \frac{0.774}{\sqrt{15}} = 2.145 \times 0.1999 \approx 0.429\ \mathrm{L}. \]
\[ \bar{x} \pm E = 12.787 \pm 0.429 = [12.36\,;\,13.22]\ \mathrm{L}. \]

\textbf{Conclusion.} We are $95\%$ confident that the true mean daily milk yield lies between approximately $12.36\ \mathrm{L}$ and $13.22\ \mathrm{L}$.
\end{corrige}

\begin{corrige}[Exercise 7 — Confidence Interval for a Proportion]
\textbf{Given:} $n = 400$, $x = 210$, so $\hat{p} = \dfrac{210}{400} = 0.525$. For $95\%$ confidence, $z = 1.96$.

\textbf{a) Wald (standard) interval:}
\[ \hat{p} \pm z\sqrt{\frac{\hat{p}(1-\hat{p})}{n}} = 0.525 \pm 1.96\sqrt{\frac{0.525 \times 0.475}{400}}. \]
\[ \sqrt{\frac{0.249375}{400}} = \sqrt{0.0006234} \approx 0.02497. \]
\[ E = 1.96 \times 0.02497 \approx 0.0489. \]
\[ \mathrm{CI}_{\text{Wald}} = 0.525 \pm 0.0489 = [0.476\,;\,0.574]. \]

\textbf{b) Wilson score interval:}
\[ \frac{\hat{p} + \dfrac{z^2}{2n} \pm z\sqrt{\dfrac{\hat{p}(1-\hat{p})}{n} + \dfrac{z^2}{4n^2}}}{1 + \dfrac{z^2}{n}}, \qquad z^2 = 3.8416. \]
Centre adjustment: $\dfrac{z^2}{2n} = \dfrac{3.8416}{800} = 0.004802$; denominator: $1 + \dfrac{3.8416}{400} = 1.009604$.
\[ \text{Numerator centre: } 0.525 + 0.004802 = 0.529802. \]
\[ \text{Under the root: } \frac{0.525 \times 0.475}{400} + \frac{3.8416}{4 \times 160000} = 0.0006234 + 0.0000060 = 0.0006294, \quad \sqrt{\phantom{x}} \approx 0.025088. \]
\[ \text{Half-width: } \frac{1.96 \times 0.025088}{1.009604} = \frac{0.049173}{1.009604} \approx 0.048705. \]
\[ \text{Centre: } \frac{0.529802}{1.009604} \approx 0.524763. \]
\[ \mathrm{CI}_{\text{Wilson}} = 0.5248 \pm 0.0487 = [0.476\,;\,0.573]. \]

\textbf{c) Comparison.} Here the two intervals are almost identical because $n = 400$ is large and $\hat{p} = 0.525$ is close to $0.5$ — conditions under which the Wald interval behaves well. The Wilson interval is slightly shifted towards $0.5$ and is generally \textbf{more appropriate}, especially for small samples or proportions near 0 or 1, where the Wald interval can have true coverage well below the nominal $95\%$ and can even produce bounds outside $[0,1]$. For this data set, either interval is acceptable; the conclusion is that support for Candidate A lies between roughly $47.6\%$ and $57.4\%$.
\end{corrige}

\begin{corrige}[Exercise 8 — Determining Sample Size]
\textbf{a)} The required sample size for estimating a mean with margin of error $E$ at $95\%$ confidence is
\[ n = \left(\frac{z_{\alpha/2}\,\sigma}{E}\right)^2. \]
Using $z_{0.025} = 1.96$, $\sigma \approx s = 2.1$ and $E = 0.5$:
\[ n = \left(\frac{1.96 \times 2.1}{0.5}\right)^2 = (8.232)^2 \approx 67.77. \]
Since the sample size must be a whole number and the margin of error must \emph{not exceed} $0.5\ \mathrm{L}$, we round \textbf{up}:
\[ \boxed{n = 68 \text{ cows}}. \]

\textbf{b)} Since $n \propto \dfrac{1}{E^2}$, halving the margin of error (from $0.5$ to $0.25$) multiplies the required sample size by
\[ \left(\frac{0.5}{0.25}\right)^2 = 4. \]
The sample size must increase by a factor of \textbf{4}, i.e. $n = (16.464)^2 \approx 271.1$, so $n = 272$ cows. This illustrates the general rule: to double the precision (halve the error), one must quadruple the sample size.
\end{corrige}

\courssub{I prepare for the GCE}

\begin{corrige}[Exercise 9 — Properties of Estimators and Simulation]
\textbf{a) Definitions (3 marks).}
\begin{itemize}
\item $\hat{\theta}$ is \textbf{unbiased} if $E(\hat{\theta}) = \theta$: on average, over repeated samples, the estimator equals the true parameter. \hfill (1 mark)
\item $\hat{\theta}$ is \textbf{consistent} if it converges in probability to $\theta$ as $n \to \infty$: for every $\varepsilon > 0$, $P(|\hat{\theta}_n - \theta| > \varepsilon) \to 0$. In practice, a sufficient condition is unbiasedness together with $\mathrm{Var}(\hat{\theta}_n) \to 0$. \hfill (1 mark)
\item Among unbiased estimators, $\hat{\theta}$ is \textbf{efficient} if it has the smallest possible variance; relatively, $\hat{\theta}_1$ is more efficient than $\hat{\theta}_2$ if $\mathrm{Var}(\hat{\theta}_1) < \mathrm{Var}(\hat{\theta}_2)$. \hfill (1 mark)
\end{itemize}

\textbf{b) Unbiasedness and variance of $\bar{X}$ (3 marks).}

Each $X_i$ satisfies $E(X_i) = \mu$ and $\mathrm{Var}(X_i) = \sigma^2$, and the $X_i$ are independent. By linearity of expectation:
\[ E(\bar{X}) = E\left(\frac{1}{n}\sum_{i=1}^{n} X_i\right) = \frac{1}{n}\sum_{i=1}^{n} E(X_i) = \frac{1}{n}(n\mu) = \mu. \]
Hence $\bar{X}$ is an unbiased estimator of $\mu$. \hfill (2 marks)

Using independence (variances add):
\[ \mathrm{Var}(\bar{X}) = \mathrm{Var}\left(\frac{1}{n}\sum_{i=1}^{n} X_i\right) = \frac{1}{n^2}\sum_{i=1}^{n} \mathrm{Var}(X_i) = \frac{n\sigma^2}{n^2} = \frac{\sigma^2}{n}. \]
\hfill (1 mark)

\textbf{c) Why $\bar{X}$ is preferred to $\tilde{X}$ (2 marks).}

Both are unbiased for $\mu$ (the normal distribution is symmetric, so mean and median coincide). However, for normal data,
\[ \mathrm{Var}(\bar{X}) = \frac{\sigma^2}{n} \quad < \quad \mathrm{Var}(\tilde{X}) \approx \frac{\pi\sigma^2}{2n} \approx 1.57\,\frac{\sigma^2}{n}. \]
The sample mean has a strictly smaller variance: it is the \emph{more efficient} estimator (indeed, for a normal population $\bar{X}$ attains the minimum possible variance among unbiased estimators). Hence $\bar{X}$ is preferred. \hfill (2 marks)

\textbf{d) Simulation results (4 marks).}

Theoretical standard error of the mean: $\dfrac{\sigma}{\sqrt{n}} = \dfrac{5}{\sqrt{20}} \approx 1.118$, which matches the simulated value $1.12$ — this validates the theory. \hfill (1 mark)

The simulated standard deviation of the medians is $1.40$, larger than $1.12$. The relative efficiency is
\[ \frac{\mathrm{Var}(\bar{X})}{\mathrm{Var}(\tilde{X})} \approx \left(\frac{1.12}{1.40}\right)^2 = (0.8)^2 = 0.64 \approx \frac{2}{\pi} \approx 0.637, \]
in excellent agreement with the theoretical asymptotic relative efficiency of the median to the mean. \hfill (2 marks)

Conclusion: the sample mean is about $1/0.64 \approx 1.56$ times more efficient; equivalently, the median would need about $57\%$ more observations to achieve the same precision as the mean. The simulation empirically confirms that $\bar{X}$ is the more efficient estimator of $\mu$ for normal data. \hfill (1 mark)

\emph{Examiner expectations:} definitions stated precisely with expectations/limits; clear use of linearity of $E$ and additivity of variance under independence; comparison based on variance, not on bias; correct interpretation of simulation output.
\end{corrige}

\begin{corrige}[Exercise 10 — Confidence Intervals: Interpretation and Common Mistakes]
\textbf{a) Construction of the interval (5 marks).}

Given: $n = 40$, $\bar{x} = 752\ \mathrm{ml}$, $s = 8\ \mathrm{ml}$; $\sigma$ unknown.

\textbf{Assumptions:} the sample is random and independent; fill volumes are (approximately) normally distributed, or $n = 40$ is large enough for the Central Limit Theorem to apply. \hfill (1 mark)

Since $\sigma$ is unknown, we use the $t$-distribution with $\nu = 39$ degrees of freedom: $t_{0.025,\,39} \approx 2.023$ (using $2.021$ for $\nu = 40$ is also acceptable). \hfill (1 mark)

Standard error: $\dfrac{s}{\sqrt{n}} = \dfrac{8}{\sqrt{40}} \approx 1.2649\ \mathrm{ml}$. \hfill (1 mark)

Margin of error: $E = 2.023 \times 1.2649 \approx 2.56\ \mathrm{ml}$. \hfill (1 mark)

Confidence interval:
\[ 752 \pm 2.56 = [749.4\,;\,754.6]\ \mathrm{ml}. \]
\hfill (1 mark)

\textbf{b) Calibration (2 marks).}

The target value $750\ \mathrm{ml}$ \emph{lies inside} the $95\%$ confidence interval $[749.4\,;\,754.6]\ \mathrm{ml}$. Therefore the data are \textbf{consistent} with a correctly calibrated machine at the $5\%$ level: there is no significant evidence that the true mean fill differs from $750\ \mathrm{ml}$. (However, since the interval extends up to $754.6\ \mathrm{ml}$, the manager may wish to monitor the machine, as the sample mean is on the high side.) \hfill (2 marks)

\textbf{c) Three common mistakes (6 marks — 2 each).}
\begin{enumerate}
\item \textbf{Mistake:} ``There is a $95\%$ probability that $\mu$ lies in $[749.4\,;\,754.6]$.'' \textbf{Correction:} $\mu$ is fixed, not random; the computed interval either contains it or not. The $95\%$ refers to the long-run success rate of the \emph{method}: $95\%$ of intervals built this way from repeated samples capture $\mu$.
\item \textbf{Mistake:} ``$95\%$ of the bottles contain between $749.4$ and $754.6\ \mathrm{ml}$.'' \textbf{Correction:} the interval concerns the \emph{mean} $\mu$, not individual values. The spread of individual bottles is described by $s$ (about $8\ \mathrm{ml}$), which is much wider; a statement about individuals would require a prediction/tolerance interval.
\item \textbf{Mistake:} ``If I take another sample, $95\%$ of the time its mean will fall in this interval'' (or: ``the interval contains $95\%$ of all possible sample means''). \textbf{Correction:} the interval was built from one sample to estimate $\mu$; sample means vary around $\mu$ with standard error $\sigma/\sqrt{n}$, and a different sample yields a different interval. The $95\%$ confidence describes the procedure's coverage of $\mu$, not the location of future sample means.
\end{enumerate}
(Other acceptable mistakes: believing a larger confidence level gives a narrower interval; ignoring the assumptions of randomness/normality; using $z$ instead of $t$ when $\sigma$ is unknown and $n$ is small.)
\end{corrige}

\begin{corrige}[Exercise 11 — Sample Size Determination and Practical Application]
\textbf{a) Pilot confidence interval (3 marks).}

$\hat{p} = 0.68$, $n = 100$, $z = 1.96$.
\[ \mathrm{SE} = \sqrt{\frac{0.68 \times 0.32}{100}} = \sqrt{0.002176} \approx 0.04665. \]
\[ E = 1.96 \times 0.04665 \approx 0.0914. \]
\[ \mathrm{CI} = 0.68 \pm 0.0914 = [0.589\,;\,0.771]. \]
We are $95\%$ confident that the true proportion of households using ITNs lies between about $58.9\%$ and $77.1\%$. \hfill (3 marks)

\textbf{b) Minimum sample size (5 marks).}

Formula: $n = \dfrac{z^2\,\hat{p}(1-\hat{p})}{E^2}$ with $z = 1.96$ ($z^2 = 3.8416$) and $E = 0.03$.

a) Using $\hat{p} = 0.68$:
\[ n = \frac{3.8416 \times 0.68 \times 0.32}{0.03^2} = \frac{3.8416 \times 0.2176}{0.0009} = \frac{0.83593}{0.0009} \approx 928.8. \]
Round up: $n = 929$ households. \hfill (2.5 marks)

b) Conservative approach: use $\hat{p} = 0.5$, which maximises $p(1-p) = 0.25$:
\[ n = \frac{3.8416 \times 0.25}{0.0009} = \frac{0.9604}{0.0009} \approx 1067.1. \]
Round up: $n = 1068$ households. \hfill (2.5 marks)

The conservative choice guarantees the precision whatever the true proportion, at the cost of about 139 extra households.

\textbf{c) Hypothesis test (7 marks).}

Sample: $n = 1200$, $x = 810$, so $\hat{p} = \dfrac{810}{1200} = 0.675$.

\textbf{Hypotheses} (two-tailed test): \hfill (1 mark)
\[ H_0 : p = 0.68 \qquad \text{against} \qquad H_1 : p \neq 0.68. \]

\textbf{Test statistic.} Under $H_0$:
\[ Z = \frac{\hat{p} - p_0}{\sqrt{\dfrac{p_0(1-p_0)}{n}}} = \frac{0.675 - 0.68}{\sqrt{\dfrac{0.68 \times 0.32}{1200}}} = \frac{-0.005}{\sqrt{0.00018133}} = \frac{-0.005}{0.013466} \approx -0.371. \]
(2 marks: correct standard error under $H_0$; correct value.)

\textbf{Critical region.} At the $5\%$ level for a two-tailed test, reject $H_0$ if $|Z| > 1.96$, i.e. $Z < -1.96$ or $Z > 1.96$. \hfill (2 marks)

\textbf{Decision and conclusion.} Since $|-0.371| = 0.371 < 1.96$, the test statistic does not fall in the critical region; we \textbf{do not reject} $H_0$. \hfill (1 mark)

In context: at the $5\%$ significance level, the survey of 1200 households provides \emph{no significant evidence} that the true proportion of households using insecticide-treated nets differs from $0.68$. The observed value $0.675$ is consistent with the pilot estimate; the small difference is attributable to sampling variability. \hfill (1 mark)

\emph{Examiner expectations:} hypotheses stated in terms of $p$ with correct tail(s); standard error computed with $p_0$ (not $\hat{p}$) for the test; critical value $1.96$ for a two-tailed $5\%$ test; conclusion stated in context and non-definitive (``no evidence against'' rather than ``$H_0$ is true'').
\end{corrige}

\newpage

\coursec{Summary sheet}

\begin{popfigure}
\begin{tikzpicture}[
    scale=0.9,
    transform shape,
    central/.style={circle, draw=popdark, fill=popblueL, thick, minimum size=1.2cm, text width=2.2cm, align=center, font=\bfseries},
    branch/.style={rectangle, rounded corners, draw=popdark, thick, minimum height=1cm, text width=2.5cm, align=center, font=\small\bfseries},
    branch1/.style={branch, fill=popblue!20, text=popblue},
    branch2/.style={branch, fill=poporange!20, text=poporange},
    branch3/.style={branch, fill=popgreen!20, text=popgreen},
    branch4/.style={branch, fill=poppurple!20, text=poppurple},
    branch5/.style={branch, fill=poppink!20, text=poppink},
    branch6/.style={branch, fill=popgold!20, text=popgold},
    line/.style={draw=popline, thick, -stealth}
]
\node[central] (center) {Sampling and Estimation};

% Branch 1
\node[branch1] (b1) at (90:4.5) {1. Foundations of Sampling};
\draw[line] (center) -- (b1);

% Branch 2
\node[branch2] (b2) at (30:4.5) {2. Sampling Errors \& Point Estimation};
\draw[line] (center) -- (b2);

% Branch 3
\node[branch3] (b3) at (-30:4.5) {3. Properties of Good Estimators};
\draw[line] (center) -- (b3);

% Branch 4
\node[branch4] (b4) at (-90:4.5) {4. CI for Mean (Known Variance)};
\draw[line] (center) -- (b4);

% Branch 5
\node[branch5] (b5) at (-150:4.5) {5. CI for Mean (Unknown Variance) \& Proportion};
\draw[line] (center) -- (b5);

% Branch 6
\node[branch6] (b6) at (150:4.5) {6. Margin of Error, Sample Size \& Applications};
\draw[line] (center) -- (b6);
\end{tikzpicture}
\legende{Mind map of Chapter PMS15: Sampling and Estimation}
\end{popfigure}

\begin{bilan}

\textbf{Key Definitions}
\begin{itemize}
\item \cle{Population}: The entire set of individuals or items under study (e.g., all Form 5 students in Cameroon).
\item \cle{Sample}: A subset of the population selected for observation.
\item \cle{Parameter}: A numerical characteristic of the population (e.g., population mean $\mu$, proportion $p$).
\item \cle{Statistic}: A numerical characteristic computed from the sample (e.g., sample mean $\bar{x}$, sample proportion $\hat{p}$).
\item \cle{Sampling Frame}: The list from which the sample is drawn.
\item \cle{Sampling Error}: The difference between a sample statistic and the corresponding population parameter due to chance.
\item \cle{Point Estimate}: A single value used to estimate a population parameter.
\item \cle{Confidence Interval (CI)}: An interval of plausible values for a parameter, with a stated confidence level $(1-\alpha)$.
\item \cle{Margin of Error (E)}: Half the width of a confidence interval.
\end{itemize}

\textbf{Laws, Formulas \& Theorems to Know}
\begin{itemize}
\item \textbf{Central Limit Theorem (CLT)}: For large $n$ ($n\ge 30$), $\bar{X}\sim N\!\left(\mu,\frac{\sigma^2}{n}\right)$ approximately, regardless of population distribution.
\item \textbf{CI for $\mu$ (known $\sigma$)}: $\bar{x}\pm z_{\alpha/2}\frac{\sigma}{\sqrt{n}}$.
\item \textbf{CI for $\mu$ (unknown $\sigma$, $n\ge 30$)}: $\bar{x}\pm z_{\alpha/2}\frac{s}{\sqrt{n}}$ (using sample standard deviation $s$).
\item \textbf{CI for $\mu$ (unknown $\sigma$, $n<30$, normal population)}: $\bar{x}\pm t_{\alpha/2,\,n-1}\frac{s}{\sqrt{n}}$.
\item \textbf{CI for population proportion $p$}: $\hat{p}\pm z_{\alpha/2}\sqrt{\frac{\hat{p}(1-\hat{p})}{n}}$, valid when $n\hat{p}\ge5$ and $n(1-\hat{p})\ge5$.
\item \textbf{Margin of Error}: $E = z_{\alpha/2}\frac{\sigma}{\sqrt{n}}$ (or with $s$, $t$).
\item \textbf{Sample size for mean}: $n = \left(\frac{z_{\alpha/2}\sigma}{E}\right)^2$ (round up).
\item \textbf{Sample size for proportion}: $n = \frac{z_{\alpha/2}^2 \hat{p}(1-\hat{p})}{E^2}$ (use $\hat{p}=0.5$ for maximum $n$ if no prior estimate).
\item \textbf{Properties of good estimators}: Unbiasedness ($E(\hat{\theta})=\theta$), Consistency ($\hat{\theta}\xrightarrow{P}\theta$ as $n\to\infty$), Efficiency (minimum variance among unbiased estimators), Sufficiency.
\end{itemize}

\textbf{Methods (Step-by-Step)}
\begin{enumerate}
\item \textbf{Identify the parameter} of interest ($\mu$ or $p$) and whether $\sigma$ is known.
\item \textbf{Check conditions}: random sample, independence ($n\le 10\%$ of population), normality or large $n$ (CLT), success/failure counts for proportion.
\item \textbf{Choose the correct distribution}: $Z$ (known $\sigma$ or large $n$), $t$ (unknown $\sigma$, small $n$, normal population).
\item \textbf{Find the critical value} $z_{\alpha/2}$ or $t_{\alpha/2,\,n-1}$ from tables.
\item \textbf{Compute the point estimate} ($\bar{x}$ or $\hat{p}$) and standard error.
\item \textbf{Construct the CI}: point estimate $\pm$ (critical value $\times$ standard error).
\item \textbf{Interpret in context}: ``We are $95\%$ confident that the true mean height of Upper Sixth students in Yaoundé lies between ...''
\item \textbf{Determine sample size} if required: use the appropriate formula, round up to the next integer.
\end{enumerate}

\textbf{Common Mistakes (Traps)}
\begin{itemize}
\item Using $z$ when $\sigma$ is unknown and $n<30$ (must use $t$).
\item Forgetting to check $n\hat{p}\ge5$ and $n(1-\hat{p})\ge5$ for proportion CI.
\item Confusing \emph{confidence level} with \emph{probability that the parameter lies in the interval} (the parameter is fixed; the interval is random).
\item Using the sample standard deviation $s$ in the formula for known-$\sigma$ CI.
\item Rounding the sample size down instead of up.
\item Not stating the assumptions (random sampling, independence, normality/CLT) before constructing a CI.
\item Misreading the $t$-table: using one-tailed instead of two-tailed critical value.
\end{itemize}

\textbf{GCE Examination Tips}
\begin{itemize}
\item \textbf{Show all steps}: state the formula, substitute values, give the critical value from the table (quote the table value, e.g., $z_{0.025}=1.96$, $t_{0.025,24}=2.064$).
\item \textbf{Units matter}: if the variable is in cm, the CI is in cm.
\item \textbf{Interpretation phrases}: ``We are $95\%$ confident that ...'' not ``There is a $95\%$ chance that ...''.
\item \textbf{Sample size questions}: always round \textbf{up} to the next whole number; a fractional person is impossible.
\item \textbf{When $\sigma$ unknown and $n<30$}, explicitly mention the assumption that the population is normally distributed.
\item \textbf{For proportion CI}, use $\hat{p}$ in the standard error, not $p$.
\item \textbf{Time management}: CI questions are routine; aim to finish them in 5--7 minutes each.
\item \textbf{Real-life context}: Cameroonian examples (e.g., estimating the proportion of households with access to electricity in Douala) often appear; adapt your interpretation accordingly.
\end{itemize}

\end{bilan}

\findecours

\end{document}
